% El estatuto de la mujer y los derechos del niño — Espagnol 1ère A — Cours signé OnBuch+
% Programme officiel MINESEC. Compiler avec : tectonic EA02-mujer-derechos-nino.tex
\newcommand{\DOCMATIERE}{Espagnol}
\newcommand{\DOCNIVEAU}{1ère A}
\newcommand{\DOCMODULE}{Módulo 1 : Vie familiale et intégration sociale}
\newcommand{\DOCLECON}{EA02}
\newcommand{\DOCTITRE}{El estatuto de la mujer y los derechos del niño}
% OnBuch+ — préambule « cours signés » ENRICHI (compilé avec tectonic / XeLaTeX)
% Système visuel « Pop » d'OnBuch+ : fond crème (jamais blanc), bordures encre
% épaisses, ombres « dures » décalées, titres Archivo Black en capitales.
% Version MATHÉMATIQUES : graphes (pgfplots/tikz), boîtes pédagogiques
% (définition, propriété, méthode, attention, le savais-tu, exercice résolu,
% exercice, TP, bilan), page de garde et signature OnBuch+ ancrée en bas.
%
% Macros attendues dans main.tex AVANT \input{preamble} :
%   \DOCMATIERE  \DOCNIVEAU  \DOCMODULE  \DOCLECON (numéro)  \DOCTITRE
\documentclass[11pt]{article}

\usepackage{fontspec}
\usepackage[spanish,french]{babel} % français = langue principale ; l'espagnol via \esp{..} / espagnol
\usepackage[a4paper,margin=2cm,top=3.4cm,bottom=2.8cm,headheight=1.4cm,footskip=1.3cm]{geometry}
\usepackage{amsmath,amssymb,amsfonts}
\usepackage{mathtools} % \xmapsto, \coloneqq... souvent utilisés par les modèles
\usepackage{cancel} % simplifications barrées (bilans de forces)
% \abs{z} (module, valeur absolue) et \norm{u} (norme), utilisés par les modèles
\providecommand{\abs}{}\let\abs\relax\DeclarePairedDelimiter{\abs}{\lvert}{\rvert}
\providecommand{\norm}{}\let\norm\relax\DeclarePairedDelimiter{\norm}{\lVert}{\rVert}
\usepackage[table]{xcolor} % [table] : fournit aussi \rowcolors (alternance de lignes)
\usepackage{pagecolor}
\usepackage{tikz}
\usetikzlibrary{arrows.meta,positioning,calc,shapes.geometric,shapes.misc,shapes.arrows,decorations.pathmorphing,decorations.pathreplacing,decorations.markings,patterns,shadows,mindmap,trees,fit,backgrounds,matrix,intersections,angles,quotes}
\usepackage{pgfplots}
\usepackage[version=4]{mhchem} % équations nucléaires \ce{^{235}_{92}U}
\usepackage[european,siunitx]{circuitikz} % schémas électriques
\pgfplotsset{compat=1.18}
\usepgfplotslibrary{fillbetween,groupplots} % aires entre courbes (calcul intégral)
\usepackage{enumitem}
\usepackage{fancyhdr}
% [most] charge toutes les bibliothèques tcolorbox (listings, minted, raster,
% documentation...) et gonfle inutilement la mémoire TeX sur un document long
% (~300 boîtes) : on ne charge que ce qui est réellement utilisé.
\usepackage[skins,breakable]{tcolorbox}
\usepackage{graphicx}
\usepackage{colortbl}
\usepackage{booktabs}
\usepackage{tabularx}
\usepackage{diagbox} % case d'en-tête diagonale des tableaux à double entrée
% Colonnes extensibles centrée / gauche / droite (C, L, R), souvent utilisées par les modèles
\newcolumntype{C}{>{\centering\arraybackslash}X}
\newcolumntype{L}{>{\raggedright\arraybackslash}X}
\newcolumntype{R}{>{\raggedleft\arraybackslash}X}
\usepackage{array}
\usepackage{multicol}
\usepackage{siunitx}
% parse-numbers=false : accepte aussi \SI{2,5 \times 10^{-3}}{...}
\sisetup{locale=FR,output-decimal-marker={,},group-minimum-digits=4,parse-numbers=false}
\DeclareSIUnit\ohm{\ensuremath{\symup{\Omega}}} % U+2126 absent de la police
\DeclareSIUnit\division{div} % réglages d'oscilloscope (ms/div, V/div)
\DeclareSIUnit\tr{tr} % tours (vitesse de rotation, ex. \SI{1500}{\tr\per\minute})
\DeclareSIUnit\molar{M} % concentration molaire (ex. \SI{3}{\molar}, \SI{1}{\micro\molar})
\DeclareSIUnit\an{an} % année (le modèle confond parfois avec \year, un compteur TeX) — ex. \SI{5}{\kilo\meter\per\an}
\DeclareSIUnit\FCFA{FCFA} % franc CFA (ex. \SI{500}{\FCFA\per\kilo\gram})
\DeclareSIUnit\degreeBrix{\textdegree Brix} % teneur en sucre d'un sirop/jus (réfractométrie)
\DeclareSIUnit\month{mois} % le modèle utilise parfois \month, un compteur TeX, comme unité de durée
\usepackage{qrcode}
% \degree : commande fréquemment utilisée par les modèles pour un angle ou une
% température en dehors de \SI{}{} (non fournie par les paquets chargés).
\providecommand{\degree}{^{\circ}}
% \qed (fin de démonstration), utilisé par les modèles sans amsthm
\providecommand{\qed}{\hfill$\square$}
% \barre : double barre des valeurs interdites dans un tableau de variations (tkz-tab)
\providecommand{\barre}{\ensuremath{\|}}
% environnement proof (amsthm non chargé) utilisé par les modèles
\ifdefined\proof\else\newenvironment{proof}[1][Démonstration]{\par\noindent\textit{#1.}\ }{\qed\par}\fi
\usepackage{lastpage}
\usepackage{hyperref}
\hypersetup{hidelinks}

% ============================================================
% Palette « Pop » OnBuch+ — mêmes hex que la classe P (plus_ui.dart)
% ============================================================
\definecolor{poporange}{HTML}{F59321}
\definecolor{popdark}{HTML}{E07A0C}
\definecolor{popcream}{HTML}{FFF6E8}
\definecolor{popink}{HTML}{1C1714}
\definecolor{poppurple}{HTML}{7A5AE0}
\definecolor{popgreen}{HTML}{2E9E6A}
\definecolor{poppink}{HTML}{E0457B}
\definecolor{popblue}{HTML}{2F7DE1}
\definecolor{popgold}{HTML}{C98A12}
\definecolor{popmuted}{HTML}{8A7F76}
\definecolor{popline}{HTML}{EADFD2}
\definecolor{popgray}{HTML}{8A7F76}
\colorlet{popgrey}{popgray}
% Alias tolérants (le modèle nomme parfois les couleurs différemment)
\colorlet{popwhite}{white}
\colorlet{popblack}{popink}
\colorlet{popred}{poppink}
\colorlet{popyellow}{popgold}
% Teintes claires (fonds de tableaux/encadrés) — le modèle nomme parfois
% les couleurs avec un suffixe L (ex. popblueL) sans les avoir définies.
\colorlet{poporangeL}{poporange!20}
\colorlet{popdarkL}{popdark!20}
\colorlet{poppurpleL}{poppurple!20}
\colorlet{popgreenL}{popgreen!20}
\colorlet{poppinkL}{poppink!20}
\colorlet{popblueL}{popblue!20}
\colorlet{popgoldL}{popgold!20}
\colorlet{popredL}{popred!20}
\colorlet{popgrayL}{popgray!20}
% Teintes claires pour fonds de tableaux / zones
\colorlet{popblueL}{popblue!10!white}
\colorlet{poporangeL}{poporange!14!white}
\colorlet{popgreenL}{popgreen!12!white}
\colorlet{poppurpleL}{poppurple!12!white}
\colorlet{poppinkL}{poppink!12!white}
\colorlet{popgoldL}{popgold!14!white}

\pagecolor{popcream}

% ============================================================
% Typographie
% ============================================================
\defaultfontfeatures{Ligatures=TeX}
\setmainfont{PlusJakartaSans-Medium.ttf}[Path=../../../fonts/,BoldFont=PlusJakartaSans-Bold.ttf]
% Maths en sans-serif (Fira Math) avec chiffres et lettres droites de Plus
% Jakarta, pour que les formules se fondent dans le texte.
\usepackage[math-style=ISO,bold-style=ISO]{unicode-math}
\setmathfont{FiraMath-Regular.otf}[Path=../../../fonts/]
\setmathfont{PlusJakartaSans-Medium.ttf}[Path=../../../fonts/,range={up/num,up/{Latin,latin},\mathup}]
\usepackage{icomma} % virgule décimale sans espace en mode math
% Caractères mathématiques Unicode absents des polices de texte (Plus Jakarta,
% Archivo Black) : rendus par leur équivalent mathématique, en texte comme en
% formule — sinon ils s'affichent en carré vide (ex. « Arithmétique dans ℤ »).
\usepackage{newunicodechar}
\newunicodechar{ℝ}{\ensuremath{\mathbb{R}}}
\newunicodechar{ℤ}{\ensuremath{\mathbb{Z}}}
\newunicodechar{ℂ}{\ensuremath{\mathbb{C}}}
\newunicodechar{ℕ}{\ensuremath{\mathbb{N}}}
\newunicodechar{ℚ}{\ensuremath{\mathbb{Q}}}
\newunicodechar{𝔻}{\ensuremath{\mathbb{D}}}
\newunicodechar{α}{\ensuremath{\alpha}}
\newunicodechar{β}{\ensuremath{\beta}}
\newunicodechar{γ}{\ensuremath{\gamma}}
\newunicodechar{δ}{\ensuremath{\delta}}
\newunicodechar{ε}{\ensuremath{\varepsilon}}
\newunicodechar{θ}{\ensuremath{\theta}}
\newunicodechar{λ}{\ensuremath{\lambda}}
\newunicodechar{μ}{\ensuremath{\mu}}
\newunicodechar{σ}{\ensuremath{\sigma}}
\newunicodechar{τ}{\ensuremath{\tau}}
\newunicodechar{φ}{\ensuremath{\varphi}}
\newunicodechar{ψ}{\ensuremath{\psi}}
\newunicodechar{ω}{\ensuremath{\omega}}
\newunicodechar{Σ}{\ensuremath{\Sigma}}
% majuscules produites par \MakeUppercase dans les titres (α-aminés -> Α)
\newunicodechar{Α}{\ensuremath{\alpha}}
\newunicodechar{Β}{\ensuremath{\beta}}
\newunicodechar{Γ}{\ensuremath{\Gamma}}
\newunicodechar{Δ}{\ensuremath{\Delta}}
\newunicodechar{Ε}{\ensuremath{\varepsilon}}
\newunicodechar{Θ}{\ensuremath{\Theta}}
\newunicodechar{Λ}{\ensuremath{\Lambda}}
\newunicodechar{Μ}{\ensuremath{\mu}}
\newunicodechar{Τ}{\ensuremath{\tau}}
\newunicodechar{Φ}{\ensuremath{\Phi}}
\newunicodechar{Ψ}{\ensuremath{\Psi}}
\newunicodechar{Ω}{\ensuremath{\Omega}}
\newunicodechar{↦}{\ensuremath{\mapsto}}
\newunicodechar{∈}{\ensuremath{\in}}
\newunicodechar{∉}{\ensuremath{\notin}}
\newunicodechar{∧}{\ensuremath{\wedge}}
\newunicodechar{⇔}{\ensuremath{\Leftrightarrow}}
\newunicodechar{⟺}{\ensuremath{\Longleftrightarrow}}
\newunicodechar{⇒}{\ensuremath{\Rightarrow}}
\newunicodechar{⟹}{\ensuremath{\Longrightarrow}}
\newunicodechar{∘}{\ensuremath{\circ}}
\newunicodechar{∪}{\ensuremath{\cup}}
\newunicodechar{∩}{\ensuremath{\cap}}
\newunicodechar{≡}{\ensuremath{\equiv}}
\newunicodechar{⊂}{\ensuremath{\subset}}
\newunicodechar{⊥}{\ensuremath{\perp}}
\newunicodechar{∃}{\ensuremath{\exists}}
\newunicodechar{∀}{\ensuremath{\forall}}
\newunicodechar{∛}{\ensuremath{\sqrt[3]{\,}}}
\newunicodechar{∜}{\ensuremath{\sqrt[4]{\,}}}
\newunicodechar{⃗}{\ensuremath{{}^{\to}}}
\newunicodechar{ⁿ}{\ifmmode^{n}\else\textsuperscript{\ensuremath{n}}\fi}
\newunicodechar{ˣ}{\ifmmode^{x}\else\textsuperscript{\ensuremath{x}}\fi}
\newunicodechar{ᵇ}{\ifmmode^{b}\else\textsuperscript{\ensuremath{b}}\fi}
\newunicodechar{ʳ}{\ifmmode^{r}\else\textsuperscript{\ensuremath{r}}\fi}
\newunicodechar{ᵘ}{\ifmmode^{u}\else\textsuperscript{\ensuremath{u}}\fi}
\newunicodechar{ʸ}{\ifmmode^{y}\else\textsuperscript{\ensuremath{y}}\fi}
\newunicodechar{ᵃ}{\ifmmode^{a}\else\textsuperscript{\ensuremath{a}}\fi}
\newunicodechar{ᵉ}{\ifmmode^{e}\else\textsuperscript{\ensuremath{e}}\fi}
\newunicodechar{ᵏ}{\ifmmode^{k}\else\textsuperscript{\ensuremath{k}}\fi}
\newunicodechar{ᵖ}{\ifmmode^{p}\else\textsuperscript{\ensuremath{p}}\fi}
\newunicodechar{ᴺ}{\ifmmode^{N}\else\textsuperscript{\ensuremath{N}}\fi}
\newunicodechar{ᶜ}{\ifmmode^{c}\else\textsuperscript{\ensuremath{c}}\fi}
\newunicodechar{ʰ}{\ifmmode^{h}\else\textsuperscript{\ensuremath{h}}\fi}
\newunicodechar{ᵠ}{\ifmmode^{q}\else\textsuperscript{\ensuremath{q}}\fi}
\newunicodechar{ₙ}{\ifmmode_{n}\else\textsubscript{\ensuremath{n}}\fi}
\newunicodechar{ₐ}{\ifmmode_{a}\else\textsubscript{\ensuremath{a}}\fi}
\newunicodechar{ᵢ}{\ifmmode_{i}\else\textsubscript{\ensuremath{i}}\fi}
\newunicodechar{ᵣ}{\ifmmode_{r}\else\textsubscript{\ensuremath{r}}\fi}
\newunicodechar{ₖ}{\ifmmode_{k}\else\textsubscript{\ensuremath{k}}\fi}
\newunicodechar{ᵦ}{\ifmmode_{\beta}\else\textsubscript{\ensuremath{\beta}}\fi}
\newunicodechar{ₓ}{\ifmmode_{x}\else\textsubscript{\ensuremath{x}}\fi}
\newfontfamily\popdisplay{ArchivoBlack-Regular.ttf}[Path=../../../fonts/]
\newfontfamily\poplogo{SpaceGrotesk-Bold.ttf}[Path=../../../fonts/]
\newfontfamily\popsemi{PlusJakartaSans-SemiBold.ttf}[Path=../../../fonts/]
\newfontfamily\popxbold{PlusJakartaSans-ExtraBold.ttf}[Path=../../../fonts/]
\newfontfamily\popxboldls{PlusJakartaSans-ExtraBold.ttf}[Path=../../../fonts/,LetterSpace=3.0]

\setlist[enumerate]{leftmargin=1.7em,itemsep=5pt,topsep=4pt}
\setlist[itemize]{leftmargin=1.7em,itemsep=4pt,topsep=4pt,label={\color{poporange}\textbullet}}
\setlength{\parindent}{0pt}
\setlength{\parskip}{5pt}
\renewcommand{\arraystretch}{1.25}

% pgfplots : style Pop par défaut
\pgfplotsset{
  every axis/.append style={axis line style={popink,line width=1pt},tick style={popink},
    grid style={popline,line width=0.5pt},label style={font=\small},tick label style={font=\footnotesize}},
  popcurve/.style={poporange,line width=1.8pt,no markers,smooth},
}
% Même style disponible dans un \draw TikZ simple (hors pgfplots)
\tikzset{popcurve/.style={poporange,line width=1.8pt}}
\tikzset{popgrid/.style={popline,very thin}}
\colorlet{popcurve}{poporange} % le nom du style est parfois utilisé comme couleur

% ============================================================
% Pill / squiggle
% ============================================================
\newcommand{\poppill}[3]{%
  \tikz[baseline=-0.55ex]\node[draw=popink,line width=1.2pt,rounded corners=2.6pt,%
    fill=#2,inner xsep=6.5pt,inner ysep=3pt]{%
    \popxboldls\fontsize{7}{7}\selectfont\color{#3}\MakeUppercase{#1}};%
}
\newcommand{\popsquiggle}[1]{%
  \tikz[baseline=-0.4ex]{
    \draw[#1,line width=2.6pt,line cap=round]
      plot [smooth] coordinates {(0,0.09) (0.34,0.01) (0.68,0.21) (1.02,0.01) (1.36,0.10)};
  }%
}
% Mot-clé surligné (marqueur orange sous le texte)
\newcommand{\cle}[1]{{\popxbold\color{popdark}#1}}

% ============================================================
% En-tête / pied
% ============================================================
\pagestyle{fancy}
\fancyhf{}
\renewcommand{\headrulewidth}{0pt}
\renewcommand{\footrulewidth}{0pt}
\lhead{\raisebox{-0.15ex}{\poplogo\fontsize{13}{13}\selectfont\color{popink}OnBuch\color{poporange}+}}
\rhead{\poppill{\DOCMATIERE\ \textbullet\ \DOCNIVEAU\ \textbullet\ Leçon \DOCLECON}{white}{popink}}
\lfoot{\popxbold\fontsize{7.6}{7.6}\selectfont\color{popmuted}\MakeUppercase{Cours signé OnBuch+}\ \textbullet\ {\popsemi\DOCTITRE}}
\rfoot{\tikz[baseline=-0.6ex]\node[rounded corners=8.5pt,fill=popink,minimum height=17pt,minimum width=17pt,inner xsep=5pt,inner sep=0]{\popxbold\fontsize{8}{8}\selectfont\color{popcream}\thepage\,/\,\pageref*{LastPage}};}

% ============================================================
% Titres
% ============================================================
\newcommand{\chaptitre}[2]{%
  \begin{tcolorbox}[enhanced,colback=poporange,colframe=popink,boxrule=2.6pt,arc=11pt,
      drop shadow=popink,left=15pt,right=15pt,top=12pt,bottom=11pt,before skip=3pt,after skip=18pt]
    \poppill{#2}{popink}{popcream}\par\vspace{9pt}
    {\raggedright\hyphenpenalty=10000\exhyphenpenalty=10000\popdisplay\fontsize{22}{24}\selectfont\color{popcream}\MakeUppercase{#1}\par}
  \end{tcolorbox}
}
% Section numérotée : pastille orange avec le numéro + titre Archivo
\newcounter{popsec}
\newcommand{\coursec}[1]{%
  \stepcounter{popsec}\setcounter{popsub}{0}%
  \par\vspace{16pt}\noindent
  \begin{minipage}{\linewidth}
  \tikz[baseline=-0.7ex]\node[draw=popink,line width=1.6pt,fill=poporange,rounded corners=4pt,minimum size=22pt,inner sep=0]{\popdisplay\fontsize{12}{12}\selectfont\color{popcream}\thepopsec};%
  \hspace{8pt}{\popdisplay\fontsize{15}{17}\selectfont\color{popink}\MakeUppercase{#1}}\par
  \vspace{4pt}\noindent{\color{poporange}\rule{36pt}{3.6pt}}
  \end{minipage}\par\nopagebreak\vspace{8pt}
}
\newcounter{popsub}
\newcommand{\courssub}[1]{%
  \stepcounter{popsub}%
  \par\vspace{9pt}\noindent{\popxbold\fontsize{11.5}{13}\selectfont\color{popink}{\color{poporange}\thepopsec.\thepopsub}\hspace{6pt}#1}\par\nopagebreak\vspace{4pt}
}

% ============================================================
% Boîtes pédagogiques — même recette : fond blanc, bordure épaisse colorée,
% ombre dure encre, badge plein. Titre optionnel [#1].
% ============================================================
\tcbset{popbase/.style={breakable,enhanced,colback=white,boxrule=2.3pt,arc=9pt,
  shadow={3pt}{-3pt}{0pt}{popink},left=14pt,right=14pt,top=11pt,bottom=11pt,before skip=11pt,after skip=15pt,
  fontupper=\popsemi\fontsize{10.6}{14.5}\selectfont\color{popink},
  fontlower=\popsemi\fontsize{10.6}{14.5}\selectfont\color{popink}}}
% Coche dessinée (le glyphe ✓ n'existe pas dans les polices du document)
\newcommand{\popcheck}[1]{\tikz[baseline=-0.5ex]\draw[#1,line width=1.6pt,line cap=round,line join=round](-0.13,0.0)--(-0.04,-0.09)--(0.13,0.1);}
\providecommand{\coche}{\popcheck{popgreen}}
\providecommand{\resultat}[1]{\par\smallskip\noindent\fcolorbox{popgreen}{popgreenL}{\parbox{\dimexpr\linewidth-2\fboxsep-2\fboxrule\relax}{\textbf{Résultat.} #1}}\par\smallskip}
\newcommand{\popboxhead}[3]{\poppill{#1}{#2}{white}\ifx\relax#3\relax\else\hspace{6pt}{\popxbold\fontsize{10.6}{12}\selectfont\color{popink}#3}\fi\par\vspace{7pt}}

\newtcolorbox{aretenir}[1][]{popbase,colframe=popgreen,before upper={\popboxhead{À retenir}{popgreen}{#1}}}
% Texte en espagnol (document de lecture, dialogue, poème, extrait) : cadre
% d'étude. Un titre optionnel (source, auteur) s'affiche dans l'en-tête.
\newtcolorbox{textebox}[1][]{popbase,colframe=popblue,before upper={\popboxhead{Texto}{popblue}{#1}\begin{otherlanguage}{spanish}},after upper={\end{otherlanguage}}}
% Marques ✓ / ✗ (forme correcte / fautive) souvent utilisées par les modèles
\providecommand{\cross}{\textcolor{poppink}{\ensuremath{\times}}}
\providecommand{\xmark}{\cross}\providecommand{\crossmark}{\cross}\providecommand{\cmark}{\ensuremath{\checkmark}}
% Ligne de réponse / justification à compléter (inventée par les modèles)
\providecommand{\justification}[1]{\par\noindent\textit{Justification~:} \dotfill\par}
% \textipa (package tipa non chargé) : la police ne couvre pas l'API -> simple passage
\providecommand{\textipa}[1]{#1}
% Soulignement (ulem non chargé) : \uline, \ul
\providecommand{\uline}[1]{\underline{#1}}\providecommand{\ul}[1]{\underline{#1}}
% Mot ou phrase en espagnol à l'intérieur d'un paragraphe français
\newcommand{\esp}[1]{\ifmmode\text{\textit{#1}}\else\begingroup\selectlanguage{spanish}\itshape #1\endgroup\fi}
\newtcolorbox{exemplebox}[1][]{popbase,colframe=poporange,before upper={\popboxhead{Exemple}{poporange}{#1}}}
\newtcolorbox{definition}[1][]{popbase,colframe=popblue,before upper={\popboxhead{Définition}{popblue}{#1}}}
\newtcolorbox{propriete}[1][]{popbase,colframe=poppurple,before upper={\popboxhead{Propriété}{poppurple}{#1}}}
\newtcolorbox{methode}[1][]{popbase,colframe=popgold,before upper={\popboxhead{Méthode}{popgold}{#1}}}
\newtcolorbox{attention}[1][]{popbase,colframe=poppink,before upper={\popboxhead{Attention}{poppink}{#1}}}
\newtcolorbox{experience}[1][]{popbase,colframe=popblue,colback=popblueL,before upper={\popboxhead{Expérience}{popblue}{#1}}}
% « Le savais-tu ? » : carte violette pleine (contexte, culture, Cameroun)
\newtcolorbox{savaistu}[1][]{popbase,colback=poppurple,colframe=popink,
  fontupper=\popsemi\fontsize{10.4}{14.2}\selectfont\color{white},
  before upper={\poppill{Le savais-tu\,?}{white}{poppurple}\ifx\relax#1\relax\else\hspace{6pt}{\popxbold\color{white}#1}\fi\par\vspace{7pt}}}
% Exercice résolu : énoncé en haut, \tcblower puis la solution sur fond crème
\newcounter{popexo}\newcounter{popexr}
\newtcolorbox{exoresolu}[1][]{popbase,colframe=poporange,colbacklower=popcream,
  segmentation style={popink,line width=1.2pt,dashed},
  before upper={\stepcounter{popexr}\popboxhead{Exercice résolu \thepopexr}{poporange}{#1}},
  before lower={\poppill{Solution}{popink}{popcream}\par\vspace{6pt}}}
% Exercice (non résolu en place) — la correction est dans la partie Corrigés
\newtcolorbox{exercice}[1][]{popbase,colframe=popink,
  before upper={\stepcounter{popexo}\popboxhead{Exercice \thepopexo}{popink}{#1}}}
% Corrigé
\newtcolorbox{corrige}[1][]{popbase,colframe=popgreen,colback=popgreenL,
  before upper={\popboxhead{Corrigé}{popgreen}{#1}}}
% Objectifs / prérequis / bilan
\newtcolorbox{objectifs}{popbase,colframe=popink,colback=poporangeL,
  before upper={\popboxhead{Objectifs de la leçon}{popink}{}}}
\newtcolorbox{prerequis}{popbase,colframe=popink,colback=popblueL,
  before upper={\popboxhead{Prérequis}{popblue}{}}}
\newtcolorbox{bilan}{popbase,colframe=popink,colback=white,boxrule=2.8pt,
  before upper={\popboxhead{Fiche bilan}{poporange}{L'essentiel à connaître pour l'examen}}}
% Figure encadrée avec légende
\newtcolorbox{popfigure}[1][]{enhanced,breakable=false,colback=white,colframe=popline,boxrule=1.4pt,arc=8pt,
  left=10pt,right=10pt,top=10pt,bottom=8pt,before skip=10pt,after skip=12pt,halign=center,
  lower separated=false,halign lower=center,
  fontlower=\popsemi\fontsize{9}{11.5}\selectfont\color{popmuted},
  #1}
\newcommand{\legende}[1]{\par\vspace{6pt}{\popsemi\fontsize{9}{11.5}\selectfont\color{popmuted}#1}}

% ============================================================
% Signature OnBuch+
% ============================================================
\newcommand{\poplogoinline}[1]{{\poplogo\fontsize{#1}{#1}\selectfont\color{popink}OnBuch\color{poporange}+}}
\newcommand{\signatureonbuch}{%
  \begin{tcolorbox}[enhanced,colback=popink,colframe=popink,boxrule=2.2pt,arc=10pt,
      drop shadow=poporange,left=14pt,right=14pt,top=10pt,bottom=10pt,before skip=0pt,after skip=0pt]
    \tikz[baseline=-0.7ex]\node[circle,fill=popgreen,draw=popcream,line width=1.3pt,minimum size=22pt,inner sep=0]{\popcheck{white}};%
    \hspace{9pt}\begin{minipage}[c]{0.62\linewidth}
      {\popxbold\fontsize{12}{13}\selectfont\color{popcream}Signé\ \textbullet\ {\poplogo OnBuch\color{poporange}+}}\par\vspace{2pt}
      {\raggedright\popsemi\fontsize{8}{10.5}\selectfont\color{popline}\MakeUppercase{Équipe pédagogique OnBuch+}\\\MakeUppercase{Conforme au programme officiel MINESEC}\par}
    \end{minipage}\hfill
    {\poplogo\fontsize{22}{22}\selectfont\color{popcream}OnBuch\color{poporange}+}
  \end{tcolorbox}%
}

% ============================================================
% Page de garde
% #1 titre · #2 matière · #3 niveau · #4 module · #5 n° leçon · #6 durée ·
% #7 description / objectifs (texte ou itemize)
% ============================================================
\newcommand{\pagegarde}[7]{%
  \thispagestyle{empty}
  \newgeometry{a4paper,margin=1.7cm,top=1.6cm,bottom=1.4cm}%
  \begin{tikzpicture}[remember picture,overlay]
    \fill[poporange!18] ([xshift=-3.2cm,yshift=-3.6cm]current page.north east) circle (4.6cm);
    \fill[poppurple!14] ([xshift=2.4cm,yshift=7.6cm]current page.south west) circle (3.8cm);
  \end{tikzpicture}%
  {\poplogo\fontsize{30}{30}\selectfont\color{popink}OnBuch\color{poporange}+}\hfill
  \raisebox{0.6ex}{\poppill{Cours signé \textbullet\ Premium}{popink}{popcream}}\par
  \vspace{1pt}\popsquiggle{poppurple}\par
  \vspace{8pt}
  \begin{tcolorbox}[enhanced,colback=poporange,colframe=popink,boxrule=2.8pt,arc=13pt,
      drop shadow=popink,left=18pt,right=18pt,top=12pt,bottom=13pt,after skip=10pt]
    \poppill{#2 \textbullet\ #3}{popink}{popcream}\hspace{5pt}\poppill{#4}{white}{popink}\par\vspace{8pt}
    {\popdisplay\fontsize{14}{14}\selectfont\color{popink}LEÇON #5}\par\vspace{4pt}
    {\raggedright\hyphenpenalty=10000\exhyphenpenalty=10000\popdisplay\fontsize{25}{27}\selectfont\color{popcream}\MakeUppercase{#1}\par}
    \vspace{8pt}
    {\popxbold\fontsize{9.5}{9.5}\selectfont\color{popink}\MakeUppercase{Durée conseillée : #6}}
  \end{tcolorbox}
  \begin{tcolorbox}[enhanced,colback=white,colframe=popink,boxrule=2.2pt,arc=10pt,
      drop shadow=popink,left=16pt,right=16pt,top=10pt,bottom=10pt,after skip=8pt]
    \poppill{Dans cette leçon}{poppurple}{white}\par\vspace{5pt}
    \popsemi\fontsize{9.3}{12.6}\selectfont\color{popink}#7
  \end{tcolorbox}
  \noindent\begin{minipage}[t]{0.487\linewidth}
    \begin{tcolorbox}[enhanced,colback=popgreen,colframe=popink,boxrule=2.2pt,arc=10pt,
        drop shadow=popink,left=11pt,right=11pt,top=10pt,bottom=10pt,height=3.1cm,valign=center]
      \tcbox[colback=white,colframe=popink,boxrule=1.3pt,arc=5pt,boxsep=0pt,nobeforeafter,
        left=3pt,right=3pt,top=3pt,bottom=3pt,valign=center]{\qrcode[height=1.9cm]{https://play.google.com/store/apps/details?id=com.luvvix.onbuch&hl=fr}}%
      \hspace{8pt}\begin{minipage}[c]{0.5\linewidth}
        {\popdisplay\fontsize{11}{12.5}\selectfont\color{white}\MakeUppercase{Télécharge OnBuch}}\par\vspace{4pt}
        {\raggedright\popsemi\fontsize{8.4}{11}\selectfont\color{white}Gratuit \textbullet\ Google Play\\Tuteur IA, annales, concours, résultats\par}
      \end{minipage}
    \end{tcolorbox}
  \end{minipage}\hfill
  \begin{minipage}[t]{0.487\linewidth}
    \begin{tcolorbox}[enhanced,colback=poppurple,colframe=popink,boxrule=2.2pt,arc=10pt,
        drop shadow=popink,left=11pt,right=11pt,top=10pt,bottom=10pt,height=3.1cm,valign=center]
      \tikz[baseline=-0.7ex]\node[circle,fill=white,draw=popink,line width=1.3pt,minimum size=1.6cm,inner sep=0]{\popdisplay\fontsize{24}{24}\selectfont\color{poppurple}+};%
      \hspace{8pt}\begin{minipage}[c]{0.6\linewidth}
        {\popdisplay\fontsize{11}{12.5}\selectfont\color{white}\MakeUppercase{Passe à OnBuch+}}\par\vspace{4pt}
        {\raggedright\popsemi\fontsize{8.4}{11}\selectfont\color{white}Tous les cours signés, Tuteur IA illimité, fiches PDF — onglet OnBuch+ de l'app\par}
      \end{minipage}
    \end{tcolorbox}
  \end{minipage}
  \vspace{8pt}
  \popcontenu
  \vfill
  \signatureonbuch
  \restoregeometry
  \newpage
}

% Rangée « Ce cours contient » (page de garde)
\newcommand{\poptile}[3]{% #1 couleur #2 gros texte #3 légende
  \begin{tcolorbox}[enhanced,colback=#1,colframe=popink,boxrule=2pt,arc=9pt,shadow={2.5pt}{-2.5pt}{0pt}{popink},
      left=6pt,right=6pt,top=9pt,bottom=9pt,halign=center,valign=center,height=1.75cm,width=\linewidth,nobeforeafter]
    {\popdisplay\fontsize{12.5}{13}\selectfont\color{white}\MakeUppercase{#2}}\par\vspace{3pt}
    {\centering\popsemi\fontsize{7.8}{9.5}\selectfont\color{white}#3\par}
  \end{tcolorbox}}
\newcommand{\popcontenu}{%
  \par\noindent{\popxboldls\fontsize{7.5}{7.5}\selectfont\color{popmuted}CE COURS SIGNÉ CONTIENT}\par\vspace{7pt}
  \noindent\begin{minipage}[t]{0.235\linewidth}\poptile{popblue}{Cours}{complet, illustré, pas à pas}\end{minipage}\hfill
  \begin{minipage}[t]{0.235\linewidth}\poptile{poporange}{Méthodes}{et exercices résolus}\end{minipage}\hfill
  \begin{minipage}[t]{0.235\linewidth}\poptile{poppink}{Exercices}{type examen + corrigés}\end{minipage}\hfill
  \begin{minipage}[t]{0.235\linewidth}\poptile{popgreen}{Fiche bilan}{l'essentiel à retenir}\end{minipage}\par}

% Sommaire
\newtcolorbox{sommaire}{enhanced,colback=white,colframe=popink,boxrule=2.2pt,arc=10pt,
  drop shadow=popink,left=18pt,right=18pt,top=13pt,bottom=13pt,before skip=6pt,after skip=20pt,
  before upper={\poppill{Sommaire}{poppurple}{white}\par\vspace{9pt}\popsemi\fontsize{10.6}{14.5}\selectfont\color{popink}}}

% Signature de fin de cours — poussée tout en bas de la dernière page
\newcommand{\findecours}{%
  \par\vfill\nopagebreak
  \begin{center}\popsquiggle{poporange}\end{center}\vspace{4pt}
  \signatureonbuch
}
\AtBeginDocument{\def\llbracket{\mathopen{[\mkern-3.5mu[}}\def\rrbracket{\mathclose{]\mkern-3.5mu]}}} % intervalles d'entiers (glyphe ⟦ absent de la police)
% Glyphes absents de Fira Math : redéfinis à partir de glyphes disponibles
\AtBeginDocument{%
  \def\setminus{\mathbin{\backslash}}\let\smallsetminus\setminus
  \def\vdots{\mathord{\vbox{\baselineskip3.2pt\lineskiplimit0pt\kern4pt\hbox{.}\hbox{.}\hbox{.}}}}%
  \def\ddots{\mathinner{\mkern1mu\raise6.5pt\hbox{.}\mkern2mu\raise3.6pt\hbox{.}\mkern2mu\raise0.7pt\hbox{.}\mkern1mu}}%
  \def\triangle{\mathord{\scalebox{0.9}{$\symup{\Delta}$}}}%
  \def\nabla{\mathord{\raisebox{\depth}{\scalebox{1}[-1]{$\symup{\Delta}$}}}}%
  \def\checkmark{\popcheck{popdark}}%
  \def\star{\mathbin{\ast}}\def\bigstar{\mathord{\ast}}%
  \def\diamond{\mathbin{\scalebox{0.6}{$\Diamond$}}}%
  \def\bigcup{\mathop{\mathchoice{\raisebox{-0.3em}{\scalebox{1.7}{$\cup$}}}{\scalebox{1.25}{$\cup$}}{\cup}{\cup}}\displaylimits}%
  \def\bigcap{\mathop{\mathchoice{\raisebox{-0.3em}{\scalebox{1.7}{$\cap$}}}{\scalebox{1.25}{$\cap$}}{\cap}{\cap}}\displaylimits}%
  \def\lesseqqgtr{\mathrel{\lessgtr}}\def\bigoplus{\mathop{\oplus}\displaylimits}%
}
\providecommand{\ssi}{si et seulement si} % abréviation fréquente
\AtBeginDocument{\providecommand{\wideparen}{\overparen}} % arc de cercle

%% FIGFIX-BEGIN v5 — correctifs globaux des figures (docs/onbuch-generation/tools/figfix)
\usepackage{adjustbox}\usepackage{etoolbox}\tcbuselibrary{raster}
% toute figure TikZ/pgfplots plus large que la ligne est réduite à la largeur disponible
\BeforeBeginEnvironment{tikzpicture}{\begin{adjustbox}{max width=\linewidth}}
\AfterEndEnvironment{tikzpicture}{\end{adjustbox}}
% légendes pgfplots lisibles par-dessus les courbes
\pgfplotsset{every axis/.append style={legend style={fill=white,fill opacity=0.92,text opacity=1,draw=black!15,inner sep=3pt}}}
% étiquettes d'arêtes (syntaxe edge["..."]) avec fond blanc
\tikzset{every edge quotes/.append style={fill=white,inner sep=1.2pt}}
%% FIGFIX-END
\begin{document}

\pagegarde{El estatuto de la mujer y los derechos del niño}{Espagnol}{1ère A}{Módulo 1 : Vie familiale et intégration sociale}{EA02}{2 h}{Cette leçon propose la lecture et le commentaire d'un texte sur le statut de la femme et les droits de l'enfant, au Cameroun et dans le monde hispanique. Elle permet d'acquérir le lexique des droits et de l'égalité, et d'exprimer son opinion et la nécessité avec les structures creo que, es necesario que et tener que + infinitif.
\vspace{4pt}
\begin{multicols}{2}\small
\begin{itemize}
\item À la fin de cette leçon, je sais lire et comprendre un texte sur la condition de la femme et les droits de l'enfant
\item À la fin de cette leçon, je sais repérer les idées principales et les idées secondaires d'un texte argumentatif
\item À la fin de cette leçon, je sais utiliser le lexique des droits : igualdad, derechos, escolarización, trabajo infantil, maltrato, protección
\item À la fin de cette leçon, je sais présenter la Convention relative aux droits de l'enfant et citer des exemples au Cameroun et dans le monde hispanique
\item À la fin de cette leçon, je sais exprimer mon opinion avec creo que, pienso que, me parece que + indicatif
\item À la fin de cette leçon, je sais exprimer la nécessité et l'obligation avec es necesario que + subjonctif et tener que + infinitif
\end{itemize}\end{multicols}}

\begin{sommaire}
\begin{multicols}{2}
\begin{enumerate}
\item Texte support : La mujer y los derechos del niño
\item Compréhension du texte
\item Lexique thématique : la mujer, el niño y los derechos
\item Structures : opinion, nécessité et obligation
\item Méthode du commentaire de texte et commentaire modèle
\item Production guidée : opinar y defender los derechos
\item Activité d'intégration
\item Méthodes et exercices résolus
\item Exercices
\item Corrigés des exercices
\item Fiche bilan
\end{enumerate}
\end{multicols}
\end{sommaire}

\begin{objectifs}
\begin{itemize}
\item À la fin de cette leçon, je sais lire et comprendre un texte sur la condition de la femme et les droits de l'enfant
\item À la fin de cette leçon, je sais repérer les idées principales et les idées secondaires d'un texte argumentatif
\item À la fin de cette leçon, je sais utiliser le lexique des droits : igualdad, derechos, escolarización, trabajo infantil, maltrato, protección
\item À la fin de cette leçon, je sais présenter la Convention relative aux droits de l'enfant et citer des exemples au Cameroun et dans le monde hispanique
\item À la fin de cette leçon, je sais exprimer mon opinion avec creo que, pienso que, me parece que + indicatif
\item À la fin de cette leçon, je sais exprimer la nécessité et l'obligation avec es necesario que + subjonctif et tener que + infinitif
\item À la fin de cette leçon, je sais rédiger un court commentaire de texte structuré (introduction, développement, conclusion)
\item À la fin de cette leçon, je sais défendre l'égalité et la protection de l'enfant dans une production écrite ou orale
\end{itemize}
\end{objectifs}

\begin{prerequis}
\begin{itemize}
\item Le présent de l'indicatif des verbes réguliers et irréguliers courants (Seconde)
\item Les structures de base de l'opinion : me gusta, estoy de acuerdo (Seconde)
\item Le vocabulaire de la famille et de l'école (Seconde)
\item L'initiation au subjonctif présent après il faut que / es necesario que (début d'année de Première)
\item La lecture méthodique d'un texte simple : titre, source, thème (Seconde)
\end{itemize}
\end{prerequis}

\newpage

\coursec{Texte support : La mujer y los derechos del niño}

\courssub{Situation d'accroche et problématique}

En 2023, au Cameroun, près d'une fille sur trois quitte l'école avant la fin du premier cycle, souvent pour se marier ou pour travailler. Pendant ce temps, en Espagne ou au Mexique, des lois garantissent l'égalité entre les femmes et les hommes et protègent les enfants grâce à la Convention relative aux droits de l'enfant. Mais que dit exactement cette Convention ? Et comment parler de ces problèmes en espagnol ?

Dans cette leçon, nous allons lire un texte sur la condition de la femme et les droits de l'enfant, enrichir notre vocabulaire et apprendre à donner notre opinion : \esp{Creo que todas las niñas deben ir a la escuela.} « Je pense que toutes les filles doivent aller à l'école. » \esp{¿Y tú, qué piensas?} « Et toi, qu'en penses-tu ? »

\begin{methode}[Problématique de la leçon]
\begin{enumerate}
\item De quoi parle le texte ? Quel problème dénonce-t-il ?
\item Quels mots nouveaux nous permettent de parler des droits de la femme et de l'enfant ?
\item Comment exprimer son opinion et une obligation en espagnol ?
\end{enumerate}
\end{methode}

\courssub{Présentation du texte}

Avant de lire, observons le document. Il s'agit d'un \cle{article de presse}, c'est-à-dire un texte court que l'on pourrait trouver dans un journal ou sur un site d'information, en Espagne, en Amérique latine ou en Guinée équatoriale. Ce type de texte informe le lecteur sur un problème de société et cherche souvent à le convaincre.

\begin{itemize}
\item \textbf{Source :} un article de presse (texte d'information et d'opinion).
\item \textbf{Thème :} le statut de la femme (\esp{el estatuto de la mujer}) et les droits de l'enfant (\esp{los derechos del niño}).
\item \textbf{Personnages évoqués :} les femmes (\esp{las mujeres}), les filles (\esp{las niñas}), les enfants (\esp{los niños}), les associations (\esp{las asociaciones}) et les Nations Unies (\esp{las Naciones Unidas}).
\item \textbf{Lieux évoqués :} le monde entier, le Cameroun et les pays hispanophones.
\end{itemize}

\courssub{Lecture du texte}

\textbf{Première étape : lecture silencieuse.} Lis le texte seul, sans t'arrêter sur les mots inconnus. Cherche d'abord le sens général.

\textbf{Deuxième étape : lecture à voix haute.} Le professeur lit le texte, puis des élèves lisent chacun un paragraphe. Attention à la prononciation : \esp{niño} se prononce « ni-gnio », \esp{mujer} « mou-rère », \esp{protección} « pro-tek-sione ».

\begin{textebox}[La mujer y los derechos del niño — article de presse]
En muchos países del mundo, la mujer todavía no tiene los mismos derechos que el hombre. En algunas regiones, las niñas no van a la escuela porque tienen que trabajar en casa o casarse muy jóvenes. Sin embargo, la educación de la mujer es fundamental para el desarrollo de la sociedad.

Por su parte, los niños también sufren: el trabajo infantil y el maltrato existen todavía en muchos lugares. La Convención sobre los Derechos del Niño, adoptada por las Naciones Unidas en 1989, protege a todos los menores de dieciocho años. Afirma que todo niño tiene derecho a la vida, a la educación, a la salud y a la protección contra la violencia.

En Camerún y en los países hispanos, muchas asociaciones luchan por la igualdad y contra el trabajo infantil. Creo que es necesario que todos los niños vayan a la escuela y que las mujeres tengan las mismas oportunidades que los hombres.
\end{textebox}

\textbf{Traduction du texte.} « Dans de nombreux pays du monde, la femme n'a pas encore les mêmes droits que l'homme. Dans certaines régions, les filles ne vont pas à l'école parce qu'elles doivent travailler à la maison ou se marier très jeunes. Pourtant, l'éducation de la femme est fondamentale pour le développement de la société. De leur côté, les enfants souffrent aussi : le travail des enfants et la maltraitance existent encore dans beaucoup d'endroits. La Convention relative aux droits de l'enfant, adoptée par les Nations Unies en 1989, protège tous les mineurs de moins de dix-huit ans. Elle affirme que tout enfant a droit à la vie, à l'éducation, à la santé et à la protection contre la violence. Au Cameroun et dans les pays hispanophones, de nombreuses associations luttent pour l'égalité et contre le travail des enfants. Je pense qu'il est nécessaire que tous les enfants aillent à l'école et que les femmes aient les mêmes opportunités que les hommes. »

\begin{definition}[Glossaire du texte]
\begin{itemize}
\item \esp{igualdad} : l'égalité (féminin : \esp{la igualdad})
\item \esp{derechos} : les droits (singulier : \esp{el derecho})
\item \esp{escolarización} : la scolarisation
\item \esp{trabajo infantil} : le travail des enfants
\item \esp{maltrato} : la maltraitance
\item \esp{protección} : la protection
\item \esp{matrimonio precoz} : le mariage précoce
\item \esp{ley} : la loi
\item \esp{menor} : le mineur (personne de moins de 18 ans)
\item \esp{luchar por / contra} : lutter pour / contre
\item \esp{casarse} : se marier (verbe pronominal)
\item \esp{sin embargo} : pourtant, cependant
\end{itemize}
\end{definition}

\begin{definition}[La Convención sobre los Derechos del Niño]
La \cle{Convención sobre los Derechos del Niño} (Convention relative aux droits de l'enfant) est un traité international adopté par l'Organisation des Nations Unies (ONU) le 20 novembre 1989 et ratifié par le Cameroun en 1993. Elle garantit les droits de toute personne de moins de dix-huit ans : droit à la vie, à l'éducation, à la santé, à un nom, à une famille et à la protection contre la violence et l'exploitation.
\end{definition}

\begin{savaistu}[Le 20 novembre : Día Universal del Niño]
Le 20 novembre est la \esp{Journée mondiale de l'enfance}, en espagnol \esp{el Día Universal del Niño}. Cette date a été choisie parce que c'est le jour de l'adoption de la Convention relative aux droits de l'enfant en 1989. Chaque année, dans les écoles d'Espagne, d'Amérique latine et du Cameroun, des activités rappellent que tout enfant a des droits.
\end{savaistu}

\courssub{Repérage des mots transparents}

Les \cle{mots transparents} sont des mots qui se ressemblent en espagnol et en français. Ils nous aident à comprendre un texte rapidement. Cherchons-les dans notre texte.

\begin{center}
\begin{tabularx}{\linewidth}{|X|X|X|}
\hline
\rowcolor{popblueL} Español & Français & Remarque \\
\hline
protección & protection & se prononce « pro-tek-sione » \\
\hline
educación & éducation & accent sur la ó \\
\hline
convención & convention & se prononce « kon-ben-sione » \\
\hline
discriminación & discrimination & attention à l'accent \\
\hline
sociedad & société & pas tout à fait transparent \\
\hline
violencia & violence & se prononce « bio-len-sia » \\
\hline
fundamental & fondamental & identique \\
\hline
\end{tabularx}
\end{center}

\begin{attention}[Faux amis et pièges de vocabulaire]
Ne confonds pas \esp{derecho} et \esp{derecha} !
\begin{itemize}
\item \esp{el derecho} = le droit (juridique) : \esp{los derechos humanos} = les droits de l'homme.
\item \esp{la derecha} = la droite (direction) : \esp{a la derecha} = à droite.
\end{itemize}
Autre piège : \esp{la sociedad} signifie « la société » au sens de la communauté humaine, mais aussi « la société » commerciale (une entreprise). Enfin, \esp{infantil} se prononce « in-fan-til » et signifie « des enfants » : \esp{el trabajo infantil} = le travail des enfants.
\end{attention}

\courssub{Le champ lexical dominant et le ton du texte}

Le \cle{champ lexical} est l'ensemble des mots qui appartiennent au même thème. Ici, le champ lexical dominant est celui des \cle{droits} et de la \cle{justice sociale}.

\begin{popfigure}
\begin{tikzpicture}[mindmap, grow cyclic, every node/.style={concept, align=center, font=\small}, root concept/.append style={concept color=popblue, font=\bfseries}, level 1/.append style={level distance=3.2cm, sibling angle=90}]
\node[root concept]{Derechos}
  child[concept color=poporange]{ node[align=center]{mujer\\ igualdad\\ educación\\ trabajo} }
  child[concept color=popgreen]{ node[align=center]{niño\\ escuela\\ salud\\ protección} }
  child[concept color=popgold]{ node[align=center]{leyes\\ Convención\\ ONU\\ menores} }
  child[concept color=poppink]{ node[align=center]{sociedad\\ asociaciones\\ lucha\\ desarrollo} };
\end{tikzpicture}
\legende{Carte mentale du champ lexical des droits : quatre branches autour du mot central \esp{Derechos}.}
\end{popfigure}

Le \cle{ton} du texte est double :
\begin{itemize}
\item un ton de \textbf{dénonciation} : le texte montre des problèmes graves (\esp{el trabajo infantil}, \esp{el maltrato}, \esp{casarse muy jóvenes}) ;
\item un ton d'\textbf{espoir} : le texte parle de solutions (\esp{la Convención}, \esp{las asociaciones luchan}, \esp{es necesario que...}).
\end{itemize}

\courssub{Première compréhension globale}

Répondons maintenant aux questions essentielles.

\begin{exoresolu}[Compréhension globale du texte]
Énoncé. Réponds en français, puis en espagnol avec une phrase simple :
\begin{enumerate}
\item De quoi parle le texte ?
\item Quel problème dénonce-t-il ?
\item Quelle solution le texte propose-t-il ?
\end{enumerate}
\tcblower
Solution détaillée.
\begin{enumerate}
\item Le texte parle du statut de la femme et des droits de l'enfant dans le monde. \esp{El texto habla de la mujer y de los derechos del niño.}
\item Il dénonce l'inégalité entre les femmes et les hommes, le travail des enfants et la maltraitance. \esp{Denuncia el trabajo infantil y el maltrato.}
\item Il propose l'éducation pour tous, l'application de la Convention de 1989 et l'action des associations. \esp{Propone la educación y la protección de los niños.}
\end{enumerate}
\end{exoresolu}

\begin{exercice}[À toi de jouer]
\begin{enumerate}
\item Trouve dans le texte trois mots du champ lexical des droits.
\item Vrai ou faux ? \esp{La Convención protege a los menores de veinte años.}
\item Traduis en français : \esp{Las niñas tienen que trabajar en casa.}
\end{enumerate}
\end{exercice}

\begin{corrige}[Exercice « À toi de jouer »]
\begin{enumerate}
\item Par exemple : \esp{derechos}, \esp{protección}, \esp{igualdad} (aussi : \esp{ley}, \esp{educación}, \esp{salud}).
\item Faux : la Convention protège les mineurs de \textbf{dix-huit} ans (\esp{dieciocho años}).
\item « Les filles doivent travailler à la maison. » Attention : \esp{tener que} + infinitif = devoir faire quelque chose.
\end{enumerate}
\end{corrige}

\begin{aretenir}[Ce qu'il faut retenir de cette première étape]
\begin{itemize}
\item Le texte est un article de presse qui dénonce l'inégalité femmes-hommes et la violation des droits de l'enfant, avec un ton d'espoir.
\item La \esp{Convención sobre los Derechos del Niño} (ONU, 1989) protège toute personne de moins de 18 ans.
\item Le vocabulaire clé : \esp{igualdad}, \esp{derechos}, \esp{escolarización}, \esp{trabajo infantil}, \esp{maltrato}, \esp{protección}, \esp{matrimonio precoz}, \esp{ley}.
\item Les mots transparents (\esp{protección}, \esp{educación}, \esp{convención}) aident à comprendre vite, mais attention aux faux amis comme \esp{derecho} / \esp{derecha}.
\end{itemize}
\end{aretenir}

\coursec{Compréhension du texte}

Dans la section précédente, nous avons lu et découvert le texte support \emph{La mujer y los derechos del niño}, qui présentait la situation des femmes et des enfants dans le monde, au Cameroun et dans les pays hispanophones. Lire un texte ne suffit pas : au baccalauréat comme dans la vie, il faut le \cle{comprendre} en profondeur. Cette deuxième section vous apprend à répondre aux questions de compréhension : d'abord globales (de quoi parle le texte ?), puis détaillées (que dit exactement l'auteur ?), puis personnelles (qu'en pensez-vous ?).

\courssub{La méthode pour répondre à une question de compréhension}

Avant de répondre aux questions, il faut savoir \emph{comment} répondre. Beaucoup d'élèves perdent des points non pas parce qu'ils n'ont pas compris, mais parce qu'ils répondent mal.

\begin{methode}[Répondre à une question de compréhension en quatre étapes]
\begin{enumerate}
\item \textbf{Lire la question attentivement} et souligner le mot interrogatif : \esp{¿qué?} (quoi), \esp{¿quién?} (qui), \esp{¿por qué?} (pourquoi), \esp{¿dónde?} (où), \esp{¿cuándo?} (quand).
\item \textbf{Repérer le passage du texte} qui contient la réponse. Ne répondez jamais de mémoire : retournez toujours au texte.
\item \textbf{Répondre par une phrase complète en espagnol}, avec un sujet et un verbe conjugué. On peut reprendre les mots de la question : \esp{¿Por qué no van las niñas a la escuela?} $\rightarrow$ \esp{Las niñas no van a la escuela porque...}
\item \textbf{Citer le texte entre guillemets} pour justifier : \esp{El texto dice: «...»}. La citation prouve que votre réponse vient du texte et non de votre imagination.
\end{enumerate}
\end{methode}

\begin{attention}[La réponse d'un seul mot est interdite]
En français comme en espagnol, un examinateur n'accepte pas une réponse comme \esp{Sí}, \esp{No} ou \esp{La escuela}. En espagnol, on répond toujours par une \cle{phrase complète} : sujet + verbe conjugué + complément.

Mauvaise réponse : \esp{¿Quién lucha por la igualdad? — Las mujeres.}

Bonne réponse : \esp{Las mujeres, las asociaciones y los gobiernos luchan por la igualdad.}

Autre piège : ne recopiez pas trois lignes du texte sans les comprendre. Citez seulement le passage utile, entre guillemets.
\end{attention}

\courssub{Compréhension globale : de quoi parle le texte ?}

La compréhension globale cherche les grandes idées. Trois questions suffisent presque toujours : le \cle{thème}, le \cle{problème dénoncé}, la \cle{solution proposée}.

\begin{exemplebox}[Questions de compréhension globale et réponses modèles]
\textbf{1. ¿De qué habla el texto?} (De quoi parle le texte ?)

\esp{El texto habla del estatuto de la mujer y de los derechos del niño en el mundo, en Camerún y en los países hispanohablantes.}

« Le texte parle du statut de la femme et des droits de l'enfant dans le monde, au Cameroun et dans les pays hispanophones. »

\textbf{2. ¿Qué problema denuncia el autor?} (Quel problème l'auteur dénonce-t-il ?)

\esp{El autor denuncia la discriminación contra las mujeres y los niños: muchas niñas no van a la escuela y algunos niños tienen que trabajar.}

« L'auteur dénonce la discrimination contre les femmes et les enfants : beaucoup de filles ne vont pas à l'école et certains enfants doivent travailler. »

\textbf{3. ¿Qué solución propone el texto?} (Quelle solution le texte propose-t-il ?)

\esp{El texto propone la educación de las niñas, la protección de los niños y el respeto de la Convención sobre los Derechos del Niño de 1989.}

« Le texte propose l'éducation des filles, la protection des enfants et le respect de la Convention relative aux droits de l'enfant de 1989. »
\end{exemplebox}

L'organigramme suivant résume le parcours de lecture d'un texte argumentatif. Utilisez-le pour tout texte du programme.

\begin{popfigure}
\begin{center}
\begin{tikzpicture}[
box/.style={draw=popblue, fill=popblueL, rounded corners, align=center, text width=4.2cm, minimum height=0.9cm, font=\small},
fleche/.style={-{Stealth[length=2.5mm]}, thick, popdark}]
\node[box] (titre) {\textbf{Título} : \emph{La mujer y los derechos del niño}};
\node[box, below=0.55cm of titre] (tema) {\textbf{Tema} : la igualdad y la protección};
\node[box, below=0.55cm of tema] (problema) {\textbf{Problema} : discriminación, niñas sin escuela, trabajo infantil};
\node[box, below=0.55cm of problema] (argumentos) {\textbf{Argumentos} : la Convención de 1989, ejemplos de Camerún y del mundo hispánico};
\node[box, below=0.55cm of argumentos, fill=popgreenL, draw=popgreen] (conclusion) {\textbf{Conclusión del autor} : educar a las niñas y proteger a los niños};
\draw[fleche] (titre) -- (tema);
\draw[fleche] (tema) -- (problema);
\draw[fleche] (problema) -- (argumentos);
\draw[fleche] (argumentos) -- (conclusion);
\end{tikzpicture}
\end{center}
\legende{Organigramme de lecture : du titre à la conclusion de l'auteur}
\end{popfigure}

\courssub{Compréhension détaillée : les questions précises}

Les questions détaillées portent sur un passage précis du texte. Voici les trois questions classiques de cette leçon, avec leurs réponses modèles.

\begin{exemplebox}[Questions de compréhension détaillée]
\textbf{1. ¿Por qué no van algunas niñas a la escuela?} (Pourquoi certaines filles ne vont-elles pas à l'école ?)

\esp{Algunas niñas no van a la escuela porque tienen que trabajar en casa, cuidar a sus hermanos pequeños o ayudar a su madre en el mercado.}

« Certaines filles ne vont pas à l'école parce qu'elles doivent travailler à la maison, garder leurs petits frères et sœurs ou aider leur mère au marché. »

\textbf{2. ¿Qué dice la Convención de 1989?} (Que dit la Convention de 1989 ?)

\esp{La Convención de 1989 dice que todo niño tiene derecho a la educación, a la salud, a un nombre y a la protección contra el maltrato y el trabajo infantil.}

« La Convention de 1989 dit que tout enfant a droit à l'éducation, à la santé, à un nom et à la protection contre la maltraitance et le travail des enfants. »

\textbf{3. ¿Quién lucha por la igualdad?} (Qui lutte pour l'égalité ?)

\esp{Luchan por la igualdad las mujeres, las asociaciones, los profesores, los gobiernos y las organizaciones internacionales.}

« Luttent pour l'égalité les femmes, les associations, les professeurs, les gouvernements et les organisations internationales. »
\end{exemplebox}

Observez la structure de la réponse 2 : \esp{todo niño tiene derecho a la educación}. Cette construction est essentielle dans le lexique des droits.

\begin{propriete}[La structure « tener derecho a »]
\esp{Tener derecho a} + nom ou + infinitif signifie « avoir le droit de / avoir droit à ».

\begin{itemize}
\item \esp{tener derecho a} + \textbf{nom} : \esp{Todo niño tiene derecho a la educación.} « Tout enfant a droit à l'éducation. »
\item \esp{tener derecho a} + \textbf{infinitif} : \esp{Los niños tienen derecho a ir a la escuela.} « Les enfants ont le droit d'aller à l'école. »
\end{itemize}

Attention à la préposition : c'est toujours \esp{a}, jamais \esp{de}. On dit \esp{derecho a la educación} et non \esp{derecho de la educación}. Autres exemples : \esp{tener derecho a la salud} (avoir droit à la santé), \esp{tener derecho a un nombre} (avoir droit à un nom), \esp{tener derecho a jugar} (avoir le droit de jouer), \esp{tener derecho a vivir en familia} (avoir le droit de vivre en famille).
\end{propriete}

\begin{exemplebox}[Exemples traduits à retenir]
\esp{Las niñas tienen que trabajar en casa.} « Les filles doivent travailler à la maison. »

\esp{Todo niño tiene derecho a la educación.} « Tout enfant a droit à l'éducation. »

\esp{Los niños tienen derecho a jugar y a descansar.} « Les enfants ont le droit de jouer et de se reposer. »

\esp{Muchas mujeres luchan por la igualdad en el trabajo.} « Beaucoup de femmes luttent pour l'égalité au travail. »
\end{exemplebox}

\courssub{Questions de justification : relever les passages du texte}

Une question de justification demande de \cle{citer} le texte. La formulation typique est : \esp{Busca en el texto un pasaje que muestre...} (« Cherche dans le texte un passage qui montre... »). La réponse doit contenir la citation exacte entre guillemets, puis une courte explication avec vos propres mots.

\begin{exemplebox}[Justifier par une citation]
\textbf{1. Cita un pasaje que muestre la discriminación.} (Cite un passage qui montre la discrimination.)

\esp{El texto dice: «muchas niñas no van a la escuela porque tienen que trabajar en casa». Este pasaje muestra que las niñas no tienen las mismas oportunidades que los niños.}

« Le texte dit : "beaucoup de filles ne vont pas à l'école parce qu'elles doivent travailler à la maison". Ce passage montre que les filles n'ont pas les mêmes chances que les garçons. »

\textbf{2. Cita un pasaje que muestre la esperanza.} (Cite un passage qui montre l'espoir.)

\esp{El texto dice: «cada día más niñas van a la escuela y más mujeres trabajan como médicas, profesoras o ingenieras». Este pasaje muestra que la situación mejora poco a poco.}

« Le texte dit : "chaque jour davantage de filles vont à l'école et davantage de femmes travaillent comme médecins, professeures ou ingénieures". Ce passage montre que la situation s'améliore petit à petit. »
\end{exemplebox}

Remarquez le vocabulaire utile pour commenter une citation : \esp{este pasaje muestra que...} (ce passage montre que...), \esp{estas palabras significan que...} (ces mots signifient que...), \esp{el autor quiere decir que...} (l'auteur veut dire que...).

\courssub{Vrai ou faux : vérifier sa compréhension}

L'exercice de \esp{verdadero o falso} (vrai ou faux) est un classique du baccalauréat. Règle d'or : chaque réponse doit être justifiée par une \cle{citation du texte}, même quand la réponse est « vrai ».

\begin{exoresolu}[Verdadero o falso : exercice résolu]
Lisez les phrases suivantes et dites si elles sont \esp{verdaderas} (V) ou \esp{falsas} (F). Justifiez chaque réponse avec une citation du texte.

1. Todas las niñas del mundo van a la escuela. (Toutes les filles du monde vont à l'école.)

2. La Convención sobre los Derechos del Niño es de 1989. (La Convention relative aux droits de l'enfant date de 1989.)

3. El trabajo infantil no existe en Camerún. (Le travail des enfants n'existe pas au Cameroun.)

4. La situación de las mujeres mejora poco a poco. (La situation des femmes s'améliore petit à petit.)
\tcblower
\textbf{Corrección :}

1. \textbf{Falso.} \esp{El texto dice: «muchas niñas no van a la escuela».} La phrase est fausse : le texte affirme le contraire.

2. \textbf{Verdadero.} \esp{El texto dice: «la Convención sobre los Derechos del Niño de 1989».} La date 1989 figure dans le texte.

3. \textbf{Falso.} \esp{El texto dice: «algunos niños trabajan en los mercados o en el campo».} Le travail des enfants existe donc au Cameroun, par exemple dans les marchés ou les champs.

4. \textbf{Verdadero.} \esp{El texto dice: «cada día más niñas van a la escuela».} L'expression \esp{cada día más} (chaque jour davantage) montre un progrès lent mais réel, ce qui correspond à \esp{poco a poco} (petit à petit).

\textbf{Leçon de méthode :} pour le vrai ou faux, repérez les mots pièges comme \esp{todas} (toutes), \esp{nunca} (jamais), \esp{no existe} (n'existe pas) : ce sont souvent eux qui rendent une phrase fausse.
\end{exoresolu}

\begin{attention}[Les mots pièges du vrai ou faux]
Les phrases fausses contiennent souvent une \cle{généralisation} que le texte ne dit pas : \esp{todos} (tous), \esp{siempre} (toujours), \esp{nunca} (jamais), \esp{ningún} (aucun). Si le texte dit \esp{muchas niñas} (beaucoup de filles) et que la phrase dit \esp{todas las niñas} (toutes les filles), la phrase est \textbf{fausse}. Comparez toujours mot à mot la phrase et le texte.
\end{attention}

\courssub{Reformulation : résumer le texte avec ses propres mots}

Résumer (\esp{resumir}) consiste à dire l'essentiel du texte en deux ou trois phrases, \textbf{sans recopier} les phrases de l'auteur. On garde les idées, on change les mots.

\begin{methode}[Résumer un texte en trois étapes]
\begin{enumerate}
\item Repérez les trois idées principales : le problème, la règle ou la loi, l'espoir ou la solution.
\item Remplacez les mots du texte par des synonymes : \esp{los niños} $\rightarrow$ \esp{los menores} ; \esp{ir a la escuela} $\rightarrow$ \esp{estar escolarizados} ; \esp{mejora} $\rightarrow$ \esp{progresa}.
\item Rédigez deux ou trois phrases complètes, sans citation, sans détail, sans exemple.
\end{enumerate}
\end{methode}

\begin{exemplebox}[Résumé modèle du texte support]
\esp{El texto denuncia la discriminación que sufren las mujeres y los niños: muchas niñas no están escolarizadas y algunos menores trabajan en vez de estudiar. Sin embargo, la Convención de 1989 protege los derechos de todos los niños. El autor termina con una nota de esperanza: la igualdad progresa gracias a la educación.}

« Le texte dénonce la discrimination que subissent les femmes et les enfants : beaucoup de filles ne sont pas scolarisées et certains mineurs travaillent au lieu d'étudier. Cependant, la Convention de 1989 protège les droits de tous les enfants. L'auteur termine sur une note d'espoir : l'égalité progresse grâce à l'éducation. »
\end{exemplebox}

\courssub{Entraînement : un texte court et ses questions}

Pour vous entraîner seul, voici un petit texte original suivi de questions de compréhension. Appliquez la méthode en quatre étapes.

\begin{textebox}[Una niña de Malabo]
Carmen tiene doce años y vive en Malabo, la capital de Guinea Ecuatorial, el único país hispanohablante de África. Todos los días, su hermano pequeño va a la escuela, pero Carmen se queda en casa: tiene que cocinar, lavar la ropa y cuidar a su abuela enferma. Su madre vende pescado en el mercado y no puede pagar los libros de los dos niños.

Un día, una maestra visita a la familia y explica a la madre que Carmen tiene derecho a la educación. «La escuela pública es gratuita», dice la maestra. «Además, una niña educada puede ayudar mejor a su familia en el futuro». La madre escucha, piensa y finalmente acepta.

Hoy, Carmen va a la escuela por la mañana y ayuda a su madre por la tarde. Quiere ser enfermera. «Estudio para mis hijos del futuro», dice con una sonrisa. Su historia muestra que la educación de las niñas es posible cuando la familia, la escuela y la ley trabajan juntas.
\end{textebox}

\textbf{Glossaire :} \esp{se queda en casa} = reste à la maison ; \esp{lavar la ropa} = laver le linge ; \esp{gratuita} = gratuite ; \esp{acepta} = accepte ; \esp{enfermera} = infirmière ; \esp{una sonrisa} = un sourire ; \esp{juntas} = ensemble.

\begin{exercice}[Comprensión del texto « Una niña de Malabo »]
1. ¿Dónde vive Carmen? (Où vit Carmen ?)

2. ¿Por qué no va Carmen a la escuela al principio del texto? (Pourquoi Carmen ne va-t-elle pas à l'école au début du texte ?)

3. ¿Qué explica la maestra a la madre? (Qu'explique la maîtresse à la mère ?)

4. Verdadero o falso : \esp{La escuela pública es cara en Guinea Ecuatorial.} Justifica con una cita.

5. Cita un pasaje que muestre la esperanza. (Cite un passage qui montre l'espoir.)

6. Resume el texto en dos frases. (Résume le texte en deux phrases.)
\end{exercice}

\begin{corrige}[Exercice « Una niña de Malabo »]
1. \esp{Carmen vive en Malabo, la capital de Guinea Ecuatorial.} Phrase complète avec sujet et verbe conjugué.

2. \esp{Carmen no va a la escuela porque tiene que cocinar, lavar la ropa y cuidar a su abuela enferma.} On reprend la structure de la question : \esp{¿por qué?} $\rightarrow$ \esp{porque}.

3. \esp{La maestra explica a la madre que Carmen tiene derecho a la educación y que la escuela pública es gratuita.}

4. \textbf{Falso.} \esp{El texto dice: «la escuela pública es gratuita».} \esp{Gratuita} signifie « gratuite », donc l'école n'est pas chère.

5. \esp{El texto dice: «Hoy, Carmen va a la escuela por la mañana». Este pasaje muestra que la situación de Carmen mejora.}

6. \esp{Carmen no puede ir a la escuela porque trabaja en casa, pero una maestra convence a su madre. Finalmente, Carmen estudia y sueña con ser enfermera.}
\end{corrige}

\courssub{La question personnelle : donner son avis}

La dernière question d'une épreuve de compréhension est souvent personnelle : elle prépare le débat et l'expression orale. Elle ne demande plus de citer le texte, mais de donner votre \cle{opinion} avec une justification.

\begin{exemplebox}[Question personnelle et réponses modèles]
\textbf{¿Estás de acuerdo con el autor? ¿Por qué?} (Es-tu d'accord avec l'auteur ? Pourquoi ?)

Réponse affirmative : \esp{Sí, estoy de acuerdo con el autor porque creo que todos los niños tienen derecho a ir a la escuela. En mi barrio de Yaoundé, veo que las niñas educadas ayudan mucho a su familia.}

« Oui, je suis d'accord avec l'auteur parce que je crois que tous les enfants ont le droit d'aller à l'école. Dans mon quartier de Yaoundé, je vois que les filles éduquées aident beaucoup leur famille. »

Réponse nuancée : \esp{Estoy de acuerdo con el autor, pero pienso que no es fácil: algunas familias pobres necesitan el trabajo de sus hijos para vivir.}

« Je suis d'accord avec l'auteur, mais je pense que ce n'est pas facile : certaines familles pauvres ont besoin du travail de leurs enfants pour vivre. »
\end{exemplebox}

\begin{aretenir}[L'essentiel de la compréhension de texte]
\begin{itemize}
\item Répondez toujours par une \textbf{phrase complète} en espagnol : sujet + verbe conjugué.
\item Justifiez par une \textbf{citation entre guillemets} : \esp{El texto dice: «...»}.
\item Pour le vrai ou faux, méfiez-vous des mots pièges : \esp{todos}, \esp{siempre}, \esp{nunca}.
\item Pour résumer, gardez les idées et changez les mots : pas de recopie.
\item Structure clé : \esp{tener derecho a} + nom ou infinitif = « avoir le droit de ».
\item Pour l'opinion : \esp{estoy de acuerdo con... porque creo que...} (je suis d'accord avec... parce que je crois que...).
\end{itemize}
\end{aretenir}

Vous savez maintenant lire, comprendre et interroger le texte support. La section suivante approfondira le \cle{lexique thématique} de la femme, de l'enfant et des droits, pour enrichir vos réponses et préparer l'expression écrite et orale.

\coursec{Lexique thématique : la mujer, el niño y los derechos}

Dans la section précédente, nous avons lu et compris le texte \esp{La mujer y los derechos del niño} : nous avons identifié les idées principales (l'égalité entre hommes et femmes, la scolarisation des filles, la protection des enfants) et répondu aux questions de compréhension. Pour aller plus loin et pouvoir, à votre tour, parler et écrire sur ces sujets — à l'oral comme à l'écrit, et notamment dans les épreuves de commentaire et d'expression personnelle —, il faut maintenant \cle{enrichir votre vocabulaire}. Cette section vous donne les mots-clés du thème, organisés en quatre champs lexicaux, avec leur traduction, des exemples et des exercices.

\courssub{Observer avant d'apprendre : le lexique dans son contexte}

Avant de mémoriser des listes de mots, observons-les dans un texte authentique. Lisez ce court article de presse et repérez les mots en gras : vous verrez qu'ils se regroupent naturellement autour de quatre idées.

\begin{textebox}[Un artículo de prensa : « Educar para proteger »]
En muchos países del mundo, la \textbf{mujer} todavía no tiene las mismas \textbf{oportunidades} que el hombre. En algunas regiones, las niñas no van a la escuela porque sus familias prefieren \textbf{educar} solamente a los hijos varones. Sin embargo, la \textbf{educación} de las niñas es un \textbf{derecho} fundamental : una mujer educada puede encontrar un buen \textbf{trabajo}, recibir un \textbf{salario} justo y \textbf{defender} sus derechos.

La situación de los \textbf{niños} también preocupa a las organizaciones internacionales. Millones de menores sufren el \textbf{trabajo infantil} : trabajan en las minas, en los campos de cacao o en los mercados, en condiciones peligrosas, en vez de ir a la escuela. Otros niños son víctimas del \textbf{maltrato} dentro de su propia familia, o del \textbf{abandono} : viven en la calle, sin \textbf{protección} ni \textbf{salud}. En algunos países existe además el \textbf{matrimonio precoz} : niñas de doce o trece años se casan con hombres adultos y abandonan sus estudios.

Frente a estos problemas, la comunidad internacional actúa. La \textbf{Convención} sobre los Derechos del Niño, adoptada en 1989, obliga a los Estados a \textbf{proteger} a todos los menores. Las \textbf{leyes} nacionales prohíben el trabajo infantil y castigan la \textbf{violencia} contra los niños y las mujeres. Las asociaciones \textbf{luchan} por la \textbf{igualdad}, \textbf{denuncian} la \textbf{explotación} y ayudan a las víctimas.

En Camerún, muchas ONG trabajan para la \textbf{escolarización} de las niñas en el norte del país, y en Guinea Ecuatorial, país hispanohablante, las mujeres reclaman más respeto y más justicia. Educar a los niños y a las niñas es la mejor manera de construir un mundo más justo : hay que \textbf{respetar} los derechos de todos, sin \textbf{discriminación}.
\end{textebox}

\textbf{Glossaire :} \esp{las oportunidades} = les chances, les opportunités ; \esp{varones} = de sexe masculin ; \esp{peligrosas} = dangereuses ; \esp{en vez de} = au lieu de ; \esp{castigan} = punissent ; \esp{las ONG} = les ONG ; \esp{reclaman} = réclament ; \esp{justo / justa} = juste.

\textbf{Questions de repérage :}
\begin{enumerate}
\item Quels mots du texte désignent des \emph{problèmes} ? (Il y en a au moins cinq.)
\item Quels verbes expriment des \emph{actions pour défendre les droits} ?
\item Quel mot espagnol signifie « l'action de mettre les enfants à l'école » ?
\end{enumerate}

\textbf{Éléments de réponse :} 1. \esp{el trabajo infantil, el maltrato, el abandono, el matrimonio precoz, la violencia, la explotación, la discriminación}. 2. \esp{educar, defender, proteger, luchar, denunciar, respetar}. 3. \esp{la escolarización}.

\courssub{Champ lexical 1 : La mujer y la igualdad}

Ce premier champ rassemble les mots qui permettent de parler de la \cle{condition de la femme} : son éducation, son travail, et le grand principe de l'\esp{igualdad} (l'égalité).

\begin{definition}[Los derechos de la mujer]
\esp{Los derechos de la mujer} désignent l'ensemble des droits que possèdent les femmes au même titre que les hommes : le droit à l'éducation, au travail, à un salaire égal, à la santé et à la participation à la vie politique. On parle aussi d'\esp{igualdad entre hombres y mujeres} (égalité entre les hommes et les femmes).
\end{definition}

\begin{popfigure}
\begin{center}
\begin{tabular}{|p{5.2cm}|p{5.2cm}|}
\hline
\rowcolor{popblueL} \textbf{Español} & \textbf{Français} \\
\hline
la mujer & la femme \\
\hline
la igualdad & l'égalité \\
\hline
la discriminación & la discrimination \\
\hline
los derechos de la mujer & les droits de la femme \\
\hline
la educación & l'éducation, la formation \\
\hline
el trabajo & le travail \\
\hline
el salario & le salaire \\
\hline
la oportunidad & l'occasion, la chance \\
\hline
\end{tabular}
\end{center}
\legende{Champ lexical 1 : la mujer y la igualdad}
\end{popfigure}

\begin{exemplebox}[El léxico en frases]
\begin{itemize}
\item \esp{Las mujeres luchan por la igualdad de salarios.} « Les femmes luttent pour l'égalité des salaires. »
\item \esp{La educación de las niñas es un derecho fundamental.} « L'éducation des filles est un droit fondamental. »
\item \esp{En mi país, las mujeres tienen las mismas oportunidades de trabajo que los hombres.} « Dans mon pays, les femmes ont les mêmes chances de travail que les hommes. »
\item \esp{La discriminación contra la mujer todavía existe en algunas empresas.} « La discrimination contre la femme existe encore dans certaines entreprises. »
\end{itemize}
\end{exemplebox}

Remarquez la construction : \esp{la igualdad \textbf{de} + nom} (l'égalité \emph{de} quelque chose) : \esp{la igualdad de salarios}, \esp{la igualdad de oportunidades}, \esp{la igualdad de derechos}. Et pour comparer : \esp{las mismas \ldots\ que \ldots} « les mêmes \ldots\ que \ldots » : \esp{las mismas oportunidades que los hombres}.

\courssub{Champ lexical 2 : El niño y sus derechos}

\begin{definition}[Los derechos del niño]
\esp{Los derechos del niño} sont les droits reconnus à tout être humain de moins de dix-huit ans : droit à la vie, à un nom, à la santé (\esp{la salud}), à l'éducation (\esp{la escolarización}), à la protection (\esp{la protección}) contre toute forme de violence (\esp{la violencia}). Ils sont garantis par la Convention relative aux droits de l'enfant (1989).
\end{definition}

\begin{popfigure}
\begin{center}
\begin{tabular}{|p{5.2cm}|p{5.2cm}|}
\hline
\rowcolor{popgreenL} \textbf{Español} & \textbf{Français} \\
\hline
el niño / la niña & l'enfant (garçon / fille) \\
\hline
el menor (de edad) & le mineur \\
\hline
los derechos del niño & les droits de l'enfant \\
\hline
la escolarización & la scolarisation \\
\hline
la salud & la santé \\
\hline
la protección & la protection \\
\hline
la violencia & la violence \\
\hline
la infancia & l'enfance \\
\hline
\end{tabular}
\end{center}
\legende{Champ lexical 2 : el niño y sus derechos}
\end{popfigure}

\begin{attention}[Faux amis : escolarización et colegio]
\begin{itemize}
\item \esp{La escolarización} = la \emph{scolarisation}, c'est-à-dire l'\emph{action} de mettre les enfants à l'école. Ne dites pas \esp{escolaridad} pour traduire « scolarisation » : \esp{escolaridad} désigne plutôt le parcours scolaire ou la scolarité (les études accomplies).
\item \esp{El colegio} = l'\emph{école} (primaire ou l'établissement scolaire en général), et non « le collège » ! Pour « le collège » ou « le lycée », on dit \esp{el instituto} (ou \esp{la escuela secundaria}).
\item \esp{El niño} peut désigner « l'enfant » en général (masculin générique, comme en français) : \esp{los derechos del niño} = les droits de l'enfant (garçons et filles).
\end{itemize}
\end{attention}

\begin{exemplebox}[Ejemplos]
\begin{itemize}
\item \esp{Todos los niños tienen derecho a la salud y a la educación.} « Tous les enfants ont droit à la santé et à l'éducation. »
\item \esp{La escolarización de las niñas es una prioridad en el norte de Camerún.} « La scolarisation des filles est une priorité dans le nord du Cameroun. »
\item \esp{Los menores necesitan protección contra la violencia.} « Les mineurs ont besoin de protection contre la violence. »
\end{itemize}
\end{exemplebox}

\courssub{Champ lexical 3 : Los problemas}

Pour dénoncer une situation, il faut pouvoir nommer précisément les \cle{problèmes}. Voici le vocabulaire des atteintes aux droits.

\begin{definition}[El trabajo infantil]
\esp{El trabajo infantil} (le travail des enfants) désigne tout travail effectué par des enfants qui est dangereux pour leur santé ou qui les empêche d'aller à l'école. Il est interdit par les conventions internationales. Attention : aider ses parents quelques minutes à la maison n'est pas du \esp{trabajo infantil} ; le mot désigne un travail régulier, dur et préjudiciable.
\end{definition}

\begin{popfigure}
\begin{center}
\begin{tabular}{|p{5.2cm}|p{5.2cm}|}
\hline
\rowcolor{poporangeL} \textbf{Español} & \textbf{Français} \\
\hline
el trabajo infantil & le travail des enfants \\
\hline
el maltrato & la maltraitance \\
\hline
el matrimonio precoz & le mariage précoce \\
\hline
la pobreza & la pauvreté \\
\hline
el abandono & l'abandon \\
\hline
la explotación & l'exploitation \\
\hline
\end{tabular}
\end{center}
\legende{Champ lexical 3 : los problemas}
\end{popfigure}

\begin{exemplebox}[Ejemplos]
\begin{itemize}
\item \esp{El trabajo infantil impide la escolarización de muchos niños.} « Le travail des enfants empêche la scolarisation de beaucoup d'enfants. »
\item \esp{Hay que proteger a los niños contra el maltrato.} « Il faut protéger les enfants contre la maltraitance. »
\item \esp{El matrimonio precoz destruye el futuro de las niñas.} « Le mariage précoce détruit l'avenir des filles. »
\item \esp{La pobreza obliga a algunos menores a trabajar en los mercados.} « La pauvreté oblige certains mineurs à travailler dans les marchés. »
\end{itemize}
\end{exemplebox}

\courssub{Champ lexical 4 : Las soluciones y la lucha}

Face aux problèmes, il y a des \cle{solutions} : des lois, des institutions et des actions. Ce champ est riche en \emph{verbes} : apprenez chaque verbe avec sa préposition.

\begin{popfigure}
\begin{center}
\begin{tabular}{|p{5.2cm}|p{5.2cm}|}
\hline
\rowcolor{poppurpleL} \textbf{Español} & \textbf{Français} \\
\hline
la ley & la loi \\
\hline
la Convención & la Convention \\
\hline
proteger (a alguien) & protéger (quelqu'un) \\
\hline
luchar (por / contra) & lutter (pour / contre) \\
\hline
defender (los derechos) & défendre (les droits) \\
\hline
educar (a los niños) & éduquer (les enfants) \\
\hline
denunciar & dénoncer \\
\hline
respetar & respecter \\
\hline
\end{tabular}
\end{center}
\legende{Champ lexical 4 : las soluciones y la lucha}
\end{popfigure}

\begin{propriete}[Verbe + préposition : les constructions à mémoriser]
\begin{itemize}
\item \esp{luchar \textbf{por} + nom} = lutter pour : \esp{luchar por la igualdad} « lutter pour l'égalité ».
\item \esp{luchar \textbf{contra} + nom} = lutter contre : \esp{luchar contra el trabajo infantil}.
\item \esp{proteger \textbf{a} alguien \textbf{de/contra} algo} : \esp{proteger a los niños contra el maltrato}. N'oubliez pas le \esp{a} personnel devant la personne protégée.
\item \esp{defender} se conjugue comme \esp{perder} (e $\to$ ie) : \esp{yo defiendo, tú defiendes, él defiende, nosotros defendemos, vosotros defendéis, ellos defienden}.
\item \esp{tener derecho \textbf{a} + nom} = avoir droit à : \esp{Los niños tienen derecho a la educación.}
\item \esp{obligar \textbf{a} alguien \textbf{a} + infinitif} = obliger quelqu'un à : \esp{La pobreza obliga a los niños a trabajar.}
\end{itemize}
\end{propriete}

\begin{exemplebox}[Expresiones útiles para hablar y escribir]
\begin{itemize}
\item \esp{luchar por la igualdad} « lutter pour l'égalité » : \esp{Las asociaciones luchan por la igualdad entre hombres y mujeres.}
\item \esp{defender los derechos} « défendre les droits » : \esp{Los periodistas defienden los derechos de los niños.}
\item \esp{proteger a los niños} « protéger les enfants » : \esp{El Estado debe proteger a los niños de la calle.}
\item \esp{ir a la escuela} « aller à l'école » : \esp{Todos los niños deben ir a la escuela.}
\item \esp{tener las mismas oportunidades} « avoir les mêmes chances » : \esp{Las niñas deben tener las mismas oportunidades que los niños.}
\end{itemize}
\end{exemplebox}

\begin{savaistu}[« Ni una menos » : las mujeres se levantan]
Au Mexique et en Argentine, de grandes manifestations (\esp{marchas}) rassemblent chaque année des milliers de femmes pour défendre leurs droits. Le mouvement \esp{Ni una menos} (« Pas une de moins »), né en Argentine en 2015, dénonce les violences faites aux femmes dans toute l'Amérique latine. Son nom signifie : pas une femme de moins à cause de la violence. Ce mouvement montre que le vocabulaire de cette leçon — \esp{luchar, defender, denunciar, igualdad} — est aujourd'hui au cœur de l'actualité du monde hispanique.
\end{savaistu}

\courssub{Synthèse : carte mentale du vocabulaire}

\begin{popfigure}
\begin{center}
\begin{tikzpicture}[
  mindmap,
  every node/.style={concept, align=center, font=\small},
  root concept/.append style={fill=popblue, font=\small\bfseries},
  level 1 concept/.append style={level distance=3.6cm, sibling angle=90},
  grow cyclic
]
\node[root concept, align=center]{Derechos:\\ mujer y niño}
  child[concept color=popgreen!70]{ node[align=center]{La mujer\\ igualdad, salario,\\ educación} }
  child[concept color=poporange!80]{ node[align=center]{El niño\\ escolarización,\\ salud, protección} }
  child[concept color=poppink!80]{ node[align=center]{Problemas\\ maltrato, trabajo\\ infantil, pobreza} }
  child[concept color=poppurple!70]{ node[align=center]{Soluciones\\ ley, luchar,\\ defender, denunciar} };
\end{tikzpicture}
\end{center}
\legende{Carte mentale des quatre champs lexicaux}
\end{popfigure}

\courssub{Entraînement}

\begin{exoresolu}[Frases con huecos]
Complétez les phrases avec : \esp{escolarización} — \esp{maltrato} — \esp{igualdad} — \esp{salario} — \esp{proteger}.

1. \esp{Las mujeres luchan por la \ldots\ldots\ldots entre hombres y mujeres.}

2. \esp{El Estado debe \ldots\ldots\ldots a los niños contra la violencia.}

3. \esp{Hay que proteger a los menores contra el \ldots\ldots\ldots en la familia.}

4. \esp{Una mujer educada puede recibir un \ldots\ldots\ldots justo.}

5. \esp{La \ldots\ldots\ldots de las niñas es una prioridad en muchos países africanos.}
\tcblower
\textbf{Solution détaillée :}

1. \esp{igualdad} : on « lutte pour l'égalité » (\esp{luchar por la igualdad}).

2. \esp{proteger} : après \esp{debe}, il faut un infinitif ; le sens exige « protéger ».

3. \esp{maltrato} : la violence dans la famille, c'est la maltraitance ; on dit \esp{proteger a alguien contra el maltrato}.

4. \esp{salario} : c'est le nom masculin qui s'accorde avec \esp{un} et \esp{justo}.

5. \esp{escolarización} : attention au faux ami, on n'écrit pas \esp{escolaridad} ; il s'agit de l'action de mettre les filles à l'école.
\end{exoresolu}

\begin{exercice}[Asociación palabra / definición]
Associez chaque mot espagnol (1 à 6) à sa définition (a à f).

1. \esp{el matrimonio precoz} \hfill 2. \esp{la explotación} \hfill 3. \esp{el abandono}

4. \esp{la ley} \hfill 5. \esp{denunciar} \hfill 6. \esp{la oportunidad}

a. \esp{decir públicamente que algo es injusto o ilegal}

b. \esp{la unión de una niña muy joven con un hombre}

c. \esp{la norma escrita que obliga a todos los ciudadanos}

d. \esp{utilizar a una persona para ganar dinero sin respetarla}

e. \esp{dejar solo a un niño, sin cuidarlo}

f. \esp{la posibilidad de hacer algo}
\end{exercice}

\begin{corrige}[Exercice : asociación palabra / definición]
1 — b : le mariage précoce est l'union d'une fillette avec un homme. 2 — d : l'exploitation, c'est utiliser une personne pour le profit. 3 — e : l'abandon, c'est laisser un enfant seul. 4 — c : la loi est la norme qui oblige tous les citoyens. 5 — a : dénoncer, c'est dire publiquement qu'une chose est injuste. 6 — f : une opportunité est une possibilité, une chance.
\end{corrige}

\begin{exercice}[Traducción (thème)]
Traduisez en espagnol, en utilisant le vocabulaire de la leçon :

1. Tous les enfants ont droit à l'éducation.

2. Les femmes luttent contre la discrimination.

3. Il faut protéger les enfants contre le travail des enfants.

4. Le mariage précoce empêche les filles d'aller à l'école.
\end{exercice}

\begin{corrige}[Exercice : traducción]
1. \esp{Todos los niños tienen derecho a la educación.} (construction \esp{tener derecho a} + nom).

2. \esp{Las mujeres luchan contra la discriminación.} (verbe \esp{luchar contra} + nom).

3. \esp{Hay que proteger a los niños contra el trabajo infantil.} (« il faut » = \esp{hay que} + infinitif ; n'oubliez pas le \esp{a} personnel devant \esp{los niños}).

4. \esp{El matrimonio precoz impide a las niñas ir a la escuela.} (on accepte aussi : \esp{El matrimonio precoz impide que las niñas vayan a la escuela}, avec le subjonctif après \esp{impedir que}).
\end{corrige}

\begin{aretenir}[Lo esencial del léxico]
\begin{itemize}
\item Quatre champs : la femme et l'égalité (\esp{igualdad, salario, oportunidad}) ; l'enfant et ses droits (\esp{escolarización, salud, protección}) ; les problèmes (\esp{trabajo infantil, maltrato, matrimonio precoz, pobreza, abandono, explotación}) ; les solutions (\esp{ley, Convención, luchar, defender, denunciar, respetar}).
\item Constructions-clés : \esp{luchar por}, \esp{luchar contra}, \esp{proteger a alguien contra algo}, \esp{tener derecho a}, \esp{obligar a alguien a + infinitif}.
\item Pièges : \esp{escolarización} $\neq$ \esp{escolaridad} ; \esp{el colegio} = l'école, pas le collège ; \esp{defender} est irrégulier (\esp{yo defiendo}).
\end{itemize}
\end{aretenir}

\coursec{Structures : opinion, nécessité et obligation}

Dans la section précédente, nous avons construit le lexique thématique de la femme, de l'enfant et des droits : \esp{igualdad}, \esp{derechos}, \esp{escolarización}, \esp{trabajo infantil}, \esp{maltrato}, \esp{protección}. Ce vocabulaire ne suffit pas encore pour s'exprimer : il faut maintenant apprendre à dire ce que l'on \emph{pense}, ce qui est \emph{nécessaire} et ce que chacun \emph{doit} faire. C'est précisément l'objet de cette section : les structures de l'opinion (\esp{creo que}), de la nécessité (\esp{es necesario que}) et de l'obligation (\esp{tener que}). Ce sont les outils grammaticaux indispensables pour commenter un texte et pour défendre une opinion à l'écrit comme à l'oral.

\courssub{Observation : les structures en contexte}

Lisons d'abord un court texte original qui utilise naturellement toutes les structures que nous allons étudier. Soulignez mentalement les verbes qui suivent \esp{creo que}, \esp{no creo que}, \esp{es necesario que} et \esp{tener que}.

\begin{textebox}[Opiniones de una estudiante de Primera]
Me llamo Nadège y soy alumna de Primera en un liceo de Yaoundé. En mi clase de español, hablamos a menudo de la mujer y de los derechos del niño. Creo que la educación es fundamental para todos: cuando una niña va a la escuela, toda la familia progresa. Me parece que la igualdad es posible, pero no creo que sea fácil conseguirla. En mi barrio, todavía hay niños que trabajan en el mercado en lugar de estudiar. No creo que el trabajo infantil sea una solución para la pobreza; al contrario, pienso que destruye el futuro de los niños.

Es necesario que todos los niños vayan a la escuela y es importante que el Estado proteja a las familias pobres. También es fundamental que los padres respeten los derechos de sus hijos. Hay que luchar contra el maltrato y hay que escuchar a los niños, porque tienen mucho que decir.

En mi casa, mis padres tienen que trabajar mucho, pero siempre encuentran tiempo para ayudarme con los deberes. Yo tengo que estudiar si quiero ser médico, y mi hermano pequeño tiene que respetar las reglas de la familia. Creo que los adultos tienen que proteger a los niños, porque los niños son el futuro del país. ¿Y tú? ¿Qué piensas de los derechos del niño?
\end{textebox}

\textbf{Glossaire} : \esp{el liceo} = le lycée ; \esp{el barrio} = le quartier ; \esp{en lugar de} = au lieu de ; \esp{conseguir} = obtenir, atteindre ; \esp{los deberes} = les devoirs ; \esp{las reglas} = les règles ; \esp{el maltrato} = la maltraitance.

\textbf{Questions d'observation} :
\begin{enumerate}
\item Quelles structures Nadège utilise-t-elle pour donner son opinion ?
\item Après \esp{no creo que}, le verbe est-il à l'indicatif ou au subjonctif ?
\item Comment exprime-t-elle ce qui est nécessaire pour la société ?
\item Comment exprime-t-elle l'obligation de ses parents et la sienne ?
\end{enumerate}

Observons les faits : après \esp{creo que} et \esp{me parece que} (affirmatifs), on trouve \esp{es}, \esp{progresa}, \esp{destruye} : c'est l'\textbf{indicatif}. Après \esp{no creo que}, on trouve \esp{sea} : c'est le \textbf{subjonctif}. Après \esp{es necesario que}, on trouve \esp{vayan}, \esp{proteja}, \esp{respeten} : encore le \textbf{subjonctif}. Enfin, \esp{tener que} est toujours suivi de l'\textbf{infinitif} : \esp{tienen que trabajar}, \esp{tengo que estudiar}. Tirons maintenant les règles.

\courssub{Exprimer l'opinion : creo que, pienso que, me parece que}

\begin{propriete}[Opinion affirmative : \cle{indicatif}]
Après les verbes d'opinion à la forme \textbf{affirmative} — \esp{creo que} (je crois que), \esp{pienso que} (je pense que), \esp{me parece que} (il me semble que), \esp{opino que} (je suis d'avis que) —, on emploie l'\cle{indicatif}, car l'opinion est présentée comme un fait réel pour celui qui parle.
\par
\esp{Creo que la educación es fundamental.} « Je pense que l'éducation est fondamentale. »
\par
\esp{Pienso que la igualdad progresa en el mundo.} « Je pense que l'égalité progresse dans le monde. »
\end{propriete}

\begin{propriete}[Opinion niée : \cle{subjonctif}]
Après ces mêmes verbes à la forme \textbf{négative} — \esp{no creo que}, \esp{no pienso que}, \esp{no me parece que} —, on emploie le \cle{subjonctif}, car l'opinion est niée : le contenu de la phrase n'est plus présenté comme réel.
\par
\esp{No creo que el trabajo infantil sea una solución.} « Je ne crois pas que le travail des enfants soit une solution. »
\par
\esp{No pienso que la igualdad exista todavía en todas partes.} « Je ne pense pas que l'égalité existe déjà partout. »
\end{propriete}

\begin{exemplebox}[Exemples traduits]
\begin{itemize}
\item \esp{Creo que los niños tienen derechos.} « Je crois que les enfants ont des droits. » (\cle{indicatif} : opinion affirmée)
\item \esp{Me parece que la igualdad es posible.} « Il me semble que l'égalité est possible. » (\cle{indicatif})
\item \esp{No creo que esta ley sea suficiente.} « Je ne crois pas que cette loi soit suffisante. » (\cle{subjonctif} : opinion niée)
\item \esp{No me parece que el maltrato tenga justificación.} « Il ne me semble pas que la maltraitance ait une justification. » (\cle{subjonctif})
\end{itemize}
\end{exemplebox}

\begin{attention}[Ne pas calquer le français]
En français, on dit souvent « je ne pense pas qu'il \emph{vient} » à l'oral, mais la norme exige le subjonctif : « je ne pense pas qu'il \emph{vienne} ». En espagnol, le subjonctif est \textbf{obligatoire} et toujours respecté : on ne dit jamais \esp{no creo que es verdad}, mais \esp{no creo que sea verdad}. Autre piège : ne pas traduire « il me semble que » par \esp{me semblo que} ; la forme correcte est \esp{me parece que} (verbe \esp{parecer}, comme \esp{me gusta} : le sujet est la chose qui « semble »).
\end{attention}

\courssub{Exprimer la nécessité : es necesario que et ses variantes}

\begin{propriete}[Nécessité impersonnelle : \cle{subjonctif}]
Les expressions \textbf{impersonnelles} de nécessité — \esp{es necesario que} (il est nécessaire que), \esp{es importante que} (il est important que), \esp{es fundamental que} (il est fondamental que), \esp{es urgente que} (il est urgent que) — sont toujours suivies du \cle{subjonctif}. Le sujet du verbe subordonné est précisé dans la phrase.
\par
\esp{Es necesario que todos los niños vayan a la escuela.} « Il est nécessaire que tous les enfants aillent à l'école. »
\par
\esp{Es importante que el Estado proteja a las familias.} « Il est important que l'État protège les familles. »
\end{propriete}

\begin{propriete}[Nécessité générale sans sujet précis : \cle{hay que} + \cle{infinitif}]
Quand la nécessité est \textbf{générale} et qu'on ne désigne personne en particulier, on emploie \esp{hay que} + \cle{infinitif} (= « il faut »). Cette forme est invariable : pas de subjonctif, pas de sujet.
\par
\esp{Hay que respetar los derechos del niño.} « Il faut respecter les droits de l'enfant. »
\par
\esp{Hay que luchar contra el trabajo infantil.} « Il faut lutter contre le travail des enfants. »
\end{propriete}

\begin{attention}[L'erreur classique du francophone]
Ne dites jamais \esp{es necesario que los niños van a la escuela} avec l'indicatif : c'est la faute la plus fréquente chez les élèves francophones, par calque du français oral. Après \esp{es necesario que}, le subjonctif est \textbf{obligatoire} : \esp{es necesario que los niños vayan a la escuela}. De même, ne confondez pas \esp{hay que} (il faut, général) et \esp{tiene que} (il/elle doit, personnel) : \esp{hay} est impersonnel et ne change jamais.
\end{attention}

\courssub{Exprimer l'obligation personnelle : tener que + infinitif}

\begin{propriete}[Obligation d'une personne précise]
Pour exprimer l'obligation d'une personne déterminée, on emploie \esp{tener que} + \cle{infinitif} (= « devoir »). Le verbe \esp{tener} se conjugue selon la personne, \esp{que} reste invariable, et le second verbe reste à l'infinitif. On peut aussi employer \esp{deber} + infinitif, plus soutenu.
\par
\esp{Los padres tienen que proteger a sus hijos.} « Les parents doivent protéger leurs enfants. »
\par
\esp{Las niñas tienen que ir a la escuela.} « Les filles doivent aller à l'école. »
\par
\esp{El gobierno debe garantizar la escolarización.} « Le gouvernement doit garantir la scolarisation. »
\end{propriete}

Conjugaison de \esp{tener que} au présent :

\begin{center}
\begin{tabular}{|l|l|l|}\hline
\rowcolor{popblueL} Personne & Forme & Exemple \\ \hline
yo & tengo que & Tengo que estudiar. \\ \hline
tú & tienes que & Tienes que escuchar. \\ \hline
él / ella / usted & tiene que & Tiene que trabajar. \\ \hline
nosotros & tenemos que & Tenemos que ayudar. \\ \hline
vosotros & tenéis que & Tenéis que respetar. \\ \hline
ellos / ustedes & tienen que & Tienen que proteger. \\ \hline
\end{tabular}
\end{center}

\courssub{Rappel : le subjonctif présent des verbes fréquents}

Les structures de la nécessité et de l'opinion niée exigent le subjonctif présent. Voici les formes des verbes les plus utiles de cette leçon (troisième personne, la plus fréquente dans les phrases impersonnelles) :

\begin{center}
\begin{tabular}{|l|l|l|}\hline
\rowcolor{popblueL} Infinitif & Subjonctif (él/ellos) & Exemple \\ \hline
ir & vaya / vayan & Es necesario que vayan a la escuela. \\ \hline
ser & sea / sean & No creo que sea justo. \\ \hline
tener & tenga / tengan & Es importante que tengan derechos. \\ \hline
hacer & haga / hagan & Es necesario que hagan los deberes. \\ \hline
poder & pueda / puedan & Es fundamental que puedan estudiar. \\ \hline
estar & esté / estén & No creo que estén protegidos. \\ \hline
haber & haya & Es necesario que haya más escuelas. \\ \hline
\end{tabular}
\end{center}

Rappel de formation : pour les verbes réguliers, on prend la première personne du présent de l'indicatif, on retire \esp{-o}, et on ajoute les terminaisons « croisées » : \esp{-e/-es/-e/-emos/-éis/-en} pour \esp{-ar} (\esp{trabajar} → \esp{trabaje}), \esp{-a/-as/-a/-amos/-áis/-an} pour \esp{-er/-ir} (\esp{proteger} → \esp{proteja} ; \esp{escribir} → \esp{escriba}).

\courssub{Comparaison français / espagnol et tableau récapitulatif}

\begin{center}
\begin{tabularx}{\linewidth}{|X|X|X|}\hline
\rowcolor{popblueL} Français & Español & Remarque \\ \hline
Je pense que l'éducation est fondamentale. & Creo que la educación es fundamental. & \cle{indicatif} après opinion affirmative \\ \hline
Je ne crois pas que ce soit juste. & No creo que sea justo. & \cle{subjonctif} après opinion niée \\ \hline
Il faut que les enfants aillent à l'école. & Es necesario que los niños vayan a la escuela. & \cle{subjonctif} dans les deux langues \\ \hline
Il faut respecter les droits. & Hay que respetar los derechos. & \cle{infinitif}, forme impersonnelle \\ \hline
Les parents doivent protéger leurs enfants. & Los padres tienen que proteger a sus hijos. & \esp{tener que} / \esp{deber} + \cle{infinitif} \\ \hline
\end{tabularx}
\end{center}

\begin{popfigure}
\begin{center}
\begin{tikzpicture}[node distance=0.9cm and 1.6cm,
  caja/.style={rectangle, rounded corners, draw=popblue, fill=popblueL, align=center, minimum height=1cm, text width=4.2cm, font=\small},
  flecha/.style={-{Stealth[length=2.5mm]}, thick, popdark}]
\node[caja, fill=poporangeL, draw=poporange, align=center] (op){\textbf{Opinion personnelle}\\ \esp{creo que / pienso que}\\ \esp{me parece que} + \cle{indicatif}\\ \esp{no creo que} + \cle{subjonctif}};
\node[caja, fill=popgreenL, draw=popgreen, right=of op, align=center] (nec){\textbf{Nécessité générale}\\ \esp{es necesario que} + \cle{subjonctif}\\ \esp{hay que} + \cle{infinitif}};
\node[caja, fill=poppurpleL, draw=poppurple, right=of nec, align=center] (obl){\textbf{Obligation d'une personne}\\ \esp{tener que} + \cle{infinitif}\\ \esp{deber} + \cle{infinitif}};
\node[caja, fill=popcream, draw=popgold, below=1.1cm of nec, text width=12.5cm] (ej) {\esp{Creo que la igualdad es posible.} \quad \esp{Es necesario que los niños vayan a la escuela.} \quad \esp{Los padres tienen que proteger a sus hijos.}};
\draw[flecha] (op.south) -- (ej.north -| op.south);
\draw[flecha] (nec.south) -- (ej.north);
\draw[flecha] (obl.south) -- (ej.north -| obl.south);
\end{tikzpicture}
\end{center}
\legende{Choisir la bonne structure : opinion, nécessité ou obligation.}
\end{popfigure}

\begin{methode}[Comment choisir la bonne structure]
\begin{enumerate}
\item \textbf{Identifier l'intention} : est-ce que je donne mon avis, est-ce que je dis ce qui est nécessaire pour tous, ou est-ce que je dis ce qu'une personne précise doit faire ?
\item \textbf{Opinion personnelle} → \esp{creo que / pienso que / me parece que} + \cle{indicatif} ; si l'opinion est niée (\esp{no creo que}), passer au \cle{subjonctif}.
\item \textbf{Nécessité générale} → \esp{es necesario que / es importante que} + \cle{subjonctif} si le sujet est précisé ; \esp{hay que} + \cle{infinitif} si l'on ne désigne personne.
\item \textbf{Obligation d'une personne précise} → \esp{tener que} + \cle{infinitif}, en conjuguant \esp{tener} selon la personne.
\item \textbf{Vérifier le mode} avant de rédiger : \cle{indicatif} ou \cle{subjonctif} ? C'est le point que le correcteur regarde en premier.
\end{enumerate}
\end{methode}

\begin{exoresolu}[Compléter avec le bon mode]
Énoncé. Complétez avec la forme correcte du verbe entre parenthèses :
\begin{enumerate}
\item Creo que la escolarización (ser) \ldots\ muy importante.
\item No creo que los niños (tener) \ldots\ que trabajar.
\item Es necesario que el gobierno (hacer) \ldots\ un esfuerzo.
\item Las niñas (tener que) \ldots\ ir a la escuela.
\item (Haber) \ldots\ que respetar a las mujeres.
\end{enumerate}
\tcblower
Solution détaillée :
\begin{enumerate}
\item \esp{Creo que la escolarización \textbf{es} muy importante.} Opinion affirmative → \cle{indicatif}.
\item \esp{No creo que los niños \textbf{tengan} que trabajar.} Opinion niée → \cle{subjonctif} de \esp{tener}.
\item \esp{Es necesario que el gobierno \textbf{haga} un esfuerzo.} Nécessité impersonnelle → \cle{subjonctif} de \esp{hacer}.
\item \esp{Las niñas \textbf{tienen que} ir a la escuela.} Obligation personnelle → \esp{tener que} conjugué + \cle{infinitif}.
\item \esp{\textbf{Hay} que respetar a las mujeres.} Nécessité générale sans sujet → \esp{hay que} + \cle{infinitif}.
\end{enumerate}
\end{exoresolu}

\begin{exercice}[À vous de jouer]
\begin{enumerate}
\item Traduisez en espagnol : a) « Je pense que l'égalité est possible. » b) « Je ne crois pas que le travail des enfants soit une solution. » c) « Il faut que tous les enfants aillent à l'école. » d) « Il faut protéger les droits de l'enfant. » e) « Les parents doivent écouter leurs enfants. »
\item Rédigez trois phrases originales sur la condition de la femme au Cameroun, en utilisant une fois \esp{me parece que}, une fois \esp{es importante que} et une fois \esp{tener que}.
\end{enumerate}
\end{exercice}

\begin{corrige}[Exercice : traduction]
a) \esp{Pienso que la igualdad es posible.} (\cle{indicatif}) ; b) \esp{No creo que el trabajo infantil sea una solución.} (\cle{subjonctif}) ; c) \esp{Es necesario que todos los niños vayan a la escuela.} (\cle{subjonctif} de \esp{ir}) ; d) \esp{Hay que proteger los derechos del niño.} (\cle{infinitif}) ; e) \esp{Los padres tienen que escuchar a sus hijos.} (attention à la \esp{a} personnelle devant \esp{sus hijos}).
\end{corrige}

\begin{aretenir}[L'essentiel de la section]
Opinion affirmative : \esp{creo que} + \cle{indicatif}. Opinion niée : \esp{no creo que} + \cle{subjonctif}. Nécessité : \esp{es necesario que} + \cle{subjonctif}, ou \esp{hay que} + \cle{infinitif} (général). Obligation personnelle : \esp{tener que} + \cle{infinitif}. Le choix du mode (\cle{indicatif} / \cle{subjonctif} / \cle{infinitif}) est la clé de la correction grammaticale : c'est la première chose à vérifier dans vos productions.
\end{aretenir}

\coursec{Méthode du commentaire de texte et commentaire modèle}

Dans la section précédente, nous avons appris à exprimer l'opinion, la nécessité et l'obligation : \esp{Creo que la educación es un derecho.} « Je pense que l'éducation est un droit. » ; \esp{Es necesario que la sociedad proteja a los niños.} « Il est nécessaire que la société protège les enfants. » ; \esp{Tenemos que luchar por la igualdad.} « Nous devons lutter pour l'égalité. » Ces structures ne sont pas des exercices isolés : ce sont précisément les outils dont vous aurez besoin pour rédiger, en espagnol, un \cle{commentaire de texte} (\esp{comentario de texto}), exercice central de l'évaluation en Première et au Baccalauréat. Cette section vous donne la méthode complète, étape par étape, puis un commentaire modèle rédigé sur le texte support de la leçon.

\courssub{Qu'est-ce qu'un commentaire de texte ?}

\begin{definition}[Le comentario de texto]
Un \cle{commentaire de texte} est une rédaction courte (10 à 15 lignes en Première) dans laquelle l'élève : \textbf{présente} le texte (nature, thème), \textbf{dégage} ses idées principales en s'appuyant sur des citations courtes, puis \textbf{donne son avis personnel} argumenté. En espagnol, on parle de \esp{comentario de texto}. Il se rédige entièrement en espagnol, avec des phrases simples et correctes.
\end{definition}

Le commentaire n'est ni un résumé mot à mot, ni une traduction, ni une paraphrase. C'est un regard personnel et organisé sur le texte. Sa structure est toujours la même, en trois parties :

\begin{popfigure}
\begin{tikzpicture}[
  bloc/.style={rectangle, rounded corners=3pt, draw=popblue, fill=popblueL, text width=3.6cm, align=center, minimum height=1.1cm, font=\small},
  sub/.style={rectangle, rounded corners=3pt, draw=poporange, fill=poporangeL, text width=2.6cm, align=center, minimum height=0.8cm, font=\footnotesize},
  fleche/.style={-{Stealth[length=2.5mm]}, thick, popdark}
]
\node[bloc, align=center] (intro) at (0,0){\textbf{Introducción}\\ nature, source, thème, problématique};
\node[bloc, fill=popgreenL, draw=popgreen, align=center] (des) at (5.4,0){\textbf{Desarrollo}\\ idées principales + citations};
\node[bloc, fill=poppinkL, draw=poppink, align=center] (conc) at (10.8,0){\textbf{Conclusión}\\ bilan + avis personnel};
\node[sub] (i1) at (4.3,-1.7) {\textbf{Idea 1 :} la mujer y la igualdad};
\node[sub] (i2) at (7.3,-1.7) {\textbf{Idea 2 :} los derechos del niño};
\draw[fleche] (intro) -- (des);
\draw[fleche] (des) -- (conc);
\draw[fleche, poporange] (des.south) -- (i1.north);
\draw[fleche, poporange] (des.south) -- (i2.north);
\end{tikzpicture}
\legende{Les trois blocs du commentaire de texte : introducción, desarrollo (deux ou trois idées), conclusión.}
\end{popfigure}

\courssub{Étape 1 — L'introduction : présenter le texte}

L'introduction répond à trois questions : \textbf{Quoi ?} (la nature du texte : article de presse, lettre, discours, extrait de roman), \textbf{Qui / d'où ?} (l'auteur ou la source, si on les connaît), \textbf{De quoi parle-t-on ?} (le thème). Elle se termine par l'annonce du plan ou une problématique simple.

\begin{exemplebox}[Phrases modèles pour l'introduction]
\begin{itemize}
\item \esp{Este texto es un artículo de prensa sobre la condición de la mujer y los derechos del niño.} « Ce texte est un article de presse sur la condition de la femme et les droits de l'enfant. »
\item \esp{El autor habla de la igualdad entre hombres y mujeres.} « L'auteur parle de l'égalité entre les hommes et les femmes. »
\item \esp{El tema principal es la protección de los menores.} « Le thème principal est la protection des mineurs. »
\end{itemize}
\end{exemplebox}

\begin{attention}[Faux amis et calques dans l'introduction]
\begin{itemize}
\item \esp{la condición de la mujer} signifie « la condition / le statut de la femme », et non « la condition » au sens d'une exigence.
\item On dit \esp{un artículo de prensa} (« un article de presse ») ; ne dites pas \esp{un artículo de presión} : \esp{presión} = « pression ».
\item Pour « parler de », l'espagnol emploie \esp{hablar de} : \esp{El texto habla de los niños.} Ne calquez pas le français « parler sur ».
\end{itemize}
\end{attention}

\courssub{Étape 2 — Le développement : dégager les idées et citer}

Le développement est le cœur du commentaire. Vous dégagez \textbf{deux ou trois idées principales}, une par paragraphe. Pour chaque idée, la démarche est toujours la même : 1) vous annoncez l'idée avec un verbe de discours (\esp{denunciar}, \esp{explicar}, \esp{decir}, \esp{presentar}, \esp{defender}) ; 2) vous l'illustrez par une \textbf{citation courte} entre guillemets ; 3) vous l'expliquez ou la traduisez avec vos propres mots.

\begin{exemplebox}[Verbes et formules pour présenter les idées]
\begin{itemize}
\item \esp{El texto denuncia el trabajo infantil.} « Le texte dénonce le travail des enfants. »
\item \esp{Según el autor, la educación es un derecho.} « Selon l'auteur, l'éducation est un droit. »
\item \esp{El autor explica que muchas niñas no van a la escuela.} « L'auteur explique que beaucoup de filles ne vont pas à l'école. »
\item \esp{Después, presenta la Convención de 1989.} « Ensuite, il présente la Convention de 1989. »
\end{itemize}
\end{exemplebox}

\begin{attention}[Citer sans recopier]
Ne recopiez \textbf{jamais} de longs passages du texte : l'examinateur le sanctionne. Citez des \textbf{fragments courts} (trois à huit mots) entre guillemets, puis commentez-les. Mauvais : recopier deux lignes entières. Bon : \esp{El autor afirma que «las niñas no van a la escuela» en algunas regiones, lo que muestra la discriminación.} « L'auteur affirme que "les filles ne vont pas à l'école" dans certaines régions, ce qui montre la discrimination. »
\end{attention}

Pour enchaîner les idées, utilisez les \cle{connecteurs} (\esp{conectores}). Ils donnent du rythme et de la logique à votre texte :

\begin{center}
\begin{tabularx}{\linewidth}{|X|X|X|}
\hline
\rowcolor{popblueL} \textbf{Conector} & \textbf{Français} & \textbf{Emploi} \\
\hline
\esp{primero / en primer lugar} & d'abord / en premier lieu & annoncer la première idée \\
\hline
\esp{además / también} & de plus / aussi & ajouter une idée \\
\hline
\esp{después / luego} & ensuite / puis & passer à l'idée suivante \\
\hline
\esp{sin embargo} & cependant & opposer deux idées \\
\hline
\esp{por eso / por eso mismo} & c'est pourquoi & exprimer la conséquence \\
\hline
\esp{finalmente} & finalement & dernière idée \\
\hline
\esp{en conclusión} & en conclusion & ouvrir la conclusion \\
\hline
\end{tabularx}
\end{center}

\begin{exemplebox}[Les connecteurs en action]
\esp{Primero, el autor denuncia la discriminación contra la mujer. Además, habla del trabajo infantil. Sin embargo, el texto termina con una nota de esperanza: la Convención protege a los menores. Por eso, el mensaje es positivo.} « D'abord, l'auteur dénonce la discrimination contre la femme. De plus, il parle du travail des enfants. Cependant, le texte se termine sur une note d'espérance : la Convention protège les mineurs. C'est pourquoi le message est positif. »
\end{exemplebox}

\courssub{Étape 3 — La conclusion : bilan et avis personnel}

La conclusion fait le bilan en une phrase, puis donne votre \textbf{avis personnel argumenté}. C'est ici que les structures de la section précédente deviennent indispensables : \esp{creo que} + indicatif, \esp{me parece que} + indicatif, \esp{es necesario que} + subjonctif, \esp{hay que} + infinitif, \esp{tener que} + infinitif.

\begin{exemplebox}[Phrases modèles pour la conclusion]
\begin{itemize}
\item \esp{En conclusión, creo que hay que proteger a los niños.} « En conclusion, je pense qu'il faut protéger les enfants. »
\item \esp{Me parece que este texto es muy importante porque defiende la igualdad.} « Il me semble que ce texte est très important parce qu'il défend l'égalité. »
\item \esp{Es necesario que la sociedad luche contra el trabajo infantil.} « Il est nécessaire que la société lutte contre le travail des enfants. »
\item \esp{Tenemos que escolarizar a todas las niñas.} « Nous devons scolariser toutes les filles. »
\end{itemize}
\end{exemplebox}

\begin{attention}[Creo que : deux pièges]
\begin{itemize}
\item Après \esp{creo que} (affirmatif), on utilise l'\textbf{indicatif} : \esp{Creo que todos los niños \textbf{tienen} derecho a la educación.} Mais à la forme négative, le subjonctif s'impose : \esp{No creo que la situación \textbf{sea} fácil.} « Je ne pense pas que la situation soit facile. »
\item Ne dites pas \esp{creo de que} : le verbe \esp{creer} se construit directement avec \esp{que}.
\end{itemize}
\end{attention}

\courssub{La méthode résumée}

\begin{methode}[Rédiger un commentaire de texte au Bac]
\begin{enumerate}
\item \textbf{Présenter le texte} : nature (article, discours...), source, thème, en deux ou trois phrases. Annoncer le plan.
\item \textbf{Dégager les idées} : deux ou trois idées principales, chacune illustrée par une citation courte entre guillemets, reliées par des connecteurs (\esp{primero}, \esp{además}, \esp{sin embargo}).
\item \textbf{Donner son avis personnel} : bilan bref, puis opinion argumentée avec \esp{creo que}, \esp{me parece que}, \esp{es necesario que}.
\end{enumerate}
Rédigez entièrement en espagnol, avec des phrases simples et correctes : une phrase courte et juste vaut mieux qu'une phrase longue et fautive.
\end{methode}

\courssub{Commentaire modèle lu et analysé en classe}

Voici un commentaire modèle rédigé sur le texte support de la leçon (\emph{La mujer y los derechos del niño}). Lisez-le attentivement : nous allons ensuite y repérer les trois parties.

\begin{textebox}[Comentario de texto — modelo]
Este texto es un artículo de prensa sobre la condición de la mujer y los derechos del niño. Primero, el autor denuncia la discriminación contra la mujer: en algunas regiones, «las niñas no van a la escuela». Además, explica los problemas graves que sufren los menores, como el trabajo infantil y el maltrato. Finalmente, presenta la Convención de 1989, que garantiza su protección. En conclusión, me parece que este texto es fundamental porque creo que todos los niños tienen derecho a la educación y es necesario que la sociedad luche por la igualdad real.
\end{textebox}

\textbf{Traduction du modèle :} « Ce texte est un article de presse sur la condition de la femme et les droits de l'enfant. D'abord, l'auteur dénonce la discrimination contre la femme : dans certaines régions, "les filles ne vont pas à l'école". De plus, il explique les problèmes graves que subissent les mineurs, comme le travail des enfants et la maltraitance. Finalement, il présente la Convention de 1989, qui garantit leur protection. En conclusion, il me semble que ce texte est fondamental parce que je pense que tous les enfants ont droit à l'éducation et qu'il est nécessaire que la société lutte pour une égalité réelle. »

\begin{aretenir}[Repérage dans le modèle]
\begin{itemize}
\item \textbf{Introduction} (phrase 1) : nature du texte (\esp{artículo de prensa}) + thème (\esp{la condición de la mujer y los derechos del niño}).
\item \textbf{Idée 1} (phrase 2) : la discrimination contre la femme, avec la citation courte « \esp{las niñas no van a la escuela} ». Connecteur : \esp{Primero}.
\item \textbf{Idée 2} (phrase 3) : les problèmes des enfants (\esp{trabajo infantil}, \esp{maltrato}). Connecteur : \esp{Además}.
\item \textbf{Idée 3} (phrase 4) : la solution, la Convention de 1989. Connecteur : \esp{Finalmente}.
\item \textbf{Conclusion} (phrase 5) : avis personnel avec \esp{me parece que}, \esp{creo que} + indicatif (\esp{tienen}) et \esp{es necesario que} + subjonctif (\esp{luche}).
\item \textbf{Connecteurs repérés} : \esp{primero}, \esp{además}, \esp{finalmente}, \esp{en conclusión}, \esp{porque}.
\end{itemize}
\end{aretenir}

\courssub{Entraînement guidé}

\begin{exoresolu}[Compléter un commentaire]
Énoncé — Complétez ce commentaire avec les connecteurs et les structures suivants : \esp{además}, \esp{en conclusión}, \esp{según}, \esp{creo que}, \esp{sin embargo}.

\esp{Este texto es un artículo sobre los derechos del niño. (1) \ldots\ el autor, la educación es un derecho de todos. (2) \ldots , denuncia el trabajo infantil en las ciudades. (3) \ldots , muchos niños todavía no van a la escuela. (4) \ldots , (5) \ldots\ hay que proteger a los menores.}

\tcblower
Solution détaillée : (1) \esp{Según} — « selon l'auteur » (devant un nom, on utilise \esp{según}). (2) \esp{Además} — ajout d'une deuxième idée. (3) \esp{Sin embargo} — opposition : malgré les droits proclamés, la réalité est différente. (4) \esp{En conclusión} — ouverture de la conclusion. (5) \esp{creo que} — avis personnel suivi de l'indicatif (\esp{hay que}).

Texte complet : \esp{Este texto es un artículo sobre los derechos del niño. Según el autor, la educación es un derecho de todos. Además, denuncia el trabajo infantil en las ciudades. Sin embargo, muchos niños todavía no van a la escuela. En conclusión, creo que hay que proteger a los menores.} « Ce texte est un article sur les droits de l'enfant. Selon l'auteur, l'éducation est un droit de tous. De plus, il dénonce le travail des enfants dans les villes. Cependant, beaucoup d'enfants ne vont pas encore à l'école. En conclusion, je pense qu'il faut protéger les mineurs. »
\end{exoresolu}

\begin{exercice}[À vous de commenter]
À partir du texte support de la leçon, rédigez un commentaire de cinq à huit phrases en suivant le plan : 1) présentez le texte ; 2) dégagez deux idées avec une citation courte chacune ; 3) concluez avec \esp{creo que} et \esp{es necesario que}. Utilisez au moins trois connecteurs de la liste : \esp{primero}, \esp{además}, \esp{sin embargo}, \esp{por eso}, \esp{en conclusión}.
\end{exercice}

\begin{corrige}[Exercice — éléments de correction]
Exemple de production attendue : \esp{Este texto es un artículo de prensa sobre la mujer y los niños. Primero, el autor explica que «la educación es un derecho» para todos. Además, denuncia el maltrato y el trabajo infantil. Sin embargo, la Convención de 1989 protege a los menores. En conclusión, creo que este texto es muy útil y es necesario que todos los gobiernos respeten los derechos del niño.} Barème indicatif : présentation correcte (nature + thème), deux idées avec citations courtes entre guillemets, connecteurs variés, avis personnel avec \esp{creo que} + indicatif et \esp{es necesario que} + subjonctif (\esp{respeten}), espagnol simple et correct.
\end{corrige}

\begin{aretenir}[L'essentiel de la section]
Un commentaire de texte se construit en trois blocs : \textbf{introducción} (nature, thème), \textbf{desarrollo} (deux ou trois idées + citations courtes + connecteurs), \textbf{conclusión} (bilan + avis avec \esp{creo que} / \esp{es necesario que}). On cite des fragments brefs entre guillemets, jamais de longs passages. On rédige en espagnol, simplement et correctement.
\end{aretenir}

\coursec{Production guidée : opinar y defender los derechos}

Après avoir appris à lire et à commenter un texte sur la condition de la femme et les droits de l'enfant, il est temps de passer à la \cle{production} : c'est vous maintenant qui allez écrire et parler. Dans cette dernière section, vous allez rédiger un paragraphe d'opinion sur la scolarisation des filles au Cameroun, débattre à l'oral avec un camarade, et découvrir un genre très fréquent à l'épreuve d'expression écrite du Baccalauréat : l'article ou l'affiche de sensibilisation. Toutes les structures étudiées dans cette leçon (\esp{creo que}, \esp{es necesario que}, \esp{tener que} + infinitif) et tout le lexique des droits vont vous servir d'outils.

\courssub{Observer avant de produire : un paragraphe modèle}

Avant d'écrire, observons ensemble un paragraphe d'opinion réussi. Lisez-le attentivement et repérez : la phrase d'introduction, les deux arguments, les connecteurs et la conclusion.

\begin{textebox}[La escolarización de las niñas en Camerún — párrafo modelo]
Hoy en día, la escolarización de las niñas es un tema muy importante en Camerún. En efecto, muchas niñas no van a la escuela porque sus familias son pobres o porque deben ayudar en la casa. En primer lugar, la educación es un derecho fundamental de todos los niños y de todas las niñas: una niña educada puede encontrar un buen trabajo y ayudar a su familia. Además, la escuela protege a las niñas contra el matrimonio precoz y el trabajo infantil, dos problemas graves en algunas regiones del país. Sin embargo, todavía existen tradiciones que impiden la igualdad entre los niños y las niñas. A mi juicio, el gobierno tiene que luchar contra estas tradiciones y apoyar a las familias pobres. En conclusión, creo que la escolarización de las niñas es necesaria para el desarrollo de Camerún, y es fundamental que todos los padres envíen a sus hijas a la escuela.
\end{textebox}

\textbf{Glossaire :} \esp{la escolarización} = la scolarisation ; \esp{en efecto} = en effet ; \esp{el matrimonio precoz} = le mariage précoce ; \esp{impedir} (e $\to$ i) = empêcher ; \esp{apoyar} = soutenir ; \esp{enviar} = envoyer ; \esp{el desarrollo} = le développement ; \esp{todavía} = encore, toujours.

\textbf{Questions d'observation :}
\begin{enumerate}
\item Quelle phrase introduit le sujet ?
\item Quels sont les deux arguments de l'auteur ? Quels connecteurs les introduisent ?
\item Quelle expression marque l'opposition ?
\item Quelles structures de la leçon retrouvez-vous dans la conclusion ?
\end{enumerate}

\textbf{Réponses :} 1. \esp{Hoy en día, la escolarización de las niñas es un tema muy importante en Camerún.} 2. Argument 1 : l'éducation est un droit fondamental (\esp{En primer lugar}) ; argument 2 : l'école protège contre le mariage précoce et le travail des enfants (\esp{Además}). 3. \esp{Sin embargo} (cependant). 4. \esp{creo que} + indicatif, \esp{es fundamental que} + subjonctif (\esp{envíen}), \esp{tiene que} + infinitif (\esp{luchar}).

\courssub{La banque d'expressions de l'opinion et de l'obligation}

Voici votre boîte à outils. Apprenez ces expressions par cœur : elles sont la clé de toute production d'opinion au lycée.

\begin{center}
\begin{tabularx}{\linewidth}{|X|X|X|}
\hline
\rowcolor{popblueL} Español & Français & Structure \\
\hline
\esp{A mi juicio, ...} & À mon avis, ... & + indicatif \\
\hline
\esp{Creo que ...} & Je pense que ... & + indicatif \\
\hline
\esp{No creo que ...} & Je ne pense pas que ... & + subjonctif \\
\hline
\esp{Estoy de acuerdo con ...} & Je suis d'accord avec ... & + nom ou infinitif \\
\hline
\esp{Es necesario que ...} & Il est nécessaire que ... & + subjonctif \\
\hline
\esp{Es fundamental que ...} & Il est fondamental que ... & + subjonctif \\
\hline
\esp{Hay que ...} & Il faut ... & + infinitif \\
\hline
\esp{Tener que ...} & Devoir, être obligé de ... & + infinitif \\
\hline
\esp{Deber ...} & Devoir (conseil, devoir moral) & + infinitif \\
\hline
\end{tabularx}
\end{center}

\begin{exemplebox}[Exemples traduits]
\esp{A mi juicio, el matrimonio precoz es un problema grave.} = À mon avis, le mariage précoce est un problème grave.

\esp{Es fundamental que el gobierno proteja a los niños.} = Il est fondamental que le gouvernement protège les enfants. (subjonctif \esp{proteja})

\esp{Hay que respetar los derechos del niño.} = Il faut respecter les droits de l'enfant.

\esp{Los padres tienen que enviar a sus hijas a la escuela.} = Les parents doivent envoyer leurs filles à l'école.

\esp{Las niñas deben ir a la escuela como los niños.} = Les filles doivent aller à l'école comme les garçons.

\esp{Estoy de acuerdo con la igualdad entre hombres y mujeres.} = Je suis d'accord avec l'égalité entre les hommes et les femmes.
\end{exemplebox}

\begin{propriete}[Nuancer l'obligation : tener que, deber, hay que]
\begin{itemize}
\item \esp{Tener que} + infinitif : obligation concrète, personnelle, avec un sujet précis. \esp{Yo tengo que estudiar.} (Je dois étudier.)
\item \esp{Deber} + infinitif : devoir moral, conseil. \esp{Debemos respetar a las mujeres.} (Nous devons respecter les femmes.)
\item \esp{Hay que} + infinitif : obligation générale, impersonnelle, sans sujet. \esp{Hay que proteger a los niños.} (Il faut protéger les enfants.)
\end{itemize}
Attention : \esp{hay que} ne change jamais de forme ; on ne dit jamais « han que » ni « hay que hacerlo nosotros » avec un sujet exprimé.
\end{propriete}

\begin{attention}[Accorder l'adjectif avec le nom]
En espagnol comme en français, l'adjectif s'accorde en genre et en nombre avec le nom : \esp{los derechos humanos} (les droits humains), \esp{la igualdad social} (l'égalité sociale), \esp{las niñas educadas} (les filles éduquées). Mais attention : \esp{infantil} est \cle{invariable en genre} ! On dit \esp{el trabajo infantil} (le travail des enfants) et \esp{la protección infantil} (la protection de l'enfance). N'écrivez jamais « la protección infantila ». Autres adjectifs invariables en genre (mais qui prennent un -s ou -es au pluriel) : \esp{débil} (\esp{un niño débil}, \esp{una niña débil}), \esp{azul} (\esp{los ojos azules}), \esp{joven} (\esp{un hombre joven}, \esp{una mujer joven}, \esp{las leyes jóvenes}).
\end{attention}

\courssub{Production écrite guidée : le plan imposé}

Vous allez rédiger un paragraphe d'opinion de \cle{8 à 10 lignes} sur le thème : \esp{La escolarización de las niñas en Camerún}. Le plan est imposé :

\begin{enumerate}
\item \textbf{Une phrase d'introduction} : présentez le sujet (\esp{Hoy en día...}, \esp{En Camerún, ...}).
\item \textbf{Deux arguments avec connecteurs} : \esp{En primer lugar...} puis \esp{Además...} ou \esp{En segundo lugar...}
\item \textbf{Une conclusion avec} \esp{creo que} \textbf{ou} \esp{es necesario que} : donnez votre opinion personnelle.
\end{enumerate}

\begin{popfigure}
\begin{tikzpicture}[
  node distance=0.55cm,
  box/.style={draw=popblue, fill=popblueL, rounded corners=3pt, align=center, text width=6.2cm, inner sep=5pt, font=\small},
  arr/.style={-{Stealth[length=2.5mm]}, thick, poporange}
]
\node[box, align=center] (tema){\textbf{1. Tema / introducción}\\ \esp{Hoy en día, la escolarización de las niñas es...}};
\node[box, below=of tema, align=center] (arg1){\textbf{2. Argumento 1}\\ \esp{En primer lugar, la educación es un derecho...}};
\node[box, below=of arg1, align=center] (arg2){\textbf{3. Argumento 2}\\ \esp{Además, la escuela protege contra...}};
\node[box, below=of arg2, align=center] (op){\textbf{4. Opinión personal}\\ \esp{A mi juicio, ... / Estoy de acuerdo con...}};
\node[box, below=of op, fill=popgreenL, draw=popgreen, align=center] (conc){\textbf{5. Conclusión}\\ \esp{Creo que... / Es necesario que...}};
\draw[arr] (tema) -- (arg1);
\draw[arr] (arg1) -- (arg2);
\draw[arr] (arg2) -- (op);
\draw[arr] (op) -- (conc);
\end{tikzpicture}
\legende{Organigramme de la production guidée : du thème à la conclusion.}
\end{popfigure}

\begin{methode}[Réussir l'expression écrite au Bac]
\begin{enumerate}
\item \textbf{Respectez le nombre de mots} demandé (ici 8 à 10 lignes) : un texte trop court ou trop long est pénalisé.
\item \textbf{Utilisez au moins trois structures de la leçon} : \esp{creo que}, \esp{es necesario que} + subjonctif, \esp{tener que} + infinitif.
\item \textbf{Employez au moins cinq mots du lexique des droits} : \esp{igualdad}, \esp{derechos}, \esp{escolarización}, \esp{trabajo infantil}, \esp{maltrato}, \esp{protección}.
\item \textbf{Soignez les connecteurs} : \esp{en primer lugar}, \esp{además}, \esp{sin embargo}, \esp{en conclusión}.
\item \textbf{Relisez les accords} (article, nom, adjectif) \textbf{et les conjugaisons}, surtout le subjonctif après \esp{es necesario que}.
\end{enumerate}
\end{methode}

\begin{exoresolu}[Rédaction guidée : compléter le paragraphe]
Énoncé : complétez ce paragraphe avec les expressions de la banque : \esp{a mi juicio}, \esp{es necesario que}, \esp{tienen que}, \esp{hay que}, \esp{creo que}.

\esp{En Camerún, muchas niñas no van a la escuela. En primer lugar, la educación es un derecho fundamental. Además, la escuela protege a las niñas contra el matrimonio precoz. (1) \dots, los padres (2) \dots{} enviar a sus hijas a la escuela. (3) \dots{} luchar contra el trabajo infantil. (4) \dots{} la igualdad es posible, y (5) \dots{} el gobierno proteja a todos los niños.}

\tcblower

Solution détaillée :

\esp{En Camerún, muchas niñas no van a la escuela. En primer lugar, la educación es un derecho fundamental. Además, la escuela protege a las niñas contra el matrimonio precoz. \textbf{A mi juicio}, los padres \textbf{tienen que} enviar a sus hijas a la escuela. \textbf{Hay que} luchar contra el trabajo infantil. \textbf{Creo que} la igualdad es posible, y \textbf{es necesario que} el gobierno proteja a todos los niños.}

Remarques : (2) \esp{tienen que} + infinitif exprime l'obligation des parents, avec un sujet précis (\esp{los padres}) ; (3) \esp{hay que} + infinitif exprime une obligation générale, impersonnelle, sans sujet ; (5) après \esp{es necesario que}, le subjonctif \esp{proteja} est obligatoire — et non l'indicatif « protege ». Rappel de la formation : \esp{proteger} $\to$ radical \esp{protej-} + terminaisons du subjonctif présent des verbes en -er : \esp{proteja, protejas, proteja, protejamos, protejáis, protejan}.
\end{exoresolu}

\courssub{Production orale : el debate en parejas}

\begin{experience}[Débat : ¿Deben las niñas ir a la escuela como los niños?]
Travaillez en binômes. L'élève A défend la position « pour » (\esp{a favor}), l'élève B joue le rôle « contre » (\esp{en contra}) — même si ce n'est pas sa vraie opinion : c'est un exercice de rhétorique !

\textbf{Tour 1} — A commence : \esp{A mi juicio, las niñas deben ir a la escuela porque...} (deux arguments, une minute).

\textbf{Tour 2} — B répond : \esp{No estoy de acuerdo. Creo que... porque...} (une minute).

\textbf{Tour 3} — Réplique libre : chacun répond à l'argument de l'autre avec \esp{Sin embargo...}, \esp{Es verdad que..., pero...}.

\textbf{Tour 4} — Conclusion commune : trouvez ensemble un accord avec \esp{Estamos de acuerdo en que...} ou \esp{Es necesario que...}.

Expressions utiles pour le débat : \esp{¿Qué piensas tú?} (qu'en penses-tu ?), \esp{Tienes razón} (tu as raison), \esp{No tienes razón} (tu n'as pas raison), \esp{Por ejemplo...} (par exemple).
\end{experience}

Pendant le débat, le professeur circule et note les erreurs. Vient ensuite la \cle{correction collective} au tableau. Voici les deux familles d'erreurs à relever :

\begin{attention}[Erreurs fréquentes à la correction collective]
\textbf{1. Le subjonctif oublié.} Après \esp{es necesario que}, \esp{es fundamental que}, \esp{es importante que}, le subjonctif est obligatoire : on dit \esp{Es necesario que las niñas \textbf{vayan} a la escuela} et non « que las niñas van ». De même : \esp{No creo que sea justo} (et non « no creo que es justo »). Rappel : \esp{ir} au subjonctif présent : \esp{vaya, vayas, vaya, vayamos, vayáis, vayan} ; \esp{ser} : \esp{sea, seas, sea, seamos, seáis, sean}.

\textbf{2. Les erreurs de lexique.} Attention aux faux amis et aux calques du français : « l'éducation » se dit \esp{la educación} (avec un accent sur la ó) ; « une école » = \esp{una escuela} ou \esp{un colegio} ; « les droits » = \esp{los derechos} (et non « las derechas », qui signifie « les droites », la direction ou la droite politique !) ; « défendre » = \esp{defender} (e $\to$ ie : \esp{defiendo}) ; « protéger » = \esp{proteger} (avec un seul « t »).
\end{attention}

\courssub{Complément utile au Bac : l'article et l'affiche de sensibilisation}

\begin{aretenir}[Un genre fréquent à l'épreuve d'expression écrite]
Au Baccalauréat, on vous demande souvent de rédiger un \cle{court article} pour le journal du lycée ou une \cle{affiche de sensibilisation} (\esp{un cartel de sensibilización}) sur un sujet de société : les droits de l'enfant, l'égalité, la santé, l'environnement. Pour l'affiche, soyez bref et percutant : un titre à l'impératif ou à l'infinitif, deux ou trois slogans, une phrase d'appel à l'action. N'oubliez pas les signes d'ouverture ¡ et ¿, typiques de l'espagnol et appréciés du correcteur.
\end{aretenir}

\begin{exemplebox}[Modèle d'affiche de sensibilisation]
\esp{¡EDUCAR A UNA NIÑA ES EDUCAR A UNA NACIÓN!}

\esp{La escuela es un derecho, no un privilegio.}

\esp{Di NO al matrimonio precoz. Di NO al trabajo infantil.}

\esp{Hay que proteger a nuestros niños. ¡Todos a la escuela!}

(« Éduquer une fille, c'est éduquer une nation ! L'école est un droit, pas un privilège. Dis NON au mariage précoce. Dis NON au travail des enfants. Il faut protéger nos enfants. Tous à l'école ! »)

Remarque : \esp{di} est l'impératif de \esp{decir} (dis) ; \esp{al} = \esp{a} + \esp{el}.
\end{exemplebox}

\begin{exercice}[À vous de jouer]
\begin{enumerate}
\item Rédigez votre paragraphe d'opinion (8 à 10 lignes) sur \esp{La escolarización de las niñas en Camerún} en suivant le plan imposé et la méthode du Bac.
\item En binômes, jouez le débat devant la classe ; les autres élèves relèvent les erreurs de subjonctif et de lexique.
\item En groupes de quatre, créez une affiche de sensibilisation contre le travail des enfants avec au moins trois slogans utilisant \esp{hay que}, \esp{es necesario que} et l'impératif.
\end{enumerate}
\end{exercice}

\begin{corrige}[Exercice « À vous de jouer » — pistes]
\textbf{1.} Paragraphe possible : \esp{Hoy en día, la escolarización de las niñas es un problema importante en mi país. En primer lugar, la educación es un derecho de todos los niños. Además, una niña educada puede ayudar a su familia y a su país. Sin embargo, muchas familias pobres no pueden pagar la escuela. A mi juicio, el gobierno tiene que ayudar a estas familias. En conclusión, creo que la igualdad es necesaria, y es fundamental que todas las niñas vayan a la escuela.}

\textbf{3.} Slogans possibles : \esp{¡Hay que decir NO al trabajo infantil!} ; \esp{Es necesario que los niños estudien, no que trabajen.} ; \esp{¡Protege a los niños! ¡Respeta sus derechos!}
\end{corrige}

\begin{savaistu}[Le Cameroun défend les droits des femmes et des enfants]
Au Cameroun, le ministère de la Promotion de la Femme et de la Famille (MINPROFF), créé en 1984, mène avec de nombreuses ONG des campagnes contre le mariage précoce et pour la scolarisation des filles, en particulier dans les régions du Nord. Le Cameroun a ratifié la Convention relative aux droits de l'enfant en 1993. Dans le monde hispanique, des pays comme l'Espagne, le Mexique ou l'Argentine ont aussi adopté des lois pour l'égalité (\esp{la igualdad}) entre les femmes et les hommes. La Guinée équatoriale, voisine hispanophone du Cameroun, célèbre chaque année la Journée internationale de la femme, le 8 mars (\esp{el Día Internacional de la Mujer}), et la Journée de l'enfant africain, le 16 juin, est commémorée dans tout le continent, y compris au Cameroun.
\end{savaistu}

\newpage

\coursec{Activité d'intégration}

\courssub{Objectifs de l'activité}

À la fin de cette activité d'intégration, l'élève sera capable de :
\begin{itemize}
\item mobiliser le \cle{lexique des droits} (\esp{igualdad, derechos, escolarización, trabajo infantil, maltrato, protección}) dans une production authentique ;
\item employer correctement les structures de l'\cle{opinion} (\esp{creo que + indicatif}), de la \cle{nécessité} (\esp{es necesario que + subjonctif}) et de l'\cle{obligation} (\esp{tener que + infinitif}) ;
\item lire et comprendre un court document de sensibilisation en espagnol ;
\item créer en groupe une \cle{affiche de campagne} et la présenter oralement en deux minutes ;
\item justifier un choix à l'oral en espagnol simple.
\end{itemize}

\courssub{Situation}

Le club d'espagnol de ton lycée à Yaoundé participe à la \emph{Semana de los Derechos del Niño}. La principale du collège bilingue voisin a lancé un concours : chaque club doit créer une affiche ou un court article de campagne intitulé \esp{« Todos los niños a la escuela »}. La meilleure production sera publiée dans le journal du quartier et affichée à l'entrée de l'école primaire du secteur, où certaines filles ne sont pas scolarisées et où des enfants vendent des marchandises au marché pendant les heures de classe. Voici le texte de l'appel à projets affiché au tableau d'information du lycée :

\begin{textebox}[Convocatoria — Semana de los Derechos del Niño]
\textbf{CAMPAÑA : « TODOS LOS NIÑOS A LA ESCUELA »}\par
El club de español organiza un concurso de carteles contra el trabajo infantil y por la escolarización de las niñas.\par
En nuestro barrio, algunos niños venden frutas en el mercado durante las horas de clase y muchas niñas no van a la escuela. Esto es una violación de sus derechos.\par
Tu cartel debe contener :\par
— un eslogan original ;\par
— tres frases con \emph{creo que}, \emph{es necesario que} y \emph{tener que} ;\par
— seis palabras del vocabulario de los derechos : \emph{igualdad, derechos, escolarización, trabajo infantil, maltrato, protección}.\par
Cada grupo presenta su cartel en dos minutos. ¡Participa y defiende los derechos de los niños !
\end{textebox}

\textbf{Glossaire :} \esp{convocatoria} = appel à projets ; \esp{cartel} = affiche ; \esp{eslogan} = slogan ; \esp{barrio} = quartier ; \esp{violación} = violation ; \esp{¡Participa!} = Participe !

\courssub{Consignes}

Par groupes de quatre, réalisez les tâches suivantes dans l'ordre :

\begin{enumerate}
\item \textbf{Compréhension.} Lisez la \esp{convocatoria} et répondez par écrit en espagnol : ¿Qué problema hay en el barrio ? ¿Qué debe contener el cartel ? (Deux phrases complètes suffisent.)
\item \textbf{Lexique.} Soulignez dans le texte les mots du vocabulaire des droits, puis complétez ce mini-lexique avec la traduction française : \esp{igualdad}, \esp{escolarización}, \esp{trabajo infantil}, \esp{maltrato}, \esp{protección}, \esp{derechos}.
\item \textbf{Grammaire.} Transformez les phrases suivantes : (a) \esp{Los niños deben ir a la escuela.} $\rightarrow$ avec \esp{tener que} ; (b) \esp{Es importante proteger a los niños.} $\rightarrow$ avec \esp{es necesario que + subjonctif} ; (c) \esp{En mi opinión, la educación es un derecho.} $\rightarrow$ avec \esp{creo que}.
\item \textbf{Production écrite.} Créez votre affiche : un slogan accrocheur (avec ¡ !), trois phrases utilisant les trois structures imposées, et au moins six mots du lexique de la leçon. Répartissez les rôles : un secrétaire, un correcteur linguistique, un dessinateur, un porte-parole.
\item \textbf{Production orale.} Présentez votre affiche à la classe en deux minutes : le porte-parole lit le slogan et les phrases, puis explique le message en deux ou trois phrases simples.
\item \textbf{Vote et justification.} La classe vote pour la meilleure campagne. Chaque électeur justifie son choix avec une phrase du type : \esp{Voto por el grupo dos porque su eslogan es claro y defiende la igualdad.}
\end{enumerate}

\courssub{Ressources à mobiliser}

\begin{aretenir}[Structures de l'opinion, de la nécessité et de l'obligation]
\begin{itemize}
\item \textbf{Opinion :} \esp{creo que / pienso que + indicatif}. \esp{Creo que la educación es un derecho.} « Je pense que l'éducation est un droit. »
\item \textbf{Nécessité :} \esp{es necesario que + subjonctif}. \esp{Es necesario que todos los niños vayan a la escuela.} « Il faut que tous les enfants aillent à l'école. »
\item \textbf{Obligation :} \esp{tener que + infinitif}. \esp{Los padres tienen que escolarizar a sus hijas.} « Les parents doivent scolariser leurs filles. »
\end{itemize}
\end{aretenir}

\begin{center}
\begin{tabularx}{\linewidth}{|X|X|X|}
\hline
\rowcolor{popblueL} Español & Français & Exemple \\
\hline
la igualdad & l'égalité & \esp{la igualdad entre niños y niñas} \\
\hline
los derechos & les droits & \esp{los derechos del niño} \\
\hline
la escolarización & la scolarisation & \esp{la escolarización de las niñas} \\
\hline
el trabajo infantil & le travail des enfants & \esp{contra el trabajo infantil} \\
\hline
el maltrato & la maltraitance & \esp{decir no al maltrato} \\
\hline
la protección & la protection & \esp{la protección de los menores} \\
\hline
\end{tabularx}
\end{center}

\begin{attention}[Pièges à éviter]
\begin{itemize}
\item Après \esp{es necesario que}, n'oublie pas le \textbf{subjonctif} : \esp{es necesario que los niños \textbf{vayan}}, et non \esp{van}.
\item \esp{Tener que} est toujours suivi de l'\textbf{infinitif} : \esp{tienen que estudiar}, jamais de conjuguer le second verbe.
\item Faux ami : \esp{la educación} signifie « l'éducation » mais aussi « la politesse » ; pour « l'enseignement », on peut dire \esp{la enseñanza}.
\item Le slogan prend les signes doubles : \esp{¡Todos los niños a la escuela!}
\end{itemize}
\end{attention}

\courssub{Proposition de production et corrigé commenté}

\begin{exoresolu}[Affiche modèle du groupe 3 — Lycée de Yaoundé]
Énoncé : rédigez l'affiche complète de la campagne \esp{« Todos los niños a la escuela »} en respectant les trois contraintes de la convocatoria.
\tcblower
\textbf{Production modèle :}\par
\esp{¡TODOS LOS NIÑOS TIENEN DERECHO A LA EDUCACIÓN!}\par
\esp{Creo que la escolarización de las niñas es tan importante como la de los niños : la igualdad empieza en la escuela. Es necesario que la sociedad proteja a los menores y que luche contra el trabajo infantil y el maltrato. Los padres tienen que enviar a sus hijos a la escuela, y no al mercado. La protección de los derechos del niño es tarea de todos. ¡Un niño en la escuela es un futuro para el Camerún !}\par
\medskip
\textbf{Traduction :} « Tous les enfants ont droit à l'éducation ! Je pense que la scolarisation des filles est aussi importante que celle des garçons : l'égalité commence à l'école. Il faut que la société protège les mineurs et qu'elle lutte contre le travail des enfants et la maltraitance. Les parents doivent envoyer leurs enfants à l'école, et non au marché. La protection des droits de l'enfant est l'affaire de tous. Un enfant à l'école, c'est un avenir pour le Cameroun ! »\par
\medskip
\textbf{Commentaire des choix de langue :}
\begin{itemize}
\item Le \textbf{slogan} est court, en majuscules, avec les signes ¡ ! et le verbe \esp{tener derecho a + nom}, structure idiomatique pour parler des droits.
\item \esp{Creo que} est suivi de l'\textbf{indicatif} (\esp{es}) : c'est une opinion affirmative, donc pas de subjonctif.
\item \esp{Es necesario que} est suivi du \textbf{subjonctif présent} : \esp{proteja} (de \esp{proteger}) et \esp{luche} (de \esp{luchar}). Attention aux deux sujets différents : \esp{la sociedad} est le sujet des deux verbes, reliés par \esp{y que}.
\item \esp{Tienen que enviar} : \esp{tener que + infinitif} exprime l'obligation qui pèse sur les parents.
\item Les \textbf{six mots du lexique} sont bien présents : \esp{escolarización, igualdad, protección / proteja, trabajo infantil, maltrato, derechos} — la contrainte est même dépassée, ce qui enrichit le texte.
\item La chute \esp{¡Un niño en la escuela es un futuro para el Camerún!} ancre la campagne dans le contexte camerounais et donne une conclusion mémorable.
\end{itemize}
\end{exoresolu}

\courssub{Bilan de l'activité}

Cette activité t'a permis de passer de la lecture à l'action : tu as compris un document authentique de sensibilisation, réutilisé le lexique des droits dans un contexte réel et produit un message collectif en espagnol. Retiens l'essentiel :

\begin{itemize}
\item \textbf{Lexique :} tu sais maintenant nommer les grandes notions des droits de l'enfant et de la condition de la femme en espagnol ;
\item \textbf{Grammaire :} tu distingues les trois structures — opinion + indicatif, nécessité + subjonctif, obligation + infinitif — et tu sais les employer sans les confondre ;
\item \textbf{Communication :} tu es capable de défendre une cause à l'écrit (affiche, slogan) et à l'oral (présentation, vote justifié), compétence directement réinvestissable dans l'épreuve d'expression du baccalauréat.
\end{itemize}

Pour aller plus loin, tu peux transformer ton affiche en court article pour le journal du lycée, en ajoutant un paragraphe sur la situation des filles dans ta région, toujours avec les structures de la leçon.

\newpage

\coursec{Méthodes et exercices résolus}

\coursec{EA02 : El estatuto de la mujer y los derechos del niño}

\courssub{Texte support : La mujer y los derechos del niño}

\begin{textebox}[Texte original — Niveau A2+/B1]
En mi barrio de Yaoundé, muchas niñas venden agua en el mercado en lugar de ir a la escuela. Sus familias son pobres y necesitan dinero para sobrevivir. Sin embargo, la educación es un derecho fundamental de todos los niños, sin distinción de sexo. La Convención sobre los Derechos del Niño, firmada por casi todos los países del mundo, incluido Camerún y Guinea Ecuatorial, garantiza el derecho a la escolarización y la protección contra el trabajo infantil y el maltrato.

Cuando una niña va a la escuela, aprende a leer, a escribir y a defender sus derechos. Luego, toda su familia sale ganando: ella tendrá un trabajo mejor, cuidará mejor de su salud y de la de sus futuros hijos. Por eso, el gobierno y los padres deben trabajar juntos. Es necesario que las autoridades construyan más escuelas y que los padres comprendan que el lugar de un niño es en el aula, no en la calle.
\end{textebox}

\textbf{Traduction / Glossaire :} « Dans mon quartier de Yaoundé, beaucoup de filles vendent de l'eau au marché au lieu d'aller à l'école. Leurs familles sont pauvres et ont besoin d'argent pour survivre. Cependant, l'éducation est un droit fondamental de tous les enfants, sans distinction de sexe. La Convention relative aux droits de l'enfant, signée par presque tous les pays du monde, y compris le Cameroun et la Guinée équatoriale, garantit le droit à la scolarisation et la protection contre le travail des enfants et la maltraitance. Quand une fille va à l'école, elle apprend à lire, à écrire et à défendre ses droits. Ensuite, toute sa famille y gagne : elle aura un meilleur travail, prendra mieux soin de sa santé et de celle de ses futurs enfants. C'est pourquoi le gouvernement et les parents doivent travailler ensemble. Il est nécessaire que les autorités construisent plus d'écoles et que les parents comprennent que la place d'un enfant est dans la salle de classe, pas dans la rue. »

\cle{el barrio} (quartier), \cle{en lugar de} (au lieu de), \cle{sobrevivir} (survivre), \cle{la Convención} (la Convention), \cle{garantiza} (garantit), \cle{el trabajo infantil} (le travail des enfants), \cle{el maltrato} (la maltraitance), \cle{defender} (défendre), \cle{sale ganando} (y gagne / s'en sort gagnant), \cle{el aula} (la salle de classe).

\courssub{Compréhension du texte}

\begin{methode}[Réussir les questions de compréhension]
\begin{enumerate}
\item \textbf{Première lecture globale.} Lis le texte une fois sans dictionnaire. Repère le thème (\esp{la mujer}, \esp{los derechos del niño}), le type de texte (témoignage / article d'opinion) et la thèse défendue.
\item \textbf{Lecture analytique.} Souligne les mots-clés du lexique des droits : \esp{igualdad}, \esp{derechos}, \esp{escolarización}, \esp{trabajo infantil}, \esp{maltrato}, \esp{protección}, \cle{la Convención}. Ils portent presque toujours les réponses.
\item \textbf{Lire les questions avant de répondre.} Repère la consigne : \esp{¿Verdadero o falso?} (justifie avec une citation du texte), \esp{¿Por qué...?} (réponds avec \esp{porque} + phrase complète), \esp{¿Qué significa...?} (reformule avec tes mots).
\item \textbf{Répondre en phrases complètes}, en espagnol, en recopiant la structure de la question : \esp{¿Por qué trabajan muchas niñas?} $\rightarrow$ \esp{Muchas niñas trabajan porque sus familias son pobres.}
\item \textbf{Citer le texte entre guillemets} quand on demande une justification : \esp{El texto dice: «...»}.
\item \textbf{Ne jamais répondre en français}, sauf si la question est posée en français et demande explicitement une traduction.
\end{enumerate}
\end{methode}

\begin{exoresolu}[Compréhension de texte : type Bac]
\textbf{Texte} : \esp{En muchos países, las mujeres todavía no tienen los mismos derechos que los hombres: ganan menos por el mismo trabajo y tienen menos acceso a la educación. Además, millones de niños trabajan en lugar de ir a la escuela. La escolarización de las niñas es una prioridad, porque educar a una niña es educar a toda una familia.}

\textbf{Énoncé} : Réponds en espagnol, par des phrases complètes.
\begin{enumerate}
\item \esp{¿Tienen las mujeres los mismos derechos que los hombres?}
\item \esp{¿Por qué trabajan muchos niños?}
\item \esp{¿Por qué es una prioridad la escolarización de las niñas?}
\end{enumerate}
\tcblower
\textbf{Raisonnement.} Chaque question contient déjà la moitié de la réponse : il suffit de reprendre sa structure et de compléter avec l'information du texte. Pour \esp{¿Por qué...?}, la réponse commence obligatoirement par \esp{porque}.

\textbf{Question 1.} Le texte dit \esp{«las mujeres todavía no tienen los mismos derechos que los hombres»}. Réponse rédigée :

\esp{No, las mujeres todavía no tienen los mismos derechos que los hombres: ganan menos por el mismo trabajo y tienen menos acceso a la educación.}

\textbf{Question 2.} Le texte indique la conséquence : \esp{trabajan en lugar de ir a la escuela}. Réponse attendue (reformulation admise) :

\esp{Muchos niños trabajan en lugar de ir a la escuela, es decir, no pueden estudiar porque tienen que trabajar.}

\textbf{Question 3.} La cause est explicite : \esp{porque educar a una niña es educar a toda una familia}. Réponse :

\esp{La escolarización de las niñas es una prioridad porque educar a una niña es educar a toda una familia.}

\textbf{Conclusion méthodologique.} Une réponse complète = reprise de la question + information du texte + verbe conjugué correctement. Aucune réponse d'un seul mot n'est acceptée au Bac.
\end{exoresolu}

\courssub{Lexique thématique : la mujer, el niño y los derechos}

\begin{definition}[Vocabulaire clé du module]
\begin{itemize}
\item \esp{la igualdad} (f.) : l'égalité. \emph{Tous les noms en \esp{-dad} sont féminins.}
\item \esp{el derecho} (m.) : le droit. \esp{los derechos humanos} (les droits de l'homme).
\item \esp{la escolarización} (f.) : la scolarisation / la mise à l'école.
\item \esp{el trabajo infantil} (m.) : le travail des enfants. \emph{\esp{infantil} est invariable en genre.}
\item \esp{el maltrato} (m.) : la maltraitance, les mauvais traitements.
\item \esp{la protección} (f.) : la protection.
\item \esp{la Convención sobre los Derechos del Niño} : la Convention relative aux droits de l'enfant (CIDE, 1989).
\item \esp{la pobreza} (f.) : la pauvreté. \esp{pobre} (adj.) : pauvre.
\item \cle{el gobierno} (m.) : le gouvernement. \cle{las autoridades} (f. pl.) : les autorités.
\end{itemize}
\end{definition}

\begin{methode}[Traduire du français vers l'espagnol (thème) : vocabulaire des droits]
\begin{enumerate}
\item \textbf{Identifier le champ lexical} avant de traduire.
\item \textbf{Attention aux genres} : \esp{el derecho}, \esp{la igualdad}, \esp{el maltrato}, \esp{la protección}, \esp{la escolarización}.
\item \textbf{Attention aux faux amis} : « un enfant » = \esp{un niño} (jamais \esp{un infante} au sens courant) ; « la maltraitance » = \esp{el maltrato} ; « défendre » = \esp{defender} (et non \esp{defenderse} = se défendre).
\item \textbf{Construire la phrase} : sujet + verbe conjugué + complément. Vérifie l'accord de l'adjectif : \esp{los derechos humanos}, \esp{la educación obligatoria}.
\item \textbf{Relire} : accents (\esp{educación}, \esp{protección}, \esp{niño} avec \esp{ñ}), articles, accords.
\end{enumerate}
\end{methode}

\begin{center}
\begin{tabularx}{\linewidth}{|X|X|X|}
\hline
\rowcolor{popblueL} Français & Español & Remarque \\
\hline
l'égalité des droits & la igualdad de derechos & \esp{-dad} = féminin \\
\hline
le travail des enfants & el trabajo infantil & \esp{infantil} invariable \\
\hline
la scolarisation des filles & la escolarización de las niñas & accent sur \esp{-ción} \\
\hline
protéger les enfants & proteger a los niños & \esp{a} personnelle \\
\hline
défendre les droits & defender los derechos & verbe \esp{e$\rightarrow$ie} \\
\hline
le droit à l'éducation & el derecho a la educación & \esp{derecho a} + nom/inf. \\
\hline
la pauvreté des familles & la pobreza de las familias & \esp{pobreza} (nom) \\
\hline
\end{tabularx}
\end{center}

\begin{exoresolu}[Thème : phrases sur les droits]
\textbf{Énoncé} : Traduis en espagnol.
\begin{enumerate}
\item Tous les enfants ont le droit d'aller à l'école.
\item Nous devons protéger les enfants contre la maltraitance.
\item L'égalité entre les hommes et les femmes est nécessaire.
\end{enumerate}
\tcblower
\textbf{Phrase 1.} « Tous les enfants » = \esp{todos los niños} ; « avoir le droit de » = \esp{tener derecho a} + infinitif ; « aller à l'école » = \esp{ir a la escuela} (ou \esp{ir al colegio}).

\esp{Todos los niños tienen derecho a ir a la escuela.}

Justification : \esp{tener} est irrégulier (\esp{e$\rightarrow$ie}) : \esp{tienen} à la 3\up{e} personne du pluriel.

\textbf{Phrase 2.} « Nous devons » = \esp{debemos} ou \esp{tenemos que} ; « protéger quelqu'un » = \esp{proteger a alguien} (\esp{a} personnelle devant un complément de personne) ; « contre » = \esp{contra}.

\esp{Debemos proteger a los niños contra el maltrato.} (ou : \esp{Tenemos que proteger a los niños del maltrato.})

\textbf{Phrase 3.} « L'égalité entre... et... » = \esp{la igualdad entre... y...} ; « nécessaire » s'accorde : \esp{necesaria}.

\esp{La igualdad entre los hombres y las mujeres es necesaria.}

\textbf{Conclusion.} Trois réflexes ont fait la différence : \esp{tener derecho a} + infinitif, l'\esp{a} personnelle de \esp{proteger a}, et l'accord de \esp{necesaria} avec \esp{la igualdad}.
\end{exoresolu}

\courssub{Structures : opinion, necesidad y obligación}

\begin{popfigure}
\begin{tikzpicture}[
    node distance=1.2cm and 1.5cm,
    main/.style={draw=popblue, thick, rounded corners, fill=popblueL, text width=3.2cm, align=center, font=\small},
    arrow/.style={-Stealth, thick, popdark}
]
\node[main, align=center] (opinaff){\textbf{Opinion affirmative}\\\esp{Creo que} + \textcolor{popgreen}{INDICATIF}};
\node[main, right=of opinaff, align=center] (opinneg){\textbf{Opinion niée}\\\esp{No creo que} + \textcolor{poporange}{SUBJONCTIF}};
\node[main, below=of opinaff, align=center] (nec){\textbf{Nécessité impersonnelle}\\\esp{Es necesario que} + \textcolor{poporange}{SUBJONCTIF}};
\node[main, right=of nec, align=center] (obl){\textbf{Obligation personnelle}\\\esp{Tener que} + \textcolor{poppurple}{INFINITIF}};
\node[main, below=of nec, align=center] (gen){\textbf{Obligation générale}\\\esp{Hay que} + \textcolor{poppurple}{INFINITIF}};

\draw[arrow] (opinaff) -- (opinneg);
\draw[arrow] (opinaff) -- (nec);
\draw[arrow] (nec) -- (obl);
\draw[arrow] (obl) -- (gen);
\end{tikzpicture}
\legende{Organigramme de choix : Indicatif vs Subjonctif vs Infinitif}
\end{popfigure}

\begin{methode}[Bien employer creo que, es necesario que, tener que, hay que]
\begin{enumerate}
\item \textbf{Opinion affirmative} : \esp{creo que} / \esp{pienso que} / \esp{me parece que} + \textbf{indicatif}. \esp{Creo que la educación es un derecho.}
\item \textbf{Opinion niée / doute} : \esp{no creo que} / \esp{no pienso que} / \esp{dudo que} + \textbf{subjonctif}. \esp{No creo que la situación sea justa.}
\item \textbf{Nécessité / obligation impersonnelle} : \esp{es necesario que} / \esp{es importante que} / \esp{es obligatorio que} / \esp{es fundamental que} + \textbf{subjonctif}. \esp{Es necesario que el gobierno proteja a los niños.}
\item \textbf{Obligation personnelle (sujet précis)} : \esp{tener que} + \textbf{infinitif}. \esp{Los padres tienen que escolarizar a sus hijos.}
\item \textbf{Obligation générale (sans sujet précis)} : \esp{hay que} + \textbf{infinitif}. \esp{Hay que luchar contra el trabajo infantil.}
\item \textbf{Réflexe de vérification} : \esp{que} + sujet différent du principal $\rightarrow$ subjonctif ; infinitif $\rightarrow$ même sujet ou impersonnel (\esp{hay que}).
\end{enumerate}
\end{methode}

\begin{attention}[Le piège du subjonctif : formes irrégulières fréquentes]
Après \esp{es necesario que}, \esp{no creo que}, \esp{es importante que}, le verbe se met au \textbf{subjonctif présent}. Formes à connaître par cœur :
\begin{center}
\begin{tabular}{|l|l|l|}
\hline
\rowcolor{popblueL} Infinitif & 1ère pers. sg. (yo) & 3ème pers. sg. (él/ella) / 3ème pl. (ellos) \\
\hline
\esp{ser} & \esp{sea} & \esp{sea} / \esp{sean} \\
\hline
\esp{ir} & \esp{vaya} & \esp{vaya} / \esp{vayan} \\
\hline
\esp{tener} & \esp{tenga} & \esp{tenga} / \esp{tengan} \\
\hline
\esp{hacer} & \esp{haga} & \esp{haga} / \esp{hagan} \\
\hline
\esp{proteger} & \esp{proteja} & \esp{proteja} / \esp{protejan} \\
\hline
\esp{escolarizar} & \esp{escolarice} & \esp{escolarice} / \esp{escolaricen} (\esp{z$\rightarrow$c}) \\
\hline
\esp{desaparecer} & \esp{desaparezca} & \esp{desaparezca} / \esp{desaparezcan} (\esp{c$\rightarrow$zc}) \\
\hline
\esp{ayudar} & \esp{ayude} & \esp{ayude} / \esp{ayuden} \\
\hline
\end{tabular}
\end{center}
\emph{Rappel :} \esp{z $\rightarrow$ c} devant \esp{e} (\esp{escolarice}), \esp{c $\rightarrow$ zc} pour les verbes en \esp{-cer/-cir} (\esp{desaparezca}).
\end{attention}

\begin{exoresolu}[Transformation : indicatif ou subjonctif ?]
\textbf{Énoncé} : Complète avec la forme correcte du verbe entre parenthèses, puis traduis la phrase en français.
\begin{enumerate}
\item \esp{Creo que las mujeres (tener) los mismos derechos que los hombres.}
\item \esp{Es necesario que el Estado (proteger) a los niños.}
\item \esp{No creo que el trabajo infantil (desaparecer) pronto.}
\item \esp{Los niños (tener que) ir a la escuela y no al mercado.}
\item \esp{Es fundamental que el gobierno (ayudar) a las familias pobres.}
\end{enumerate}
\tcblower
\textbf{Phrase 1.} \esp{Creo que} affirmatif $\rightarrow$ \textbf{indicatif présent} : \esp{tienen} (\esp{e$\rightarrow$ie}).

\esp{Creo que las mujeres tienen los mismos derechos que los hombres.} « Je pense que les femmes ont les mêmes droits que les hommes. »

\textbf{Phrase 2.} \esp{Es necesario que} $\rightarrow$ \textbf{subjonctif présent} de \esp{proteger} : \esp{proteja}. Ne pas oublier l'\esp{a} personnelle : \esp{a los niños}.

\esp{Es necesario que el Estado proteja a los niños.} « Il est nécessaire que l'État protège les enfants. »

\textbf{Phrase 3.} \esp{No creo que} (opinion niée) $\rightarrow$ \textbf{subjonctif} : \esp{desaparezca} (\esp{c$\rightarrow$zc}).

\esp{No creo que el trabajo infantil desaparezca pronto.} « Je ne pense pas que le travail des enfants disparaisse bientôt. »

\textbf{Phrase 4.} \esp{Tener que} + infinitif exprime l'obligation ; on conjugue \esp{tener} : \esp{tienen que} (\esp{e$\rightarrow$ie}).

\esp{Los niños tienen que ir a la escuela y no al mercado.} « Les enfants doivent aller à l'école et non au marché. »

\textbf{Phrase 5.} \esp{Es fundamental que} $\rightarrow$ \textbf{subjonctif} : \esp{ayude}.

\esp{Es fundamental que el gobierno ayude a las familias pobres.} « Il est fondamental que le gouvernement aide les familles pauvres. »

\textbf{Conclusion.} La règle décisive : opinion affirmative = indicatif ; opinion niée ou nécessité impersonnelle = subjonctif ; obligation avec sujet = \esp{tener que} + infinitif ; obligation générale = \esp{hay que} + infinitif.
\end{exoresolu}

\courssub{Méthode du commentaire de texte et commentaire modèle}

\begin{methode}[La fiche de commentaire en 5 étapes]
\begin{enumerate}
\item \textbf{Présenter le texte} (2-3 phrases) : type de texte, thème, sujet précis, source si connue. Formule : \esp{Este texto es un testimonio / artículo que trata de... / habla de...}
\item \textbf{Dégager l'idée principale} : quelle thèse l'auteur défend-il ? \esp{La idea principal es que... / El autor defiende que... / El texto denuncia que...}
\item \textbf{Relever les arguments} : cite deux ou trois arguments du texte avec des connecteurs logiques : \esp{primero / en primer lugar}, \esp{además / además de eso}, \esp{por otro lado}, \esp{por eso / por tanto}, \esp{en conclusión}.
\item \textbf{Donner ton opinion personnelle} avec les structures de la leçon : \esp{En mi opinión...}, \esp{Creo que...} + indicatif, \esp{Es necesario que...} + subjonctif, \esp{Tenemos que...} + infinitif. Appuie-toi sur un exemple concret (Cameroun, Guinée équatoriale, monde hispanique).
\item \textbf{Conclure} en une ou deux phrases : \esp{En conclusión, este texto nos muestra que... / En definitiva, es urgente que...}
\end{enumerate}
\textbf{Réflexes} : rester en espagnol ; citer le texte entre guillemets (\esp{«...»}) ; éviter de recopier tout le texte ; soigner les connecteurs ; vérifier les accords et le subjonctif.
\end{methode}

\begin{exoresolu}[Commentaire guidé : modèle rédigé]
\textbf{Énoncé} : Lis le texte suivant, puis rédige un commentaire de 8 à 10 lignes en suivant la fiche méthode.

\esp{En mi barrio de Yaoundé, muchas niñas venden agua en el mercado en lugar de ir a la escuela. Sus familias son pobres y necesitan dinero. Sin embargo, la educación es un derecho de todos los niños. Si una niña va a la escuela, toda su familia sale ganando. El gobierno y los padres deben trabajar juntos para proteger a los niños.}
\tcblower
\textbf{Raisonnement.} On applique les 5 étapes : présentation (témoignage, thème : scolarisation des filles à Yaoundé), idée principale (l'éducation est un droit), arguments (pauvreté des familles, bénéfice pour toute la famille), opinion personnelle avec \esp{creo que} et \esp{es necesario que} + subjonctif, conclusion.

\textbf{Commentaire modèle} :

\esp{Este texto es un testimonio que habla de la escolarización de las niñas en un barrio de Yaoundé. La idea principal es que la educación es un derecho de todos los niños, también de las niñas. El autor da dos argumentos: primero, muchas niñas venden agua en el mercado porque sus familias son pobres; además, afirma que «si una niña va a la escuela, toda su familia sale ganando». En mi opinión, este texto tiene razón: creo que ninguna niña debe trabajar en lugar de estudiar. Es necesario que el gobierno ayude a las familias pobres y que los padres escolaricen a sus hijas. En conclusión, este texto nos muestra que proteger a los niños es un deber de todos.}

\textbf{Traduction du commentaire :} « Ce texte est un témoignage qui parle de la scolarisation des filles dans un quartier de Yaoundé. L'idée principale est que l'éducation est un droit de tous les enfants, y compris des filles. L'auteur donne deux arguments : d'abord, beaucoup de filles vendent de l'eau au marché parce que leurs familles sont pauvres ; en outre, il affirme que "si une fille va à l'école, toute sa famille y gagne". À mon avis, ce texte a raison : je pense qu'aucune fille ne doit travailler au lieu d'étudier. Il est nécessaire que le gouvernement aide les familles pauvres et que les parents scolarisent leurs filles. En conclusion, ce texte nous montre que protéger les enfants est un devoir de tous. »

\textbf{Points forts du modèle} : structures de la leçon utilisées (\esp{creo que} + indicatif \esp{debe/tiene} ; \esp{es necesario que} + subjonctif \esp{ayude}, \esp{escolaricen}), une citation du texte intégrée, des connecteurs variés (\esp{primero}, \esp{además}, \esp{en conclusión}), un ancrage camerounais naturel.
\end{exoresolu}

\courssub{Traduction : version (espagnol → français)}

\begin{methode}[Traduire sans calquer (version)]
\begin{enumerate}
\item \textbf{Lire la phrase entière} avant d'écrire le moindre mot : repère le sujet, le verbe, le temps, les connecteurs.
\item \textbf{Traduire les structures, pas les mots isolés} : \esp{en lugar de} + infinitif = « au lieu de » + infinitif ; \esp{los mismos... que} = « les mêmes... que » ; \esp{cada vez más} = « de plus en plus » ; \esp{sin embargo} = « cependant » ; \esp{por eso} = « c'est pourquoi / pour cette raison ».
\item \textbf{Méfie-toi des faux amis} : \esp{actualmente} = « actuellement » (jamais « en fait / réellement ») ; \esp{asistir a} = « assister à / aller à » (pas « assister » au sens de soigner) ; \esp{una carta} = « une lettre » ; \esp{realizar} = « réaliser / accomplir » (pas « réaliser » au sens de comprendre).
\item \textbf{Rends les temps correctement} : présent espagnol = présent français (souvent) ; \esp{deben} = « doivent » ; \esp{proteja} (subj.) = « protège » (subjonctif français) ou « doive protéger ».
\item \textbf{Relis ta phrase française} : elle doit être naturelle, sans trace de syntaxe espagnole (ex. éviter « il est nécessaire que le gouvernement protège » si « doit protéger » est plus naturel, mais garder le subjonctif si le style le permet).
\end{enumerate}
\end{methode}

\begin{exoresolu}[Version : un paragraphe sur l'égalité]
\textbf{Énoncé} : Traduis en français.

\esp{Actualmente, cada vez más mujeres estudian en la universidad y trabajan en lugar de quedarse en casa. Sin embargo, todavía ganan menos que los hombres por el mismo trabajo. Por eso, muchas organizaciones defienden la igualdad de derechos entre hombres y mujeres.}
\tcblower
\textbf{Analyse.} \esp{Actualmente} = « actuellement » ; \esp{cada vez más} = « de plus en plus » ; \esp{en lugar de quedarse} = « au lieu de rester » ; \esp{sin embargo} = « cependant » ; \esp{por eso} = « c'est pourquoi » ; \esp{defienden} = « défendent » (verbe \esp{defender}, \esp{e$\rightarrow$ie}) ; \esp{la igualdad de derechos} = « l'égalité des droits ».

\textbf{Traduction} :

« Actuellement, de plus en plus de femmes étudient à l'université et travaillent au lieu de rester à la maison. Cependant, elles gagnent encore moins que les hommes pour le même travail. C'est pourquoi de nombreuses organisations défendent l'égalité des droits entre les hommes et les femmes. »

\textbf{Pièges évités} : \esp{actualmente} ne signifie pas « en fait » ; \esp{ganar menos} = « gagner moins » (salaire), pas « gagner » au sens de remporter ; \esp{la igualdad de derechos} se rend naturellement par « l'égalité des droits » ; \esp{en lugar de} + infinitif $\rightarrow$ « au lieu de » + infinitif.

\textbf{Conclusion.} Une bonne version est une phrase française correcte et naturelle qui rend exactement le sens de l'espagnol, structure par structure.
\end{exoresolu}

\courssub{Production guidée : opinar y defender los derechos}

\begin{experience}[Débat oral : « La place de la fille à l'école »]
\textbf{Consigne.} En groupes de 4. L'un joue le rôle d'un père réticent, l'autre d'une mère favorable, le troisième d'un enseignant, le quatrième d'un représentant d'ONG (UNICEF, Plan International). Débattez 5 minutes en espagnol.
\textbf{Structures imposées} (chacun doit en utiliser au moins 2) :
\begin{itemize}
\item \esp{Creo que / Pienso que} + indicatif.
\item \esp{No creo que} + subjonctif.
\item \esp{Es necesario que / Es fundamental que} + subjonctif.
\item \esp{Tenemos que / Hay que} + infinitif.
\item \esp{En mi opinión / A mi juicio}.
\end{itemize}
\textbf{Exemple de réplique :} \esp{-- Padre: No creo que la niña necesite ir a la universidad. -- Madre: ¡Pero es necesario que estudie! Tenemos que cambiar la mentalidad. -- ONG: Hay que proteger sus derechos.}
\end{experience}

\begin{exercice}[Expression écrite : Article d'opinion]
\textbf{Consigne.} Rédige un article de 10-12 lignes pour le journal du lycée sur le thème : \esp{« La educación de las niñas: un derecho, no un favor »}. Utilise le texte support, le lexique et les structures de la leçon. Plan suggéré :
\begin{enumerate}
\item Accroche : situation actuelle (Cameroun / monde).
\item Argument 1 : droit légal (Convention).
\item Argument 2 : bénéfices sociaux (santé, économie).
\item Proposition concrète : \esp{Es necesario que...}, \esp{Hay que...}.
\item Conclusion engageante.
\end{enumerate}
\end{exercice}

\begin{corrige}[Exercice : Expression écrite (proposition de correction)]
\esp{La educación de las niñas: un derecho, no un favor}

\esp{En muchos barrios de Yaoundé y de Douala, todavía vemos niñas vendiendo en los mercados en lugar de estar en el aula. Sin embargo, la Convención sobre los Derechos del Niño, firmada por Camerún, garantiza el derecho a la educación para todos.}

\esp{Primero, la escolarización de las niñas rompe el ciclo de la pobreza. Una mujer educada tiene mejores oportunidades de trabajo y cuida mejor la salud de su familia. Además, la igualdad de género fortalece la sociedad entera.}

\esp{Por eso, creo que el gobierno debe actuar ya. Es necesario que las autoridades construyan más escuelas secundarias en zonas rurales y que garanticen la gratuidad real. Hay que sensibilizar a los padres: los niños no son mano de obra, son el futuro.}

\esp{En conclusión, defender los derechos de la niña no es caridad, es justicia. ¡La escuela es su lugar!}
\end{corrige}

\courssub{Erreurs fréquentes et points de vigilance (Spécial Bac)}

\begin{attention}[Les 5 erreurs qui coûtent des points au Bac]
\begin{enumerate}
\item \textbf{Oublier les accents et la \esp{ñ}} : \esp{educación}, \esp{protección}, \esp{igualdad}, \esp{niño} (jamais \esp{nino}), \esp{país}, \esp{gobierno}, \esp{también}. Un accent oublié change le sens ou le temps : \esp{está} (il est) vs \esp{esta} (cette) ; \esp{habla} (il parle) vs \esp{hablá} (impératif \emph{vos} argentin, hors programme standard).
\item \textbf{Calquer le français sur les comparatifs} : « les mêmes droits que » = \esp{los mismos derechos que} (jamais \esp{como}). « Au lieu de » = \esp{en lugar de} + infinitif (jamais \esp{en lugar de que} + indicatif).
\item \textbf{Mettre l'indicatif après \esp{es necesario que} / \esp{no creo que}} : on écrit \esp{Es necesario que los niños vayan a la escuela} (subj.) et non \esp{van} (ind.). Même piège après \esp{es importante que}, \esp{es urgente que}.
\item \textbf{Oublier l'\esp{a} personnelle} : devant un complément de personne (ou animal personifié), le verbe transitif direct exige \esp{a} : \esp{proteger a los niños}, \esp{educar a una niña}, \esp{ayudar a las familias}, \esp{ver a mis amigos}.
\item \textbf{Faute d'accord de l'adjectif} : l'adjectif s'accorde \textbf{toujours} en genre et en nombre avec le nom qu'il qualifie. \esp{La educación es obligatori\uline{a}} (f. sg.), \esp{Los derechos human\uline{os}} (m. pl.), \esp{Una ley necesari\uline{a}} (f. sg.), \esp{Niños pobres} (m. pl.).
\end{enumerate}
\end{attention}

\begin{savaistu}[Culture hispanique : La CIDE et la Guinée équatoriale]
La \textbf{Convención sobre los Derechos del Niño (CIDE)}, adoptée par l'ONU en 1989, est le traité des droits de l'homme le plus ratifié de l'histoire (196 États). Seuls les États-Unis ne l'ont pas ratifiée.
\textbf{Guinée équatoriale} (seul pays africain hispanophone) : a ratifié la CIDE en 1992. Le pays fait des efforts pour la \esp{escolarización} (taux de scolarisation primaire élevé), mais des défis persistent : \esp{trabajo infantil} dans l'agriculture et le commerce informel, \esp{matrimonio infantil} (mariage précoce des filles) dans certaines zones rurales.
\textbf{Cameroun} : bien que francophone/anglophone, le Cameroun est membre observateur de l'OIF et proche de l'espace hispanophone (frontière avec la Guinée équatoriale). La \esp{Ley de Orientación de la Educación} (1998) et la Constitution garantissent l'éducation obligatoire. Le défi majeur reste la scolarisation des filles dans le Grand-Nord (Extrême-Nord, Nord, Adamaoua) où la pauvreté et les coutumes frequent la \esp{escolarización de las niñas}.
\textbf{Citation} (Gabriela Mistral, prix Nobel chilienne) : \esp{« Muchos de los que hoy son hombres, fueron niños que no tuvieron infancia. »} (« Beaucoup de ceux qui sont aujourd'hui des hommes étaient des enfants qui n'ont pas eu d'enfance. »)
\end{savaistu}

\newpage

\coursec{Exercices}

\courssub{Je vérifie mes connaissances}

\begin{exercice}[QCM : le texte et les droits]
Entoure la bonne réponse.
\begin{enumerate}
\item La Convention relative aux droits de l'enfant a été adoptée en :
\begin{enumerate}
\item 1945 \item 1959 \item 1989 \item 2001
\end{enumerate}
\item Selon la Convention, un enfant est toute personne âgée de moins de :
\begin{enumerate}
\item 15 ans \item 18 ans \item 21 ans \item 16 ans
\end{enumerate}
\item \esp{La escolarización} signifie :
\begin{enumerate}
\item la scolarisation \item la scolarité payante \item l'école privée \item le redoublement
\end{enumerate}
\item Après \esp{es necesario que}, on emploie :
\begin{enumerate}
\item l'infinitif \item l'indicatif présent \item le subjonctif présent \item le futur
\end{enumerate}
\item \esp{El maltrato infantil} désigne :
\begin{enumerate}
\item le jeu des enfants \item la maltraitance des enfants \item le travail des enfants \item l'école des enfants
\end{enumerate}
\end{enumerate}
\end{exercice}

\begin{exercice}[Vrai ou faux justifié]
Dis si les affirmations suivantes sont vraies ou fausses, puis justifie ta réponse par une phrase en français.
\begin{enumerate}
\item La Convention relative aux droits de l'enfant a été adoptée par l'ONU en 1989.
\item La Convention ne concerne que les enfants de moins de 12 ans.
\item Tout enfant a droit à l'éducation selon la Convention.
\item Le travail des enfants est autorisé sans limite s'il aide la famille.
\item L'égalité entre l'homme et la femme est déjà une réalité partout dans le monde.
\end{enumerate}
\end{exercice}

\begin{exercice}[Vocabulaire : le champ des droits]
Complète chaque phrase avec le mot espagnol qui convient : \esp{igualdad} — \esp{derechos} — \esp{escolarización} — \esp{trabajo infantil} — \esp{maltrato} — \esp{protección}.
\begin{enumerate}
\item Los niños tienen \ldots\ldots\ldots\ldots a la educación y a la salud.
\item La \ldots\ldots\ldots\ldots de las niñas es un problema en algunas zonas rurales.
\item El \ldots\ldots\ldots\ldots impide que muchos niños vayan a la escuela.
\item La \ldots\ldots\ldots\ldots entre el hombre y la mujer es un objetivo importante.
\item El Estado debe garantizar la \ldots\ldots\ldots\ldots de los menores.
\item Hay que luchar contra el \ldots\ldots\ldots\ldots de los niños en la familia.
\end{enumerate}
\end{exercice}

\begin{exercice}[Conjugaison : indicatif ou subjonctif]
Conjugue le verbe entre parenthèses au présent de l'indicatif ou du subjonctif selon la structure.
\begin{enumerate}
\item Creo que la igualdad (ser) \ldots\ldots\ldots\ldots posible.
\item No creo que las niñas (tener) \ldots\ldots\ldots\ldots menos derechos que los niños.
\item Es necesario que todos los niños (ir) \ldots\ldots\ldots\ldots a la escuela.
\item Pienso que el gobierno (proteger) \ldots\ldots\ldots\ldots a los menores.
\item Es importante que los padres (escuchar) \ldots\ldots\ldots\ldots a sus hijos.
\item No pienso que el trabajo infantil (poder) \ldots\ldots\ldots\ldots justificarse.
\end{enumerate}
\end{exercice}

\courssub{Je m'entraîne}

\begin{exercice}[Thème : cinq phrases à traduire]
Traduis en espagnol les phrases suivantes.
\begin{enumerate}
\item Les filles doivent aller à l'école.
\item Il faut protéger les enfants contre la maltraitance.
\item Je crois que l'égalité entre l'homme et la femme est nécessaire.
\item Il est important que les enfants aient des droits.
\item Le travail des enfants est un problème grave au Cameroun.
\end{enumerate}
\end{exercice}

\begin{exercice}[Transformation : opinion et obligation]
Transforme chaque phrase en utilisant la structure indiquée entre parenthèses, sans changer le sens.
\begin{enumerate}
\item Los niños tienen derecho a la educación. (\esp{Creo que}\ldots)
\item Las niñas no siempre van a la escuela en el campo. (\esp{No creo que}\ldots)
\item El Estado protege a los menores. (\esp{Es necesario que}\ldots)
\item Los padres escolarizan a sus hijas. (\esp{Es importante que}\ldots)
\item La sociedad lucha contra el maltrato infantil. (\esp{Tenemos que}\ldots\ + infinitif)
\end{enumerate}
\end{exercice}

\begin{exercice}[Texte à trous : la condition de la femme]
Complète le texte suivant avec les mots de la liste : \esp{igualdad} — \esp{derechos} — \esp{trabajo} — \esp{educación} — \esp{hijas} — \esp{necesario}.
\begin{textebox}[La mujer camerunesa hoy]
En Camerún, la mujer ha progresado mucho, pero la \ldots\ldots\ldots\ldots total entre el hombre y la mujer no es todavía una realidad. Muchas familias envían a sus \ldots\ldots\ldots\ldots a la escuela, pero en algunas regiones rurales la \ldots\ldots\ldots\ldots de las niñas sigue siendo difícil. Las mujeres trabajan en los campos, en los mercados y en las oficinas, pero su \ldots\ldots\ldots\ldots no siempre está bien pagado. Es \ldots\ldots\ldots\ldots que la sociedad respete los \ldots\ldots\ldots\ldots de la mujer para construir un país más justo.
\end{textebox}
\end{exercice}

\begin{exercice}[Appariements : mots et définitions]
Associe chaque mot espagnol (1 à 10) à sa définition française (a à j), puis emploie cinq de ces mots dans des phrases complètes en espagnol.
\begin{center}
\begin{tabularx}{\linewidth}{|l|X|}
\hline
\rowcolor{popblueL} \textbf{Español} & \textbf{Français} \\
\hline
1. \esp{la igualdad} & a. le fait de frapper ou de blesser un enfant \\
\hline
2. \esp{los derechos} & b. le fait d'inscrire les enfants à l'école \\
\hline
3. \esp{la escolarización} & c. ce que la loi reconnaît à chaque personne \\
\hline
4. \esp{el trabajo infantil} & d. le fait d'avoir les mêmes droits et chances \\
\hline
5. \esp{el maltrato} & e. le fait de défendre les plus faibles \\
\hline
6. \esp{la protección} & f. le fait d'employer des enfants pour travailler \\
\hline
7. \esp{el matrimonio precoz} & g. un texte international signé par les États \\
\hline
8. \esp{la convención} & h. le mariage d'une fille trop jeune \\
\hline
9. \esp{la educación} & i. le fait de maltraiter, en anglais \emph{abuse} \\
\hline
10. \esp{el abuso} & j. l'instruction reçue à l'école et en famille \\
\hline
\end{tabularx}
\end{center}
\end{exercice}

\courssub{Je me prépare au Bac}

\begin{exercice}[Compréhension écrite type Bac]
Lis le texte suivant, puis réponds aux questions en espagnol, avec des phrases complètes.
\begin{textebox}[El matrimonio precoz en Camerún]
En algunas regiones del norte de Camerún, muchas niñas se casan antes de los quince años. Sus familias, a menudo pobres, piensan que el matrimonio es una solución económica. Pero esta tradición tiene consecuencias graves : la niña deja la escuela, no termina sus estudios y queda dependiente de su marido. Además, los embarazos precoces son peligrosos para su salud.

Sin embargo, la situación empieza a cambiar. Asociaciones de mujeres y jóvenes cameruneses sensibilizan a los padres en los pueblos. Explican que una niña educada puede ayudar mejor a su familia más tarde, con un trabajo digno. Algunas madres que se casaron muy jóvenes testifican : « No quiero que mi hija viva lo que yo he vivido. Quiero que estudie y que elija su vida. »

La ley camerunesa y la Convención sobre los Derechos del Niño protegen a los menores, pero hace falta aplicarlas en todas partes. La escuela es la mejor arma : cuando una niña estudia, toda la comunidad progresa. Como dice un proverbio africano : « Educar a una niña es educar a una nación. »
\end{textebox}
\begin{enumerate}
\item ¿Por qué algunas familias casan a sus hijas muy jóvenes ? (2 pts)
\item Cita del texto dos consecuencias del matrimonio precoz para la niña. (2 pts)
\item ¿Qué hacen las asociaciones para cambiar la situación ? (2 pts)
\item ¿Qué significa el testimonio de las madres ? Explica con tus palabras. (2 pts)
\item \textbf{Langue :} dans la phrase \esp{« No quiero que mi hija viva lo que yo he vivido »}, justifie l'emploi du subjonctif \esp{viva}, puis conjugue \esp{querer que} + \esp{estudiar} à toutes les personnes du présent. (4 pts)
\item \textbf{Expression personnelle :} ¿Estás de acuerdo con el proverbio final ? Responde en tres o cuatro frases. (3 pts)
\end{enumerate}
\end{exercice}

\begin{exercice}[Thème et version type Bac]
\textbf{A. Thème.} Traduis en espagnol les phrases suivantes. (5 pts)
\begin{enumerate}
\item Tous les enfants ont le droit d'aller à l'école.
\item Il est nécessaire que les parents protègent leurs enfants.
\item Je ne crois pas que le mariage précoce soit une bonne solution.
\item Les femmes doivent travailler et gagner le même salaire que les hommes.
\item Il faut lutter contre le travail des enfants en Afrique.
\end{enumerate}
\textbf{B. Version.} Traduis en français le passage suivant. (5 pts)
\begin{textebox}[Los derechos del niño]
Todo niño tiene derecho a una familia, a un nombre y a una educación. Ningún niño debe trabajar antes de la edad legal ni sufrir maltrato. Los Estados que firmaron la Convención de 1989 deben garantizar estos derechos y castigar a los que no los respetan. Proteger a los niños es proteger el futuro de la humanidad.
\end{textebox}
\end{exercice}

\begin{exercice}[Expression écrite type Bac]
Rédige en espagnol un article de \textbf{100 à 120 mots} intitulé :
\begin{center}
\esp{La igualdad entre el hombre y la mujer : ¿una realidad en Camerún?}
\end{center}
Consignes :
\begin{itemize}
\item Présente la situation de la femme au Cameroun (éducation, travail, famille).
\item Donne ton opinion personnelle avec au moins deux structures étudiées : \esp{creo que}, \esp{no creo que}, \esp{es necesario que}, \esp{tener que} + infinitif.
\item Propose une conclusion avec une ou deux solutions concrètes.
\end{itemize}
Tu seras évalué sur : le respect du sujet (4 pts), la correction grammaticale (4 pts), la richesse du vocabulaire (4 pts), l'organisation du texte (3 pts).
\end{exercice}

\newpage

\coursec{Corrigés des exercices}

\begin{corrige}[Exercice 1 — QCM : le texte et les droits]
\begin{enumerate}
\item \textbf{c. 1989.} La Convention relative aux droits de l'enfant a été adoptée par l'Assemblée générale des Nations unies le 20 novembre 1989.
\item \textbf{b. 18 ans.} Selon l'article premier de la Convention, un enfant est tout être humain âgé de moins de dix-huit ans.
\item \textbf{a. la scolarisation.} \esp{La escolarización} = le fait d'inscrire les enfants à l'école. Attention au faux ami : « la scolarité » (les études, les frais) se dit plutôt \esp{la escolaridad}.
\item \textbf{c. le subjonctif présent.} Après une expression impersonnelle de nécessité comme \esp{es necesario que}, le verbe subordonné est toujours au subjonctif : \esp{Es necesario que todos los niños vayan a la escuela.}
\item \textbf{b. la maltraitance des enfants.} \esp{El maltrato infantil} = les mauvais traitements infligés aux enfants (coups, violences, négligence).
\end{enumerate}
\textbf{Points de vigilance :} ne pas confondre \esp{la escolarización} (scolarisation) et \esp{el colegio} (l'école, l'établissement) ; retenir la date 1989, très souvent demandée.
\end{corrige}

\begin{corrige}[Exercice 2 — Vrai ou faux justifié]
\begin{enumerate}
\item \textbf{Vrai.} La Convention relative aux droits de l'enfant a bien été adoptée par l'ONU le 20 novembre 1989 ; c'est le texte international le plus ratifié au monde.
\item \textbf{Faux.} La Convention concerne toute personne âgée de moins de dix-huit ans, et non seulement les moins de douze ans.
\item \textbf{Vrai.} L'article 28 de la Convention reconnaît à tout enfant le droit à l'éducation, et l'école primaire doit être obligatoire et gratuite.
\item \textbf{Faux.} La Convention interdit le travail dangereux et le travail qui empêche la scolarisation ; aider sa famille ne justifie pas l'exploitation de l'enfant.
\item \textbf{Faux.} Malgré les progrès, des inégalités persistent partout : salaires plus faibles, accès difficile à l'éducation des filles, mariages précoces, sous-représentation politique.
\end{enumerate}
\textbf{Points de vigilance :} dans une justification, cite toujours un élément précis (un âge, un droit, un exemple) et non une simple répétition de l'affirmation.
\end{corrige}

\begin{corrige}[Exercice 3 — Vocabulaire : le champ des droits]
\begin{enumerate}
\item Los niños tienen \textbf{derechos} a la educación y a la salud. « Les enfants ont droit à l'éducation et à la santé. »
\item La \textbf{escolarización} de las niñas es un problema en algunas zonas rurales. « La scolarisation des filles est un problème dans certaines zones rurales. »
\item El \textbf{trabajo infantil} impide que muchos niños vayan a la escuela. « Le travail des enfants empêche beaucoup d'enfants d'aller à l'école. »
\item La \textbf{igualdad} entre el hombre y la mujer es un objetivo importante. « L'égalité entre l'homme et la femme est un objectif important. »
\item El Estado debe garantizar la \textbf{protección} de los menores. « L'État doit garantir la protection des mineurs. »
\item Hay que luchar contra el \textbf{maltrato} de los niños en la familia. « Il faut lutter contre la maltraitance des enfants dans la famille. »
\end{enumerate}
\textbf{Remarques grammaticales :} \esp{tener derecho a} + nom ou infinitif ; \esp{impedir que} + subjonctif (\esp{vayan}) ; \esp{hay que} + infinitif = « il faut ».
\end{corrige}

\begin{corrige}[Exercice 4 — Conjugaison : indicatif ou subjonctif]
\begin{enumerate}
\item Creo que la igualdad \textbf{es} posible. → \esp{creo que} + indicatif (opinion affirmée).
\item No creo que las niñas \textbf{tengan} menos derechos que los niños. → \esp{no creo que} + subjonctif (opinion niée = doute).
\item Es necesario que todos los niños \textbf{vayan} a la escuela. → subjonctif de \esp{ir}, irrégulier : vaya, vayas, vaya, vayamos, vayáis, vayan.
\item Pienso que el gobierno \textbf{protege} a los menores. → \esp{pienso que} affirmatif + indicatif.
\item Es importante que los padres \textbf{escuchen} a sus hijos. → expression impersonnelle + subjonctif.
\item No pienso que el trabajo infantil \textbf{pueda} justificarse. → opinion niée + subjonctif de \esp{poder} (pueda).
\end{enumerate}
\textbf{Règle à retenir :} \esp{creo que / pienso que} affirmatifs → indicatif ; les mêmes verbes à la forme négative et les tournures \esp{es necesario que}, \esp{es importante que} → subjonctif.
\end{corrige}

\begin{corrige}[Exercice 5 — Thème : cinq phrases à traduire]
\begin{enumerate}
\item \esp{Las niñas tienen que ir a la escuela.} (ou \esp{Las niñas deben ir a la escuela.})
\item \esp{Hay que proteger a los niños contra el maltrato.} (ou \esp{Es necesario proteger a los niños del maltrato.})
\item \esp{Creo que la igualdad entre el hombre y la mujer es necesaria.}
\item \esp{Es importante que los niños tengan derechos.} → subjonctif \esp{tengan} obligatoire.
\item \esp{El trabajo infantil es un problema grave en Camerún.}
\end{enumerate}
\textbf{Points de vigilance :}
\begin{itemize}
\item « Les filles doivent » : \esp{tener que} + infinitif ou \esp{deber} + infinitif, jamais \esp{deber de}.
\item « Protéger les enfants » : a personnelle obligatoire → \esp{proteger \textbf{a} los niños}.
\item « Au Cameroun » : \esp{en Camerún} (accent sur la dernière syllabe, sans article).
\item « Il faut » impersonnel : \esp{hay que} + infinitif ; si l'on précise le sujet, \esp{es necesario que} + subjonctif.
\end{itemize}
\end{corrige}

\begin{corrige}[Exercice 6 — Transformation : opinion et obligation]
\begin{enumerate}
\item \esp{Creo que los niños tienen derecho a la educación.} → indicatif après \esp{creo que}.
\item \esp{No creo que las niñas vayan siempre a la escuela en el campo.} → subjonctif \esp{vayan} après \esp{no creo que} ; attention à la place de la négation : \esp{no siempre} devient \esp{no creo que\ldots{} siempre}.
\item \esp{Es necesario que el Estado proteja a los menores.} → subjonctif \esp{proteja} (\esp{proteger} : protejo, protejas, proteja\ldots).
\item \esp{Es importante que los padres escolaricen a sus hijas.} → subjonctif \esp{escolaricen} ; changement orthographique c → c devant e (\esp{escolarizar} : escolarizo, escolari\textbf{c}en).
\item \esp{Tenemos que luchar contra el maltrato infantil.} → \esp{tener que} + infinitif, sans changement de mode.
\end{enumerate}
\textbf{Point de vigilance :} quand la tournure impersonnelle introduit un sujet nouveau (\esp{el Estado}, \esp{los padres}), le subjonctif est obligatoire ; avec \esp{tenemos que}, le sujet reste le même, donc infinitif.
\end{corrige}

\begin{corrige}[Exercice 7 — Texte à trous : la condition de la femme]
Texte complété :
\begin{quote}
\esp{En Camerún, la mujer ha progresado mucho, pero la \textbf{igualdad} total entre el hombre y la mujer no es todavía una realidad. Muchas familias envían a sus \textbf{hijas} a la escuela, pero en algunas regiones rurales la \textbf{educación} de las niñas sigue siendo difícil. Las mujeres trabajan en los campos, en los mercados y en las oficinas, pero su \textbf{trabajo} no siempre está bien pagado. Es \textbf{necesario} que la sociedad respete los \textbf{derechos} de la mujer para construir un país más justo.}
\end{quote}
Traduction : « Au Cameroun, la femme a beaucoup progressé, mais l'égalité totale entre l'homme et la femme n'est pas encore une réalité. Beaucoup de familles envoient leurs filles à l'école, mais dans certaines régions rurales l'éducation des filles reste difficile. Les femmes travaillent dans les champs, dans les marchés et dans les bureaux, mais leur travail n'est pas toujours bien payé. Il est nécessaire que la société respecte les droits de la femme pour construire un pays plus juste. »
\textbf{Indices utiles :} \esp{sus hijas} (envoyer « ses filles » à l'école) ; \esp{su trabajo} (possessif singulier devant nom singulier) ; \esp{es necesario que} + subjonctif \esp{respete}.
\end{corrige}

\begin{corrige}[Exercice 8 — Appariements : mots et définitions]
\textbf{Appariements :} 1-d ; 2-c ; 3-b ; 4-f ; 5-a ; 6-e ; 7-h ; 8-g ; 9-j ; 10-i.

\textbf{Exemples de phrases (cinq mots au choix) :}
\begin{enumerate}
\item \esp{La igualdad de salario entre hombres y mujeres es un derecho.} « L'égalité de salaire entre les hommes et les femmes est un droit. »
\item \esp{Todos los niños tienen derechos, y nadie puede quitárselos.} « Tous les enfants ont des droits, et personne ne peut les leur enlever. »
\item \esp{El trabajo infantil impide la escolarización de muchos niños.} « Le travail des enfants empêche la scolarisation de beaucoup d'enfants. »
\item \esp{El matrimonio precoz es un problema grave en algunas regiones.} « Le mariage précoce est un problème grave dans certaines régions. »
\item \esp{La Convención de 1989 garantiza la protección de los menores.} « La Convention de 1989 garantit la protection des mineurs. »
\end{enumerate}
\textbf{Point de vigilance :} \esp{el abuso} (i) et \esp{el maltrato} (a) sont proches : \esp{maltrato} insiste sur les coups et la violence physique, \esp{abuso} sur l'abus de pouvoir ou l'abus sexuel.
\end{corrige}

\begin{corrige}[Exercice 9 — Compréhension écrite type Bac]
\textbf{Corrigé détaillé avec barème indicatif :}
\begin{enumerate}
\item \esp{Algunas familias casan a sus hijas muy jóvenes porque son pobres y piensan que el matrimonio es una solución económica.} « Certaines familles marient leurs filles très jeunes parce qu'elles sont pauvres et pensent que le mariage est une solution économique. » (2 pts : 1 pt pour la pauvreté, 1 pt pour la solution économique.)
\item \esp{La niña deja la escuela y no termina sus estudios ; además, los embarazos precoces son peligrosos para su salud.} « La fille quitte l'école et ne termine pas ses études ; de plus, les grossesses précoces sont dangereuses pour sa santé. » (2 pts : 1 pt par conséquence citée.)
\item \esp{Las asociaciones sensibilizan a los padres en los pueblos y explican que una niña educada puede ayudar mejor a su familia más tarde.} « Les associations sensibilisent les parents dans les villages et expliquent qu'une fille éduquée peut mieux aider sa famille plus tard. » (2 pts.)
\item \esp{El testimonio significa que las madres que sufrieron el matrimonio precoz no quieren que sus hijas vivan la misma experiencia : quieren que estudien y que elijan su vida.} « Le témoignage signifie que les mères qui ont souffert du mariage précoce ne veulent pas que leurs filles vivent la même expérience : elles veulent qu'elles étudient et qu'elles choisissent leur vie. » (2 pts : reformulation personnelle exigée, pas de recopie.)
\item \textbf{Langue.} \esp{viva} est au subjonctif parce que \esp{querer que} exprime une volonté portant sur un autre sujet (\esp{mi hija}) ; après un verbe de volonté + \esp{que}, l'espagnol exige le subjonctif. Conjugaison de \esp{querer que} + \esp{estudiar} au présent :
\begin{center}
\begin{tabular}{|l|l|}
\hline
\rowcolor{popblueL} \textbf{Personne} & \textbf{Forme} \\
\hline
yo & quiero que (tú) estudies \\
\hline
tú & quieres que (él/ella) estudie \\
\hline
él/ella & quiere que (nosotros) estudiemos \\
\hline
nosotros & queremos que (vosotros) estudiéis \\
\hline
vosotros & queréis que (ellos) estudien \\
\hline
ellos & quieren que (yo) estudie \\
\hline
\end{tabular}
\end{center}
(4 pts : 2 pts pour la justification, 2 pts pour la conjugaison.)
\item \textbf{Expression personnelle (modèle) :} \esp{Sí, estoy de acuerdo con el proverbio. Creo que educar a una niña es educar a toda la sociedad, porque una mujer educada transmite sus conocimientos a sus hijos. Además, es necesario que las niñas estudien para tener un trabajo digno y ser independientes. Por eso pienso que la educación de las niñas es la clave del progreso de una nación.} « Oui, je suis d'accord avec le proverbe. Je crois qu'éduquer une fille, c'est éduquer toute la société, car une femme éduquée transmet ses connaissances à ses enfants. De plus, il est nécessaire que les filles étudient pour avoir un travail digne et être indépendantes. C'est pourquoi je pense que l'éducation des filles est la clé du progrès d'une nation. » (3 pts : opinion claire 1 pt, arguments 1 pt, correction 1 pt.)
\end{enumerate}
\textbf{Points de vigilance :} répondre toujours par des phrases complètes ; ne pas recopier le texte mot à mot aux questions 3 et 4 ; vérifier les accents (\esp{educación}, \esp{económica}, \esp{además}).
\end{corrige}

\begin{corrige}[Exercice 10 — Thème et version type Bac]
\textbf{A. Thème (5 pts, 1 pt par phrase).}
\begin{enumerate}
\item \esp{Todos los niños tienen derecho a ir a la escuela.}
\item \esp{Es necesario que los padres protejan a sus hijos.} → subjonctif \esp{protejan} ; a personnelle : \esp{a sus hijos}.
\item \esp{No creo que el matrimonio precoz sea una buena solución.} → subjonctif \esp{sea} après \esp{no creo que}.
\item \esp{Las mujeres tienen que trabajar y ganar el mismo salario que los hombres.}
\item \esp{Hay que luchar contra el trabajo infantil en África.}
\end{enumerate}
\textbf{Points de vigilance :} \esp{tener derecho a} + infinitif ; « le même\ldots{} que » = \esp{el mismo\ldots{} que} ; \esp{África} prend un accent ; « il faut » impersonnel = \esp{hay que} + infinitif.

\textbf{B. Version (5 pts).}
« Tout enfant a droit à une famille, à un nom et à une éducation. Aucun enfant ne doit travailler avant l'âge légal ni subir de maltraitance. Les États qui ont signé la Convention de 1989 doivent garantir ces droits et punir ceux qui ne les respectent pas. Protéger les enfants, c'est protéger l'avenir de l'humanité. »

\textbf{Points de vigilance pour la version :}
\begin{itemize}
\item \esp{ningún niño\ldots{} ni\ldots} : double négation espagnole → « aucun enfant\ldots{} ni ».
\item \esp{castigar} = punir, châtier (et non « châtier » au sens religieux) ; ne pas traduire par « castrer » (faux ami).
\item \esp{los que no los respetan} : \esp{los} reprend \esp{derechos} → « ceux qui ne les respectent pas ».
\item \esp{proteger\ldots{} es proteger\ldots} : tournure d'équivalence → « protéger\ldots, c'est protéger\ldots ».
\end{itemize}
\textbf{Barème indicatif :} 1 pt par segment de sens correctement traduit, moins 0,25 pt par faute de grammaire ou de contresens.
\end{corrige}

\begin{corrige}[Exercice 11 — Expression écrite type Bac]
\textbf{Proposition de rédaction modèle (environ 115 mots) :}
\begin{quote}
\esp{La igualdad entre el hombre y la mujer : ¿una realidad en Camerún?}

\esp{En Camerún, la mujer ha progresado mucho : trabaja en los mercados, en las oficinas y en la política. Sin embargo, la igualdad total no es todavía una realidad. En algunas regiones, las niñas dejan la escuela antes que los niños y los salarios de las mujeres son más bajos.}

\esp{Creo que la igualdad es necesaria para el progreso del país, pero no creo que sea fácil conseguirla rápidamente. Es necesario que el gobierno aplique las leyes y que los padres escolaricen a sus hijas.}

\esp{Para terminar, pienso que tenemos que sensibilizar a las familias y luchar contra el matrimonio precoz. Una sociedad justa es una sociedad donde hombres y mujeres tienen los mismos derechos.}
\end{quote}
Traduction : « Au Cameroun, la femme a beaucoup progressé : elle travaille dans les marchés, dans les bureaux et en politique. Cependant, l'égalité totale n'est pas encore une réalité. Dans certaines régions, les filles quittent l'école avant les garçons et les salaires des femmes sont plus bas. Je crois que l'égalité est nécessaire pour le progrès du pays, mais je ne crois pas qu'elle soit facile à obtenir rapidement. Il est nécessaire que le gouvernement applique les lois et que les parents scolarisent leurs filles. Pour terminer, je pense que nous devons sensibiliser les familles et lutter contre le mariage précoce. Une société juste est une société où hommes et femmes ont les mêmes droits. »

\textbf{Barème indicatif (15 pts) :}
\begin{itemize}
\item Respect du sujet et des consignes : 4 pts (situation présentée, opinion personnelle, conclusion avec solutions).
\item Correction grammaticale : 4 pts (modes corrects après \esp{creo que} / \esp{es necesario que}, accords, accents).
\item Richesse du vocabulaire : 4 pts (lexique des droits : \esp{igualdad, derechos, escolarizar, matrimonio precoz, salario}\ldots).
\item Organisation : 3 pts (introduction, développement, conclusion ; connecteurs \esp{sin embargo, además, para terminar}).
\end{itemize}
\textbf{Points de vigilance :} respecter la fourchette de 100 à 120 mots ; employer au moins deux structures de la leçon avec le bon mode (indicatif après \esp{creo que}, subjonctif après \esp{no creo que} et \esp{es necesario que}) ; soigner les accents (\esp{política, rápidamente, necesario}) ; éviter le calque français « las mujeres son más basses payées » → \esp{los salarios de las mujeres son más bajos}.
\end{corrige}

\newpage

\coursec{Fiche bilan}

\begin{popfigure}
\begin{tikzpicture}[mindmap, grow cyclic, every node/.style={concept, text=white, font=\small},
  root concept/.append style={concept color=popblue, font=\bfseries},
  level 1/.append style={level distance=3.2cm, sibling angle=60, font=\footnotesize},
  level 1 concept/.append style={concept color=poporange}]
\node[root concept]{La mujer y los derechos del niño}
  child[concept color=popgreen]{ node {Lexique : igualdad, derechos, escolarización} }
  child[concept color=poporange]{ node {Protección : maltrato, trabajo infantil} }
  child[concept color=poppurple]{ node {Opinion : creo que, pienso que, me parece que} }
  child[concept color=poppink]{ node {Obligation : tener que + inf., hay que} }
  child[concept color=popgold]{ node {Nécessité : es necesario que + subj.} }
  child[concept color=popblue]{ node {Méthode : commenter un texte} };
\end{tikzpicture}
\legende{Carte mentale de la leçon : la femme, l'enfant et les droits}
\end{popfigure}

\begin{bilan}
\textbf{1. Lexique clé à connaître par cœur}
\begin{itemize}
\item \esp{la igualdad} : l'égalité ; \esp{la desigualdad} : l'inégalité ; \esp{los derechos} : les droits
\item \esp{la escolarización} : la scolarisation ; \esp{la educación} : l'éducation ; \esp{el analfabetismo} : l'analphabétisme
\item \esp{el trabajo infantil} : le travail des enfants ; \esp{el maltrato} : la maltraitance ; \esp{el abuso} : l'abus
\item \esp{la protección} : la protection ; \esp{proteger} : protéger ; \esp{defender los derechos} : défendre les droits
\item \esp{la Convención sobre los Derechos del Niño} : la Convention relative aux droits de l'enfant (1989)
\item \esp{el matrimonio precoz / forzado} : le mariage précoce / forcé ; \esp{la violencia doméstica} : la violence domestique
\item \esp{la mujer trabajadora} : la femme active ; \esp{el salario} : le salaire ; \esp{el acceso al trabajo} : l'accès au travail
\end{itemize}

\textbf{2. Grammaire : opinion, nécessité, obligation}
\begin{center}
\begin{tabularx}{\linewidth}{|l|X|X|}
\hline
\rowcolor{popblueL} Structure & Emploi & Exemple \\
\hline
\esp{creo que / pienso que + indicatif} & exprimer une opinion affirmative & \esp{Creo que la igualdad es importante.} \\
\hline
\esp{no creo que + subjonctif} & opinion niée (doute) & \esp{No creo que sea justo.} \\
\hline
\esp{es necesario que + subjonctif} & nécessité générale & \esp{Es necesario que los niños vayan a la escuela.} \\
\hline
\esp{tener que + infinitif} & obligation personnelle & \esp{Los padres tienen que proteger a sus hijos.} \\
\hline
\esp{hay que + infinitif} & obligation impersonnelle & \esp{Hay que respetar los derechos del niño.} \\
\hline
\esp{estar de acuerdo con / en contra de} & accord / désaccord & \esp{Estoy a favor de la escolarización de las niñas.} \\
\hline
\end{tabularx}
\end{center}

\textbf{3. Méthode express : commenter un texte}
\begin{enumerate}
\item Lire deux fois ; repérer le thème, l'auteur, le type de texte.
\item Dégager l'idée principale (\esp{la idea principal}) et les idées secondaires.
\item Relever le lexique des droits et les structures d'opinion.
\item Rédiger : introduction (présentation), développement (idées + exemples), conclusion (opinion personnelle avec \esp{en mi opinión}).
\end{enumerate}

\textbf{4. Pièges de francophones}
\begin{itemize}
\item \esp{los derechos} = les droits ; attention à \esp{derecho/a} (adjectif : droit, opposé de gauche).
\item Ne pas dire \esp{es necesario de} ; toujours \esp{es necesario que + subjonctif} ou \esp{es necesario + infinitif}.
\item \esp{la escolarización} ne se traduit pas par « scolarité » au sens français strict mais par « scolarisation ».
\end{itemize}
\end{bilan}

\begin{aretenir}[Checklist avant le Bac]
\begin{itemize}
\item[$\square$] Je connais le lexique des droits : \esp{igualdad, derechos, escolarización, maltrato, protección, trabajo infantil}.
\item[$\square$] Je sais situer la Convention relative aux droits de l'enfant (Nations unies, 1989).
\item[$\square$] J'exprime mon opinion avec \esp{creo que + indicatif} et \esp{no creo que + subjonctif}.
\item[$\square$] Je forme correctement le subjonctif présent des verbes réguliers (\esp{que los niños vayan, que sea, que haya}).
\item[$\square$] J'utilise \esp{es necesario que + subjonctif} sans oublier le \esp{que}.
\item[$\square$] Je distingue \esp{tener que + inf.} (obligation personnelle) et \esp{hay que + inf.} (obligation générale).
\item[$\square$] Je sais dire « être pour / contre » : \esp{estar a favor de / estar en contra de}.
\item[$\square$] Je peux citer un exemple camerounais ou hispanophone (scolarisation des filles, lutte contre le travail des enfants).
\item[$\square$] Je respecte la structure d'un commentaire : introduction, développement, conclusion avec opinion personnelle.
\item[$\square$] Je vérifie les accents : \esp{educación, protección, igualdad, niño, niña}.
\item[$\square$] Je n'oublie pas les signes \esp{¿} et \esp{¡} dans les questions et exclamations.
\item[$\square$] Je relis ma production pour éviter les calques du français (\esp{es necesario que}, pas de \esp{de}).
\end{itemize}
\end{aretenir}

\findecours

\end{document}
